\chapter{The vertex algebra of the hyperbolic plane and the tensor theory of central charge 26}\label{ch:17}
\chaptermark{The hyperbolic plane and central charge 26}
\begingroup
\emergencystretch=3em
% Local to this chapter: a Lean identifier that may break after `_' and `.' and before capitals.
\providecommand{\chseventeenL}[1]{\texttt{\directlua{ local s = "\luaescapestring{\detokenize{#1}}" local first = true for c in s:gmatch(utf8.charpattern) do if (not first) and c:match("^[A-Z]$") then tex.sprint("\string\\allowbreak{}") end tex.sprint(-2, c) if c == "_" or c == "." then tex.sprint("\string\\allowbreak{}") end first = false end}}}

\providecommand{\chseventeenII}{{\mathrm{II}_{1,1}}}
\providecommand{\chseventeenVII}{{V_{\mathrm{II}_{1,1}}}}
\providecommand{\chseventeenvac}{\mathbb{1}}
\providecommand{\chseventeenF}{\mathcal{F}}
\providecommand{\chseventeenM}{\mathcal{M}}

This chapter documents the release of the library named in the front matter; every Lean quotation in it is checked by
\texttt{scripts/textbook/sync.py} to occur verbatim (up to whitespace and comments) in the source files of that
release.  It is
the second chapter of the part ``Monstrous moonshine'' and it builds the stage on which Borcherds' proof
of the Conway--Norton conjecture is played: the vertex algebra $\chseventeenVII$ of the even unimodular
Lorentzian lattice of rank two, and the tensor product
\begin{equation*}
  \chseventeenM \;=\; \Vnat\otimes\chseventeenVII ,
\end{equation*}
a vertex algebra of central charge $24+2=26$.  Chapter~\ref{ch:18} will extract from $\chseventeenM$ the
Monster Lie algebra (the space of ``physical states'' of $\chseventeenM$, by the no-ghost theorem), and
Chapters~\ref{ch:19}--\ref{ch:20} use that Lie algebra to prove the Conway--Norton conjecture for every
automorphism of $\Vnat$.  The present chapter documents the 91 Lean files that construct
$\chseventeenVII$ and $\chseventeenM$, their invariant bilinear forms, the ``light-cone'' linear algebra of a
momentum sector, and the general vertex-algebra theory (Borcherds identity, skew-symmetry, contragredient
adjointness) for vertex algebras whose grading is \emph{not} bounded below.

\paragraph{The results in brief.}  Everything below uses only Lean's three standard axioms.  Items
(i)--(iii) have no hypotheses; the general theorems of (iv) and (v) carry the mathematical premises displayed in
their statements, which are discharged for the moonshine tensor theory in Chapter~\ref{ch:18}.  Terms used here
are defined in the sections cited.
\begin{enumerate}[label=(\roman*)]
\item \emph{The lattice} (\S\ref{ch17:sec:lattice}).  $\chseventeenII=\Z^2$ with Gram matrix
  $\left(\begin{smallmatrix}0&1\\1&0\end{smallmatrix}\right)$, norm $2xy$: even, self-dual, not positive-definite.
\item \emph{Its vertex algebra} (\S\S\ref{ch17:sec:frame}--\ref{ch17:sec:realization}).  A vertex algebra
  \chseventeenL{lorentzianNearVOA} with conformal vector and central charge $2$ on the usual lattice state space, built
  with the Frenkel--Lepowsky--Meurman vertex operators written in a \emph{complex} frame.  Its weight grading is
  not bounded below and its weight-zero space is infinite-dimensional, so it is not a vertex operator algebra; the
  library's name for this kind of object is \emph{near-VOA} (Definition~\ref{ch17:def:nearvoa}).
\item \emph{The $c=26$ theory} (\S\ref{ch17:sec:tensor}).  The tensor product $\Vnat\otimes\chseventeenVII$
  (\chseventeenL{moonshineMatterNearVOA}) has central charge exactly $26$, an internal grading by $\chseventeenII$-momentum,
  and a symmetric, normalised, nondegenerate bilinear form for which $L_n$ is adjoint to $L_{-n}$.
\item \emph{General theory} (\S\S\ref{ch17:sec:near}--\ref{ch17:sec:fhl}).  For every near-VOA: Goddard uniqueness,
  the commutator and iterate formulas, the Borcherds identity and skew-symmetry; and, when $L_1$ is locally
  nilpotent, the Frenkel--Huang--Lepowsky (FHL) adjoint of a field, with the theorem that if the generating fields
  are adjoint for a bilinear form then all fields are (premises: local nilpotence, generation, adjointness of
  generators).
\item \emph{One momentum sector, in light-cone variables} (\S\ref{ch17:sec:hyp}).  An explicit polynomial model of a
  momentum sector on which the Sugawara operators satisfy the Virasoro algebra with $c=2$, and on which the
  negative Virasoro modes and a ``longitudinal'' subspace split the sector at every nonzero momentum (premise:
  nonzero momentum); likewise after tensoring with an arbitrary matter space (premise: the Witt relations of the
  matter operators).  This splitting is the linear algebra behind the no-ghost theorem.
\end{enumerate}

\paragraph{What this chapter does \emph{not} contain.}  The no-ghost theorem itself, the physical-state
Lie algebra, the identification of its root spaces with $\Vnat_{1+mn}$, and the Borcherds--Kac--Moody
structure are in Chapter~\ref{ch:18}.  Several files of this chapter import Chapter~\ref{ch:18} objects
(the Monster Lie algebra \chseventeenL{monsterLieAlgebra}, its bilinear form \chseventeenL{monsterLieBilinearForm}, the
subspace \chseventeenL{primaryWeightOne} of weight-one Virasoro primaries); when a statement mentions them we say
so and point forward.  Nothing in this chapter is conditional on an open input: every headline is proved
with no hypotheses beyond the ones displayed in its Lean statement, and we say for each displayed
hypothesis what it means.  As everywhere in this book, our concern is the \emph{meaning} of the Lean
statements; proofs are explained mathematically, and we do not describe how the Lean proof scripts are
organised.

\paragraph{How to read this chapter.}  A reader who wants the headline only should read
\S\ref{ch17:sec:physics} (the physical motivation and the sign conventions), \S\ref{ch17:sec:lattice},
the definitions of \S\ref{ch17:sec:target}, the theorems of \S\ref{ch17:sec:realization} and
\S\ref{ch17:sec:tensor}, and the referee's checklist (\S\ref{ch17:sec:checklist}).
Sections~\ref{ch17:sec:hyp}, \ref{ch17:sec:near} and~\ref{ch17:sec:fhl} are the technical layer that
Chapter~\ref{ch:18} consumes; they are presented in full because the no-ghost theorem and the invariance
of the Monster Lie algebra's bilinear form rest on them.  Section~\ref{ch17:sec:files} accounts for every
one of the 91 files.

\paragraph{Notation.}  The chapter uses the following conventions throughout.
\begin{center}
\small
\begin{longtable}{p{0.22\textwidth}p{0.70\textwidth}}
\toprule
symbol & meaning \\ \midrule
\endhead
$\alpha,\beta,\gamma\in L$ & lattice vectors (\emph{momenta}); for $\chseventeenII$, $\alpha=(\alpha_0,\alpha_1)\in\Z^2$ \\
$(x,y)$ & the two coordinates of \emph{one} momentum of $\chseventeenII$ (as in ``momentum $(x,y)$, weight $xy$'') \\
$\Delta(\alpha)=\langle\alpha,\alpha\rangle/2$ & the \emph{ground weight} of momentum $\alpha$ (for $\chseventeenII$: $\alpha_0\alpha_1$) \\
level & the Fock-space degree of a polynomial (oscillator energy) \\
weight & the $L_0$-eigenvalue: level $+\,\Delta(\alpha)$; ``weight $N$'' always means total weight \\
$e^\alpha$ & the \emph{charged vacuum} $1\otimes e^\alpha$, the ground state of momentum $\alpha$ \\
sector & all states of a fixed momentum $\alpha$ \\
$u_{(n)}$ & the $n$-th mode, $Y(u,z)=\sum_nu_{(n)}z^{-n-1}$ (Lean \chseventeenL{mode (Y u) n}); a Lean \chseventeenL{FormalDist} \chseventeenL{f t} is instead the coefficient of $z^t$ \\
$c$ & central charge only \\
$\varepsilon,\ \chi(\alpha)$ & the FLM cocycle, and the sign $\chi(\alpha)=(-1)^{\Delta(\alpha)}\varepsilon(\alpha,\alpha)$ of the lattice form \\
$F$ & a complex lattice frame (Definition~\ref{ch17:def:frame}) \\
$\chseventeenF_r$, $\chseventeenF^{\mathrm{iso}}$ & the Fock space in orthonormal variables; the same space in light-cone variables (\S\ref{ch17:sec:hyp}) \\
$B_{\mathcal F}$, $B_L$, $B^{\mathrm{iso}}$, $B$ & the Fock form, the lattice form, the isotropic form, the matter form (table in \S\ref{ch17:sec:tensor}) \\
$E^\pm(\alpha,z)$, $P_N(\alpha)$, $A_N(\alpha)$ & the vertex-operator exponentials and their coefficients \\
$\kappa,\lambda$ & the two light-cone zero modes of a sector (Lean \chseventeenL{κ}, \chseventeenL{ell}) \\
$\varphi\in\{a,b\}$ & a light-cone oscillator family \\
$\ell,m,n$ in identities & mode indices of the Borcherds identity \\
\bottomrule
\end{longtable}
\end{center}

%======================================================================
\section{Why the hyperbolic plane: the physical picture}\label{ch17:sec:physics}
%======================================================================

\subsection{The bosonic string in twenty-six dimensions}

A closed bosonic string moving in a flat spacetime $\mathbb R^{1,d-1}$ is described, after fixing
conformal gauge on its world-sheet, by $d$ free bosons $X^\mu(z,\bar z)$, $\mu=0,\dots,d-1$, one for each
spacetime coordinate, with the Minkowski metric $\eta_{\mu\nu}=\operatorname{diag}(-1,1,\dots,1)$.  Each
holomorphic current $i\partial X^\mu(z)=\sum_n\alpha^\mu_nz^{-n-1}$ has modes satisfying
$[\alpha^\mu_m,\alpha^\nu_n]=m\,\eta^{\mu\nu}\delta_{m+n,0}$: the \emph{Lorentzian} Heisenberg algebra.  The
chiral state space is a sum over spacetime momenta $p$ of Fock spaces; the Virasoro generators
$L_n=\frac12\sum_k\eta_{\mu\nu}{:}\alpha^\mu_{n-k}\alpha^\nu_k{:}$ satisfy the Virasoro algebra with central
charge $c=d$.  Gauge fixing leaves a residual reparametrisation symmetry, and \emph{physical states} are
the states $\psi$ with
\begin{equation}\label{ch17:eq:physical}
  L_0\psi=\psi,\qquad L_n\psi=0\quad(n\ge1),
\end{equation}
taken modulo the ``null spurious'' states $L_{-1}\chi_1+L_{-2}\chi_2$ that have zero norm.  Because the
metric is indefinite, the Fock space has states of negative norm (``ghosts'', the timelike oscillators
$\alpha^0_{-n}$).  The \emph{no-ghost theorem} of Brower and of Goddard--Thorn (1972) says that when
$c=26$ the physical states modulo null states form a positive-definite space isomorphic to the
\emph{transverse} Fock space built from the $d-2=24$ oscillators orthogonal to a light-cone pair of
directions; for $c\ne26$ and $d>26$ ghosts survive.  This is the origin of the critical dimension $26$.

\begin{physics}
Twenty-six is not arbitrary here.  The $c=26$ matter theory is exactly what cancels the $c=-26$ of the
reparametrisation ghosts $b,c$ in the BRST quantisation (Frenkel--Garland--Zuckerman 1986), so that the
BRST charge squares to zero; the no-ghost theorem is equivalently the statement that the BRST cohomology
at ghost number one is the transverse space.  In light-cone gauge one solves the constraints
\eqref{ch17:eq:physical} explicitly by eliminating the two light-cone oscillators $\alpha^\pm_n$, which
leaves the $24$ transverse bosons.  Borcherds' insight was that the $24$ transverse bosons can be replaced
by \emph{any} $c=24$ vertex operator algebra, in particular by $\Vnat$.
\end{physics}

\subsection{Borcherds' string}

Replace the $24$ transverse bosons by the Moonshine module $\Vnat$ (central charge $24$,
Chapter~\ref{ch:14}), and keep two free bosons $X^0,X^1$ for the light-cone directions.  To obtain an
algebraic object with a \emph{discrete} momentum, compactify these two directions on a Lorentzian torus
$\mathbb R^{1,1}/\chseventeenII$: the chiral momenta then range over the even unimodular Lorentzian lattice
$\chseventeenII$, the unique such lattice of signature $(1,1)$.  The chiral algebra of these two bosons is
the lattice vertex algebra $\chseventeenVII$, of central charge $2$, and the full ``matter'' algebra is
$\chseventeenM=\Vnat\otimes\chseventeenVII$ of central charge $26$.  Borcherds (1992) defined the
\emph{Monster Lie algebra} as the space of weight-one Virasoro primaries of $\chseventeenM$ modulo the
radical of its invariant form, with bracket $[u,v]=u_{(0)}v$; it is graded by $\chseventeenII$, the
Monster acts on it, and by the no-ghost theorem its piece of momentum $r\ne0$ is isomorphic to the weight
$1-\tfrac12\langle r,r\rangle$ piece of $\Vnat$.  Thus the generating function of the dimensions of
$\Vnat$ (the $J$-function) becomes the root multiplicities of a Lie algebra, whose denominator identity
is $j(p)-j(q)=p^{-1}\prod_{m>0,n\in\Z}(1-p^mq^n)^{c(mn)}$ (Chapter~\ref{ch:19}).

\begin{physics}
Physically $\chseventeenM$ is the chiral algebra of a bosonic string compactified on ``the Monster CFT
times a Lorentzian torus''.  A heuristic but faithful dictionary: a lattice vector $r=(x,y)\in\chseventeenII$
is a pair of light-cone momenta; $\tfrac12\langle r,r\rangle=xy$ is $-\tfrac12p^2$ up to sign conventions, i.e.\
the mass-shell term; the mass-shell condition $L_0=1$ on a state $v\otimes e^r$ with $v\in\Vnat_k$ reads
$k+xy=1$, so the physical states of momentum $r$ are built on $\Vnat_{1-xy}$.  The sign conventions of the
repository and of Borcherds differ and are fixed in \S\ref{ch17:sec:conventions}.
\end{physics}

\subsection{Why a new kind of vertex algebra is needed}\label{ch17:sec:whynew}

In the lattice vertex operator algebra $V_L$ of a positive-definite lattice (Chapter~\ref{ch:03},
Chapter~\ref{ch:L05}) the state $p\otimes e^\alpha$, with $p$ a polynomial of oscillator level $w\ge0$,
has conformal weight $w+\tfrac12\langle\alpha,\alpha\rangle\ge0$; for each weight $n$ only finitely many
$\alpha$ satisfy $\tfrac12\langle\alpha,\alpha\rangle\le n$, so the weight spaces are finite-dimensional, and
$V_0=\C\chseventeenvac$.  For $\chseventeenII$ all three facts fail:
\begin{itemize}
\item the charged vacuum $e^{(1,N-1)}$ has weight $N-1$, which is arbitrarily negative: the grading is
  \emph{not bounded below} (Theorem~\ref{ch17:thm:notbounded});
\item the isotropic vectors $(i,0)$, $i\in\Z$, all have norm $0$, so the $e^{(i,0)}$ are infinitely many
  independent states of weight $0$: \emph{weight spaces are infinite-dimensional}
  (Theorem~\ref{ch17:thm:weightzero});
\item the library's real orthonormal frame $\alpha\mapsto\sqrt G\,\underline\alpha$ of Chapter~\ref{ch:03}
  (Definition~\ref{ch03:def:frame}) needs $G$ positive semidefinite; for $G=\left(\begin{smallmatrix}0&1\\1&0\end{smallmatrix}\right)$
  there is no real $C$ with $C^TC=G$.
\end{itemize}
Hence $\chseventeenVII$ is not a \chseventeenL{VertexOperatorAlgebra} in the library's sense (which requires
finite-dimensional weight spaces and CFT type), nor a \chseventeenL{GradedVertexOperatorAlgebra} (which requires a
lower bound).  The library's structure for it is \chseventeenL{NearVOA} (Gannon's ``near-VOA'', Definition~5.1.3(c)):
all the vertex-algebra and Virasoro axioms, an internal $\Z$-grading by $L_0$-eigenvalue, and nothing about
the size or the lower bound of the graded pieces.  Every general theorem that the library had proved for
vertex operator algebras (Chapter~\ref{ch:L03}) has to be re-proved for near-VOAs; this is the content of
\S\ref{ch17:sec:near}.

\subsection{Sign and coordinate conventions}\label{ch17:sec:conventions}

The repository uses the Gram matrix $\left(\begin{smallmatrix}0&1\\1&0\end{smallmatrix}\right)$ on $\Z^2$, so
\begin{equation*}
  \langle(x,y),(x',y')\rangle = xy'+yx',\qquad \langle(x,y),(x,y)\rangle=2xy,\qquad
  \tfrac12\langle r,r\rangle = xy .
\end{equation*}
Borcherds writes roots of the Monster Lie algebra as $(m,n)$
with norm $-2mn$.  The two conventions are matched by
\begin{equation}\label{ch17:eq:signdict}
  \text{Borcherds' }(m,n)\quad\longleftrightarrow\quad\text{repository momentum }(x,y)=(m,-n),
\end{equation}
and the root space of $(m,n)$ is $\Vnat_{1+mn}=\Vnat_{1-xy}$.

\begin{warning}
\textbf{The sign trap $(m,n)\leftrightarrow(m,-n)$.}  Never write the denominator product over repository momenta
$x,y\ge1$.  The denominator product $\prod_{m>0,n\in\Z}(1-p^mq^n)^{c(mn)}$ is
over Borcherds' coordinates; in repository coordinates the exponent is $c(-xy)$.  Every statement in this
chapter is in repository coordinates.
\end{warning}

%======================================================================
\section{The lattice \texorpdfstring{$\mathrm{II}_{1,1}$}{II(1,1)}}\label{ch17:sec:lattice}
%======================================================================

Recall (Chapter~\ref{ch:03}, Definition~\ref{ch03:def:lattice}; Chapter~\ref{ch:L05},
Definition~\ref{l05:def:lattice}) that the library's \chseventeenL{IntegralLattice} is a structure with a carrier
type that is a free, finitely generated $\Z$-module, a $\Z$-valued bilinear form \chseventeenL{form}, and a proof
\chseventeenL{isSymm} that the form is symmetric; \chseventeenL{Even} means $2\mid\langle v,v\rangle$ for all $v$,
\chseventeenL{PosDef} means $\langle v,v\rangle>0$ for $v\ne0$, and \chseventeenL{rank} is $\operatorname{finrank}_\Z$ of the
carrier.  Nothing in that structure requires positivity.

\begin{definition}[The hyperbolic plane]\label{ch17:def:lattice}
$\chseventeenII$ is $\Z^2$ with the symmetric bilinear form $\langle\alpha,\beta\rangle=\alpha_0\beta_1+\alpha_1\beta_0$, i.e.\ Gram
matrix $\left(\begin{smallmatrix}0&1\\1&0\end{smallmatrix}\right)$ in the standard basis.  Its \emph{ground weight}
(half-norm) is $\Delta(\alpha)=\alpha_0\alpha_1$.
\end{definition}

\noindent File \leanpath{VOA/\allowbreak Moonshine/\allowbreak MoonshineCNLattice.lean}.
\begin{leancode}
/-- Gram matrix of the hyperbolic plane in its isotropic integral basis. -/
def hyperbolicGram : Matrix (Fin 2) (Fin 2) ℤ := !![0, 1; 1, 0]

/-- The integral hyperbolic plane `II_{1,1}`, with norm `2xy`. -/
noncomputable def lorentzianLattice : IntegralLattice where
  carrier := Fin 2 → ℤ
  form := Matrix.toBilin' hyperbolicGram

theorem lorentzian_form_apply (x y : lorentzianLattice.carrier) :
    lorentzianLattice.form x y = x 0 * y 1 + x 1 * y 0 := by

/-- Half the norm, the integral ground-state conformal weight. -/
def lorentzianEnergy (x : lorentzianLattice.carrier) : ℤ := x 0 * x 1
\end{leancode}

\begin{dictionary}
\chseventeenL{Fin 2 → ℤ} is $\Z^2$, a vector being a function of its two coordinates \chseventeenL{x 0}, \chseventeenL{x 1}.
\chseventeenL{!![0, 1; 1, 0]} is Mathlib's notation for the $2\times2$ matrix with rows $(0,1)$ and $(1,0)$.
\chseventeenL{Matrix.toBilin' M} is the bilinear form $(x,y)\mapsto\sum_{i,j}x_iM_{ij}y_j$; with this $M$ it is
$x_0y_1+x_1y_0$, which is exactly what \chseventeenL{lorentzian_form_apply} records, so a referee need not trust
the unfolding.  The omitted field \chseventeenL{isSymm} is the proof of symmetry.  The full names are
\chseventeenL{VOA.Moonshine.ConwayNorton.lorentzianLattice} etc.; every declaration of this chapter not in
namespace \chseventeenL{VOA.NearVOA} lives in \chseventeenL{VOA.Moonshine.ConwayNorton}.  The basis vectors $u=(1,0)$ and
$v=(0,1)$ are \emph{isotropic} ($\langle u,u\rangle=\langle v,v\rangle=0$) with $\langle u,v\rangle=1$; this is
why the library calls the standard basis ``isotropic''.
\end{dictionary}

\begin{physics}
$\chseventeenII$ is the momentum lattice of two chiral bosons compactified on a Lorentzian torus; in the
basis $u,v$ the two bosons are the light-cone combinations $X^\pm$.  The ground weight $\alpha_0\alpha_1$ is the conformal
weight of the charged ground state $e^{(\alpha_0,\alpha_1)}$ --- the ``$p^+p^-$'' mass-shell term of light-cone
quantisation.
\end{physics}

\begin{status}
Definitions.  \chseventeenL{lorentzian_form_apply} is proved; printed in
\leanpath{VOA/\allowbreak Moonshine/\allowbreak MoonshineCNAudit.lean} with axioms \chseventeenL{[propext, Classical.choice, Quot.sound]}.
\end{status}

\begin{theorem}[$\chseventeenII$ is even, unimodular, of rank two and indefinite]\label{ch17:thm:latticeprops}
The lattice $\chseventeenII$ is even; its rank is $2$; the determinant of its Gram matrix is $-1$; its form
is nondegenerate; it is not positive-definite; and it is self-dual: an element $x\in\chseventeenII\otimes\mathbb R$
has $\langle x,v\rangle\in\Z$ for all $v\in\chseventeenII$ if and only if $x\in\chseventeenII$.
\end{theorem}

\noindent File \leanpath{VOA/\allowbreak Moonshine/\allowbreak MoonshineCNLattice.lean}.
\begin{leancode}
theorem lorentzian_even : lorentzianLattice.Even := by

theorem lorentzian_rank : rank lorentzianLattice = 2 := by

theorem lorentzian_gram_det : hyperbolicGram.det = -1 := by

theorem lorentzian_form_nondegenerate : lorentzianLattice.form.Nondegenerate := by

theorem lorentzian_not_posDef : ¬ lorentzianLattice.PosDef := by

/-- Unimodularity in the repository's actual real dual-lattice sense. -/
theorem lorentzian_selfDual : lorentzianLattice.SelfDual := by
\end{leancode}

\begin{dictionary}
\chseventeenL{SelfDual} (file \leanpath{VOA/\allowbreak Foundations/\allowbreak Lattices.lean}) says: for every
$x\in\mathbb R\otimes_\Z L$, $x$ lies in the dual lattice \chseventeenL{L.dual}
$=\{x:\langle x,v\rangle\in\Z\ \forall v\in L\}$ iff $x$ is the image of a lattice vector.  This is
unimodularity in its intrinsic form; \chseventeenL{lorentzian_gram_det} gives the determinant form
$\det G=-1$ as well, and together with \chseventeenL{lorentzian_not_posDef} it pins the signature: a real symmetric
$2\times2$ matrix with negative determinant has one positive and one negative eigenvalue, so the signature
is $(1,1)$.  \chseventeenL{Nondegenerate} is Mathlib's \chseventeenL{LinearMap.BilinForm.Nondegenerate} (left and right
separating).  The witness for \chseventeenL{lorentzian_not_posDef} is $(1,-1)$, of norm $-2$.  \textbf{Fidelity.}
The six statements are exactly the six properties in the theorem; nothing is assumed.  ``Even unimodular
of signature $(1,1)$'' characterises $\chseventeenII$ up to isomorphism (classical; not formalised and not needed).
\end{dictionary}

\begin{proof}[Proof on paper]
$\langle x,x\rangle=2x_0x_1$ is even.  The standard basis has two elements.  $\det\left(\begin{smallmatrix}0&1\\1&0\end{smallmatrix}\right)=-1$.
If $\langle x,y\rangle=0$ for $y=(0,1)$ and $(1,0)$ then $x_0=x_1=0$.  $\langle(1,-1),(1,-1)\rangle=-2<0$.  For
self-duality: $\langle x,(0,1)\rangle$ and $\langle x,(1,0)\rangle$ are the two real coordinates of $x$ in reversed
order, so both are integers iff $x\in\Z^2$.
\end{proof}

\begin{physics}
Evenness makes every vertex operator $e^\alpha$ bosonic with integral weight.  Unimodularity (self-duality, as
for $E_8$ and $\Leech$) means there are no further sectors $L^*/L$ to add; for a definite lattice it is also the
condition for a modular-invariant chiral partition function, but for $\chseventeenII$ no character
$\operatorname{tr}q^{L_0}$ exists (Theorem~\ref{ch17:thm:weightzero}), so that reading does not apply.  What the
moonshine application uses is self-duality: it makes the Monster Lie algebra
graded by the \emph{whole} lattice $\chseventeenII$ rather than by a coset.  Indefiniteness is the presence of
time: one of the two directions is timelike.
\end{physics}

\begin{status}
All six proved unconditionally; printed in \leanpath{VOA/\allowbreak Moonshine/\allowbreak MoonshineCNAudit.lean}
(\chseventeenL{lorentzian_not_posDef}, \chseventeenL{lorentzian_gram_det}, \chseventeenL{lorentzian_selfDual} checked there);
axioms \chseventeenL{[propext, Classical.choice, Quot.sound]}.
\end{status}

\begin{theorem}[Energy]\label{ch17:thm:energy}
For $\alpha,\beta\in\chseventeenII$: $2\Delta(\alpha)=\langle\alpha,\alpha\rangle$; $\Delta(\alpha+\beta)=\Delta(\alpha)+\Delta(\beta)+\langle\alpha,\beta\rangle$; and for every
$N\in\Z$ there is $\alpha$ with $\Delta(\alpha)<N$ (namely $\alpha=(1,N-1)$).
\end{theorem}

\noindent File \leanpath{VOA/\allowbreak Moonshine/\allowbreak MoonshineCNLattice.lean}.
\begin{leancode}
theorem lorentzianEnergy_half_norm (x : lorentzianLattice.carrier) :
    2 * lorentzianEnergy x = lorentzianLattice.form x x := by

theorem lorentzianEnergy_add (x y : lorentzianLattice.carrier) :
    lorentzianEnergy (x + y) = lorentzianEnergy x + lorentzianEnergy y +
      lorentzianLattice.form x y := by

/-- Charged ground states have no uniform lower energy bound. -/
theorem lorentzianEnergy_unbounded_below (N : ℤ) :
    ∃ x : lorentzianLattice.carrier, lorentzianEnergy x < N := by
\end{leancode}

\begin{dictionary}
\chseventeenL{lorentzianEnergy} is $\Delta$, defined directly as $\alpha_0\alpha_1$ rather than as the library's
\chseventeenL{halfNorm L x} $=\lfloor\langle x,x\rangle/2\rfloor$ (the Lean binder \chseventeenL{x} is a whole lattice vector) (Chapter~\ref{ch:03}); the two agree by
\chseventeenL{lorentzian_halfNorm} (Theorem~\ref{ch17:thm:grading} below).  The third statement is the honest
form of ``unbounded below'': an existential for every integer bound.
\end{dictionary}

\begin{proof}[Proof on paper]
$(\alpha_0+\beta_0)(\alpha_1+\beta_1)=\alpha_0\alpha_1+\beta_0\beta_1+(\alpha_0\beta_1+\alpha_1\beta_0)$, and $1\cdot(N-1)<N$.
\end{proof}

\begin{physics}
$\Delta(\alpha+\beta)=\Delta(\alpha)+\Delta(\beta)+\langle\alpha,\beta\rangle$ is the statement that the operator product $e^\alpha(z)e^\beta(w)$ behaves like
$(z-w)^{\langle\alpha,\beta\rangle}e^{\alpha+\beta}$: weights add up to the correction $\langle\alpha,\beta\rangle$ from the OPE.  The
unbounded energy is the tachyonic/time-like direction of the Lorentzian torus: in a Lorentzian
compactification there is no lowest-energy state.
\end{physics}

\begin{status}
Proved unconditionally; \chseventeenL{lorentzianEnergy_unbounded_below} is printed in
\leanpath{VOA/\allowbreak Moonshine/\allowbreak MoonshineCNAudit.lean}; standard axioms.
\end{status}

\begin{definition}[The cocycle]\label{ch17:def:cocycle}
$\varepsilon=\chseventeenL{lorentzianCocycle}$ is the library's FLM cocycle \chseventeenL{latticeCocycle} of
Chapter~\ref{ch:03} (Definition~\ref{ch03:def:cocExp}), applied to $\chseventeenII$.  It is normalised,
satisfies the 2-cocycle law, takes values in $\{\pm1\}$ (so $\varepsilon(\alpha,\beta)^{-1}=\varepsilon(\alpha,\beta)$), and
$\varepsilon(\alpha,\beta)\varepsilon(\beta,\alpha)=(-1)^{\langle\alpha,\beta\rangle}$; together these give
$\varepsilon(\alpha,\beta)=(-1)^{\langle\alpha,\beta\rangle}\varepsilon(\beta,\alpha)$.
\end{definition}

\noindent File \leanpath{VOA/\allowbreak Moonshine/\allowbreak MoonshineCNLattice.lean}.
\begin{leancode}
/-- The existing FLM cocycle works without positive definiteness. -/
noncomputable def lorentzianCocycle :
    lorentzianLattice.carrier → lorentzianLattice.carrier → ℂ :=
  latticeCocycle lorentzianLattice

theorem lorentzianCocycle_law (x y z : lorentzianLattice.carrier) :
    lorentzianCocycle x y * lorentzianCocycle (x + y) z =
      lorentzianCocycle y z * lorentzianCocycle x (y + z) :=

theorem lorentzianCocycle_commutator (x y : lorentzianLattice.carrier) :
    lorentzianCocycle x y * lorentzianCocycle y x =
      (-1 : ℂ) ^ (lorentzianLattice.form x y) :=
\end{leancode}

\begin{dictionary}
The FLM cocycle is built from a chosen $\Z$-basis and the Gram matrix only, so positivity never enters;
the commutator property uses evenness (\chseventeenL{lorentzian_even}).  \chseventeenL{(-1 : ℂ) ^ k} with \chseventeenL{k : ℤ}
is an integer power (\chseventeenL{zpow}).  A second, simpler cocycle $(-1)^{x_1y_0}$ is introduced in
Definition~\ref{ch17:def:stdcocycle}; it is \emph{not} the one used in the vertex algebra.
\end{dictionary}

\begin{physics}
The cocycle makes the vertex operators $e^x(z)$ and $e^y(w)$ mutually local rather than local up to the
sign $(-1)^{\langle x,y\rangle}$: it implements the Klein factors that make the two compact bosons' vertex
operators commute (rather than anticommute) when the pairing is odd.
\end{physics}

\begin{status}
Proved unconditionally (re-uses \chseventeenL{latticeCocycle_cocycleLaw}, \chseventeenL{latticeCocycle_comm}); standard
axioms (\leanpath{VOA/\allowbreak Moonshine/\allowbreak MoonshineCNAudit.lean}).
\end{status}

%======================================================================
\section{Oscillators in a complex frame}\label{ch17:sec:frame}
%======================================================================

\subsection{The problem, and the fix}

The library's Fock space $\chseventeenF_r=\chseventeenL{FockN r}=\C[x_{j,a}:j\in\N,\ a\in\{1..r\}]$
(Chapter~\ref{ch:03}, Definition~\ref{ch03:def:FockN}) carries \emph{orthonormal} oscillators $h_a(n)$:
$h_a(-n)$ ($n\ge1$) is multiplication by $x_{n-1,a}$ and $h_a(n)$ ($n\ge1$) is $n\,\partial/\partial x_{n-1,a}$,
with $[h_a(m),h_b(n)]=m\delta_{ab}\delta_{m+n,0}$.  To realise the lattice Heisenberg algebra
$[\alpha(m),\beta(n)]=m\langle\alpha,\beta\rangle\delta_{m+n,0}$ one needs coordinates $\alpha^\perp\in\C^r$ with
\begin{equation}\label{ch17:eq:framepairing}
  \sum_{a}\alpha^\perp_a\beta^\perp_a = \langle\alpha,\beta\rangle\qquad(\alpha,\beta\in L),
\end{equation}
and then $\alpha(n):=\sum_a\alpha^\perp_ah_a(n)$.  For positive-definite $L$ the library used the real matrix
$\sqrt G$ (Chapter~\ref{ch:03}, \S\ref{ch03:sec:frame}).  Over $\C$ the pairing \eqref{ch17:eq:framepairing}
is \emph{bilinear} (no complex conjugation), and every nondegenerate symmetric form over $\C$ has such
coordinates.  For $\chseventeenII$ one may take
\begin{equation}\label{ch17:eq:complexframe}
  (x_0,x_1)^\perp=\Bigl(x_0+\tfrac{x_1}2,\ i\bigl(x_0-\tfrac{x_1}2\bigr)\Bigr),
\end{equation}
since $(x_0+\tfrac{x_1}2)(y_0+\tfrac{y_1}2)-(x_0-\tfrac{x_1}2)(y_0-\tfrac{y_1}2)=x_0y_1+x_1y_0$.

\begin{definition}[Complex lattice frame]\label{ch17:def:frame}
A \emph{complex lattice frame} for an integral lattice $L$ of rank $r$ is an additive map
$\alpha\mapsto\alpha^\perp$, $L\to\C^r$, satisfying \eqref{ch17:eq:framepairing}.  The \emph{frame oscillator}
of $\alpha$ is $\alpha(m)=\sum_a\alpha^\perp_a h_a(m)$ on $\chseventeenF_r$.  The Lorentzian frame is
\eqref{ch17:eq:complexframe}, reindexed along $\operatorname{rank}\chseventeenII=2$.
\end{definition}

\noindent{\raggedright Files \leanpath{VOA/\allowbreak Moonshine/\allowbreak MoonshineCNComplexFrame.lean},
\leanpath{VOA/\allowbreak Moonshine/\allowbreak MoonshineCNOscillators.lean}.\par}
\begin{leancode}
/-- Complex bilinear coordinates for the integral lattice pairing. -/
structure ComplexLatticeFrame (L : IntegralLattice) where
  coordinate : L.carrier →+ (Fin (rank L) → ℂ)
  pairing : ∀ x y, ∑ a, coordinate x a * coordinate y a = (L.form x y : ℂ)

/-- An explicit complex frame: `(x+y/2, i(x-y/2))` has bilinear norm `2xy`. -/
noncomputable def complexHyperbolicCoordinate (x : lorentzianLattice.carrier) : Fin 2 → ℂ :=
  ![(x 0 : ℂ) + (x 1 : ℂ) / 2, Complex.I * ((x 0 : ℂ) - (x 1 : ℂ) / 2)]

/-- The explicit hyperbolic frame, reindexed to the existing Fock carrier. -/
noncomputable def lorentzianComplexFrame : ComplexLatticeFrame lorentzianLattice where
  coordinate :=
    { toFun := fun x a => complexHyperbolicCoordinate x (finCongr lorentzian_rank a)

/-- The oscillator defined by an arbitrary complex lattice frame. -/
noncomputable def frameOscillator {L : IntegralLattice} (F : ComplexLatticeFrame L)
    (α : L.carrier) (m : ℤ) : Module.End ℂ (FockN (rank L)) :=
  ∑ a, F.coordinate α a • oscN (rank L) 1 a m
\end{leancode}

\begin{dictionary}
\chseventeenL{L.carrier →+ (Fin (rank L) → ℂ)} is an additive group homomorphism into $\C^r$ (indexed by
\chseventeenL{Fin (rank L)}, the Fock colours).  \chseventeenL{pairing} is \eqref{ch17:eq:framepairing} with the integer
$\langle x,y\rangle$ cast to $\C$; note there is no complex conjugate, so this is a \emph{bilinear}
orthonormal frame.  \chseventeenL{oscN (rank L) 1 a m} is $h_a(m)$ at level $k=1$ (Chapter~\ref{ch:03},
Definition~\ref{ch03:def:oscN}).  \chseventeenL{finCongr lorentzian_rank} is the bijection
$\{0,\dots,\operatorname{rank}-1\}\cong\{0,1\}$ coming from $\operatorname{rank}\chseventeenII=2$: the Fock
space is indexed by \chseventeenL{Fin (rank lorentzianLattice)}, which is only propositionally equal to
\chseventeenL{Fin 2}.  The omitted fields of \chseventeenL{lorentzianComplexFrame} are the additivity proofs and the
pairing proof (which is \chseventeenL{complexHyperbolicCoordinate_pairing}, Theorem~\ref{ch17:thm:frame}).
\textbf{Could this be vacuous?}  A \chseventeenL{ComplexLatticeFrame} exists only if the coordinates realise the form;
the Lorentzian instance is \emph{constructed}, not assumed, so the structure is inhabited for the lattice we
care about.
\end{dictionary}

\begin{physics}
Over $\C$ every nondegenerate symmetric bilinear form is a sum of squares, so two bosons with the Lorentzian
OPE can always be rewritten as two ``orthonormal'' bosons with complex coefficients.  In
\eqref{ch17:eq:complexframe} the form is $A^2-B^2$ with $A=x_0+\frac{x_1}2$ (the spacelike combination) and
$B=x_0-\frac{x_1}2$ (the timelike one); the factor $i$ in front of $B$ turns $-B^2$ into $+(iB)^2$.  The Fock
machinery is purely algebraic (polynomials and derivations), so complex coefficients cost nothing and change no
OPE: the commutators are the lattice pairing (Theorem~\ref{ch17:thm:frame}).
\end{physics}

\begin{status}
Definitions; the Lorentzian frame is constructed with no hypotheses
(\chseventeenL{lorentzianComplexFrame}, printed in \leanpath{VOA/\allowbreak Moonshine/\allowbreak MoonshineCNAudit2.lean}).
\end{status}

\begin{theorem}[Frame pairing and Heisenberg relations]\label{ch17:thm:frame}
The coordinates \eqref{ch17:eq:complexframe} satisfy \eqref{ch17:eq:framepairing} for $\chseventeenII$.  For any
complex frame of any integral lattice, the frame oscillators satisfy
$[\alpha(m),\beta(n)]=m\langle\alpha,\beta\rangle\delta_{m+n,0}$, and $\alpha(m)\chseventeenvac=0$ for $m\ge0$.
\end{theorem}

\noindent{\raggedright Files \leanpath{VOA/\allowbreak Moonshine/\allowbreak MoonshineCNOscillators.lean},
\leanpath{VOA/\allowbreak Moonshine/\allowbreak MoonshineCNComplexFrame.lean}.\par}
\begin{leancode}
theorem complexHyperbolicCoordinate_pairing (x y : lorentzianLattice.carrier) :
    (∑ a, complexHyperbolicCoordinate x a * complexHyperbolicCoordinate y a) =
      (lorentzianLattice.form x y : ℂ) := by

theorem lorentzianOscillator_commutator (x y : lorentzianLattice.carrier) (m n : ℤ) :
    ⁅lorentzianOscillator x m, lorentzianOscillator y n⁆ =
      (if m + n = 0 then (lorentzianLattice.form x y : ℂ) * (m : ℂ) else 0) •
        LinearMap.id := by

theorem frameOscillator_vacuum {L : IntegralLattice} (F : ComplexLatticeFrame L)
    (α : L.carrier) (m : ℤ) (hm : 0 ≤ m) :
    frameOscillator F α m (vacuumN (rank L)) = 0 :=

theorem frameOscillator_commutator {L : IntegralLattice} (F : ComplexLatticeFrame L)
    (α β : L.carrier) (m n : ℤ) :
    frameOscillator F α m * frameOscillator F β n -
      frameOscillator F β n * frameOscillator F α m =
      (if m + n = 0 then (L.form α β : ℂ) * (m : ℂ) else 0) • LinearMap.id := by
\end{leancode}

\begin{dictionary}
\chseventeenL{⁅A, B⁆} is the commutator $AB-BA$ of endomorphisms; the second commutator theorem writes it out.
\chseventeenL{lorentzianOscillator x m} is \chseventeenL{complexChargeOscillator (complexHyperbolicCoordinate x) m}, the
same operator as \chseventeenL{frameOscillator lorentzianComplexFrame x m} up to the reindexing by
\chseventeenL{finCongr} (on \chseventeenL{FockN 2} directly).  The \chseventeenL{if m + n = 0 then … else 0} is the Kronecker delta.
\chseventeenL{vacuumN} is the constant polynomial $1=\chseventeenvac$.  \textbf{Fidelity:} exactly the Heisenberg
relations with the \emph{lattice} form, which is what the lattice vertex algebra needs; the hypotheses are
only the data (a frame), no positivity.
\end{dictionary}

\begin{proof}[Proof on paper]
The pairing is the computation after \eqref{ch17:eq:complexframe}, using $i^2=-1$.  Expanding
$\alpha(m),\beta(n)$ in the $h_a$, only diagonal terms survive and give
$m\delta_{m+n,0}\sum_a\alpha^\perp_a\beta^\perp_a=m\delta_{m+n,0}\langle\alpha,\beta\rangle$.  Annihilation of the vacuum
holds for each $h_a(m)$, $m\ge0$ (for $m=0$ the oscillator $h_a(0)$ is zero on the Fock space; the momentum
zero mode is added separately in Definition~\ref{ch17:def:currents}).
\end{proof}

\begin{physics}
$[\alpha(m),\beta(n)]=m\langle\alpha,\beta\rangle\delta_{m+n,0}$ is the OPE $\partial X^\alpha(z)\partial X^\beta(w)\sim
\langle\alpha,\beta\rangle/(z-w)^2$ of two Lorentzian currents.  For the isotropic pair $u=(1,0)$, $v=(0,1)$ it
says $[u(m),u(n)]=[v(m),v(n)]=0$ and $[u(m),v(n)]=m\delta_{m+n,0}$: the light-cone oscillators $\alpha^+$ and
$\alpha^-$ commute among themselves and are conjugate to each other.
\end{physics}

\begin{status}
Proved unconditionally; \chseventeenL{lorentzianOscillator_commutator} printed in
\leanpath{VOA/\allowbreak Moonshine/\allowbreak MoonshineCNAudit.lean}, the frame statements in
\leanpath{VOA/\allowbreak Moonshine/\allowbreak MoonshineCNAudit2.lean}; standard axioms.
\end{status}

\begin{definition}[The coordinate cocycle]\label{ch17:def:stdcocycle}
$\varepsilon_{\mathrm{std}}(x,y)=(-1)^{x_1y_0}$ on $\chseventeenII$.
\end{definition}

\noindent File \leanpath{VOA/\allowbreak Moonshine/\allowbreak MoonshineCNComplexFrame.lean}.
\begin{leancode}
/-- Fixed-coordinate sign cocycle: the ordered isotropic basis gives `(-1)^(x₁y₀)`. -/
noncomputable def standardLorentzianCocycle (x y : lorentzianLattice.carrier) : ℂ :=
  (-1 : ℂ) ^ (x 1 * y 0)

theorem standardLorentzianCocycle_commutator (x y : lorentzianLattice.carrier) :
    standardLorentzianCocycle x y * standardLorentzianCocycle y x =
      (-1 : ℂ) ^ lorentzianLattice.form x y := by
\end{leancode}

\begin{dictionary}
The file also proves the cocycle law, normalisation and nonvanishing.  $\varepsilon_{\mathrm{std}}$ is
bimultiplicative, and $\varepsilon_{\mathrm{std}}(x,y)\varepsilon_{\mathrm{std}}(y,x)=(-1)^{x_1y_0+y_1x_0}=(-1)^{\langle x,y\rangle}$.
\textbf{Status in the construction:} it is used only for a parallel family of charged fields
(\chseventeenL{lorentzianChargedField}, Definition~\ref{ch17:def:charged}) whose locality is proved; the vertex
algebra \chseventeenL{lorentzianNearVOA} uses the library's \chseventeenL{latticeCocycle} instead, because the target
structure prescribes it (Definition~\ref{ch17:def:target}).  The two cocycles are \emph{not} claimed equal,
and no gauge transformation between them is formalised (see the warning in \S\ref{ch17:sec:traps}).
\end{dictionary}

\begin{physics}
Any two cocycles with the same commutator are cohomologous; the resulting vertex algebras are isomorphic
by rescaling each $e^\alpha$ by a sign.  Physically the choice of Klein factors is unobservable.
\end{physics}

\begin{status}
Definition and proved properties; standard axioms (\leanpath{VOA/\allowbreak Moonshine/\allowbreak MoonshineCNAudit2.lean}).
\end{status}

%======================================================================
\section{What kind of object \texorpdfstring{$V_{\mathrm{II}_{1,1}}$}{V(II(1,1))} is}\label{ch17:sec:target}
%======================================================================

\subsection{Near-VOAs}

We first print the structure that the whole chapter is about.  It lives in the core of the library and is
mentioned in Chapter~\ref{ch:01} and Chapter~\ref{ch:L03}; since every headline of this part of the book
is a statement about it, we decode it in full here.  Recall the conventions of Chapter~\ref{ch:01}: a
\chseventeenL{FormalDist (End ℂ 𝒱)}, also written \chseventeenL{(End ℂ 𝒱)⟦z^±⟧}, is a function $\Z\to\End(\mathcal V)$,
the value at $t$ being the coefficient of $z^t$; \chseventeenL{mode f n} $=f(-n-1)$, so that \chseventeenL{mode (Y u) n}
is $u_{(n)}$ in $Y(u,z)=\sum_nu_{(n)}z^{-n-1}$; \chseventeenL{Local A B} means that $(z_1-z_2)^N[A(z_1),B(z_2)]=0$
for some $N\in\N$; \chseventeenL{deriv} is the formal $\partial_z$.

\begin{definition}[Near-VOA, following Gannon]\label{ch17:def:nearvoa}
A \emph{near-VOA} on a complex vector space $\mathcal V$ consists of: a family of subspaces
$\mathcal V_n$, $n\in\Z$, with $\mathcal V=\bigoplus_n\mathcal V_n$ (internal direct sum); a linear map
$Y:\mathcal V\to\End(\mathcal V)[[z^{\pm1}]]$; a vacuum $\chseventeenvac\in\mathcal V_0$; a conformal vector
$\omega\in\mathcal V_2$; and a central charge $c\in\C$; such that, writing $L_n=\omega_{(n+1)}$:
(VA1) $u_{(n)}\mathcal V_j\subseteq\mathcal V_{k+j-n-1}$ for $u\in\mathcal V_k$;
(VA2) $Y(\chseventeenvac,z)=\mathrm{id}$;
(VA3) $u_{(-1)}\chseventeenvac=u$ and $u_{(n)}\chseventeenvac=0$ for $n\ge0$;
(VA4) all $Y(u,z),Y(v,z)$ are mutually local;
(VA5) for all $u,v$, $u_{(n)}v=0$ for $n\gg0$;
(VOA1) $[L_m,L_n]=(m-n)L_{m+n}+\tfrac{c}{12}(m^3-m)\delta_{m+n,0}$;
(VOA2) $L_0=n$ on $\mathcal V_n$;
(VOA3) $Y(L_{-1}v,z)=\partial_zY(v,z)$.
No finiteness of $\dim\mathcal V_n$ and no lower bound on $n$ is required.
\end{definition}

\noindent File \leanpath{VOA/\allowbreak Core/\allowbreak VertexOperatorAlgebra.lean}.
\begin{leancode}
structure NearVOA (𝒱 : Type u) [AddCommGroup 𝒱] [Module ℂ 𝒱] where
  /-- The internal `ℤ`-grading (pieces possibly infinite-dimensional). -/
  grading : ℤ → Submodule ℂ 𝒱
  /-- The grading is an internal direct-sum decomposition of `𝒱`. -/
  isInternal : DirectSum.IsInternal grading
  /-- The state–field correspondence. -/
  Y : 𝒱 →ₗ[ℂ] (End ℂ 𝒱)⟦z^±⟧
  /-- The vacuum vector. -/
  vacuum : 𝒱
  /-- `𝟙 ∈ 𝒱₀`. -/
  vacuum_mem : vacuum ∈ grading 0
  /-- **VA1.** -/
  grading_mode : ∀ {k ℓ : ℤ} (u : 𝒱), u ∈ grading k → ∀ (v : 𝒱), v ∈ grading ℓ →
    ∀ (n : ℤ), mode (Y u) n v ∈ grading (k + ℓ - n - 1)
  /-- **VA2.** -/
  vacuum_Y : Y vacuum = (1 : (End ℂ 𝒱)⟦z^±⟧)
  /-- **VA3.** -/
  creation : ∀ u : 𝒱, mode (Y u) (-1) vacuum = u ∧ ∀ n : ℤ, 0 ≤ n → mode (Y u) n vacuum = 0
  /-- **VA4 (locality).** -/
  locality : ∀ u v : 𝒱, Local (Y u) (Y v)
  /-- **VA5 (regularity).** -/
  regularity : ∀ u v : 𝒱, ∃ N : ℤ, ∀ n : ℤ, N ≤ n → mode (Y u) n v = 0
  /-- The conformal vector `ω`. -/
  conformalVector : 𝒱
  /-- `ω ∈ 𝒱₂`. -/
  conformal_mem : conformalVector ∈ grading 2
  /-- The central charge `c`. -/
  centralCharge : ℂ
  /-- **VOA1.** -/
  virasoro : ∀ m n : ℤ,
      ⁅mode (Y conformalVector) (m + 1), mode (Y conformalVector) (n + 1)⁆
        = ((m : ℂ) - n) • mode (Y conformalVector) (m + n + 1)
          + (if m + n = 0 then (((m : ℂ) ^ 3 - m) / 12 * centralCharge) • (1 : End ℂ 𝒱)
             else 0)
  /-- **VOA2.** -/
  L0_grading : ∀ {n : ℤ} (v : 𝒱), v ∈ grading n → mode (Y conformalVector) 1 v = (n : ℂ) • v
  /-- **VOA3.** -/
  translation : ∀ v : 𝒱, Y (mode (Y conformalVector) 0 v) = deriv (Y v)
\end{leancode}

\begin{dictionary}
\textbf{Field by field.} \chseventeenL{DirectSum.IsInternal grading} is Mathlib's statement that the canonical map
$\bigoplus_n\mathcal V_n\to\mathcal V$ is bijective (every vector is uniquely a finite sum of homogeneous
pieces).  \chseventeenL{𝒱 →ₗ[ℂ] …} is a $\C$-linear map.  In \chseventeenL{virasoro}, \chseventeenL{mode (Y conformalVector) (m + 1)}
$=\omega_{(m+1)}=L_m$; the central term is $\frac{m^3-m}{12}c\,\mathrm{id}$ when $m+n=0$.  In
\chseventeenL{L0_grading}, \chseventeenL{mode (Y conformalVector) 1} $=L_0$.  In \chseventeenL{translation},
\chseventeenL{mode (Y conformalVector) 0} $=L_{-1}$.  Locality is ordinary (bosonic) locality.
\textbf{What is absent, compared with \chseventeenL{VertexOperatorAlgebra}} (Chapter~\ref{ch:L03},
Definition~\ref{l03:def:VOA}).  The library layers its structures: \chseventeenL{PreVertexAlgebra} carries the field
\chseventeenL{finrank} (finite-dimensional weight spaces); \chseventeenL{GradedVertexOperatorAlgebra} extends it by the
conformal data and \chseventeenL{grading_bounded_below} ($\mathcal V_n=0$ for $n\ll0$); \chseventeenL{VertexOperatorAlgebra}
adds the CFT-type fields \chseventeenL{vacuum_grade} ($\mathcal V_0=\C\chseventeenvac$) and \chseventeenL{positivity}
($\mathcal V_n=0$ for $n<0$).  A near-VOA has none of \chseventeenL{finrank}, \chseventeenL{grading_bounded_below},
\chseventeenL{vacuum_grade}, \chseventeenL{positivity}.
Every VOA gives a near-VOA by \chseventeenL{VertexOperatorAlgebra.toNearVOA}, which forgets these fields; this is
how $\Vnat$ enters below (\chseventeenL{VnatVOA.toNearVOA}).
\textbf{Possible vacuity.} The zero space is a near-VOA of every central charge; non-vacuity of the specific
near-VOAs below is witnessed by their explicit nonzero states (e.g.\ \chseventeenL{lorentzianChargeState_ne_zero}) and
by their central charges being forced by the Virasoro relation on a nonzero space.  Note also that the
structure contains \emph{both} the ordinary locality axiom and translation covariance (VOA3); the counterexample
of Chapter~\ref{ch:L03} (\S\ref{l03:sec:counter}) shows that without (VOA3) the Borcherds identity can fail, so
its presence matters for \S\ref{ch17:sec:near}.
\end{dictionary}

\begin{physics}
A near-VOA is a chiral algebra of a (possibly non-unitary, non-compact) conformal field theory: states,
fields, OPEs, a stress tensor with central charge $c$, and $L_0$ as the energy.  What it does not have is a
vacuum that is the unique lowest-energy state; Lorentzian or non-compact target spaces (the light-cone
bosons, a linear dilaton, Liouville theory) produce exactly such algebras.
\end{physics}

\begin{status}
Definition (core library); \chseventeenL{toNearVOA} is axiom-clean (Chapter~\ref{ch:01}).
\end{status}

\subsection{The target structure and its obstruction}

Before constructing any field, the library states precisely which object is to be built.

\begin{definition}[Even-lattice vertex realisation]\label{ch17:def:target}
An \emph{even-lattice vertex realisation} of an integral lattice $L$ is a near-VOA
(Definition~\ref{ch17:def:nearvoa}) on the lattice state
space $V_L=\chseventeenF_r\otimes\C[L]$ (Chapter~\ref{ch:03}, Definition~\ref{ch03:def:LatModuleGen})
together with: $L$ even and nondegenerate; vacuum $=1\otimes e^0$; central charge $=\operatorname{rank}L$;
each charged vacuum $e^x=1\otimes e^x$ in weight $\langle x,x\rangle/2$; and the leading product of charged
vacua
\begin{equation}\label{ch17:eq:chargeleading}
  e^x_{(-\langle x,y\rangle-1)}e^y=\varepsilon(x,y)\,e^{x+y},\qquad e^x_{(n)}e^y=0\quad(n>-\langle x,y\rangle-1),
\end{equation}
with $\varepsilon=$ \chseventeenL{latticeCocycle}.  \chseventeenL{LorentzianVertexTarget} is the proposition that such a
realisation exists for $\chseventeenII$.
\end{definition}

\noindent File \leanpath{VOA/\allowbreak Moonshine/\allowbreak MoonshineCNVertexTarget.lean}.
\begin{leancode}
/-- Target data for an even lattice, including the correct charge-state weights.
An inhabitant, not merely this structure declaration, is still required. -/
structure EvenLatticeVertexRealization (L : IntegralLattice) extends
    UnboundedConformalVertexAlgebra (LatGen.LatModuleGen L) where
  even : L.Even
  form_nondegenerate : L.form.Nondegenerate
  vacuum_eq : vacuum = LatGen.latVacuumGen L
  centralCharge_eq : centralCharge = (IntegralLattice.rank L : ℂ)
  charge_mem : ∀ x : L.carrier,
    LatGen.latChargeGen L x ∈ grading (L.form x x / 2)
  charge_leading : ∀ x y : L.carrier,
    mode (Y (LatGen.latChargeGen L x)) (-L.form x y - 1) (LatGen.latChargeGen L y) =
      IntegralLattice.latticeCocycle L x y • LatGen.latChargeGen L (x + y)
  charge_positive_modes : ∀ (x y : L.carrier) (n : ℤ), -L.form x y - 1 < n →
    mode (Y (LatGen.latChargeGen L x)) n (LatGen.latChargeGen L y) = 0

/-- The intended hyperbolic lattice state space, reusing the existing carrier. -/
noncomputable abbrev LorentzianFockSpace := LatGen.LatModuleGen lorentzianLattice

noncomputable def lorentzianChargeState (x : lorentzianLattice.carrier) :
    LorentzianFockSpace := LatGen.latChargeGen lorentzianLattice x

def LorentzianVertexTarget : Prop := Nonempty (EvenLatticeVertexRealization lorentzianLattice)

/-- Exact conversion to the repository's conformal vertex-algebra structure. -/
def UnboundedConformalVertexAlgebra.toNearVOA {M : Type*} [AddCommGroup M] [Module ℂ M]
    (V : UnboundedConformalVertexAlgebra M) : NearVOA M where
\end{leancode}
\noindent (The last declaration is in \leanpath{VOA/\allowbreak Moonshine/\allowbreak MoonshineCNComplexFrame.lean}.)

\begin{dictionary}
\textbf{A Lean alias.} The Lean file states the realisation as an extension of a structure
\chseventeenL{UnboundedConformalVertexAlgebra} (same file, not printed) rather than of \chseventeenL{NearVOA}.
Mathematically there is nothing new: it has literally the fields of \chseventeenL{NearVOA}, field for field, the only
difference being that \chseventeenL{grading_mode} and \chseventeenL{L0_grading} take their weights as explicit rather than
implicit arguments, and \chseventeenL{toNearVOA} copies each field (its body, not shown, is sixteen lines
\chseventeenL{grading := V.grading}, …).  It is a historical duplicate (the near-VOA structure had been overlooked
when the target was written); read it as ``near-VOA''.  \chseventeenL{LatGen.LatModuleGen L} $=\chseventeenF_r\otimes\C[L]$,
\chseventeenL{latVacuumGen} $=1\otimes e^0$, \chseventeenL{latChargeGen L x} $=1\otimes e^x$ (Chapter~\ref{ch:03},
Definition~\ref{ch03:def:LatModuleGen}).  \chseventeenL{L.form x x / 2} is integer division; with \chseventeenL{even} it is exact.
\textbf{Why the extra fields.} They pin the object to be the lattice vertex algebra and not an arbitrary
near-VOA on the same space: the central charge, the weights of the charged vacua, and the leading
operator products \eqref{ch17:eq:chargeleading} with the FLM cocycle.  Together with locality and the
creation axiom they determine $Y$ uniquely on the subalgebra generated by the $e^x$ (Goddard uniqueness,
Theorem~\ref{ch17:thm:goddard}), which for a lattice vertex algebra is everything once the Heisenberg
currents are also pinned; the realisation actually constructed below pins them too
(Theorem~\ref{ch17:thm:dualbasis}).
\textbf{A stale docstring.} The docstring of \chseventeenL{EvenLatticeVertexRealization} (``An inhabitant \dots\ is still
required'') and that of \chseventeenL{LorentzianVertexTarget} (``The exact remaining existence statement'') describe
the state of the development before the construction; the inhabitant is
\chseventeenL{lorentzianVertexTarget_inhabited} (Theorem~\ref{ch17:thm:realization}).  See \S\ref{ch17:sec:traps}.
\end{dictionary}

\begin{physics}
The target says: ``the chiral algebra of two compact light-cone bosons, normalised so that the vertex
operator of $e^x$ creates the momentum state $e^x$, and so that $e^x(z)e^y(w)\sim\varepsilon(x,y)(z-w)^{\langle x,y\rangle}e^{x+y}(w)$.''
These are the defining OPEs of a lattice CFT.
\end{physics}

\begin{status}
Definitions.  \chseventeenL{UnboundedConformalVertexAlgebra.toNearVOA} is printed in
\leanpath{VOA/\allowbreak Moonshine/\allowbreak MoonshineCNAudit2.lean}; standard axioms.
\end{status}

\begin{theorem}[No lower bound is possible]\label{ch17:thm:notbounded}
Let $(\mathcal V_n)_{n\in\Z}$ be \emph{any} family of subspaces of $\chseventeenVII$'s state space such that each
charged vacuum $e^x$ lies in $\mathcal V_{\Delta(x)}$.  Then there is no $N$ with $\mathcal V_n=0$ for all $n<N$.
In particular the weight grading of $\chseventeenVII$ is not bounded below.
\end{theorem}

\noindent{\raggedright Files \leanpath{VOA/\allowbreak Moonshine/\allowbreak MoonshineCNVertexTarget.lean},
\leanpath{VOA/\allowbreak Moonshine/\allowbreak MoonshineCNGrading.lean}.\par}
\begin{leancode}
theorem lorentzian_grading_not_bounded_below
    (grading : ℤ → Submodule ℂ LorentzianFockSpace)
    (charge_mem : ∀ x : lorentzianLattice.carrier,
      lorentzianChargeState x ∈ grading (lorentzianEnergy x)) :
    ¬ ∃ N : ℤ, ∀ n : ℤ, n < N → grading n = ⊥ := by

theorem lorentzianWeightGrading_not_bounded_below :
    ¬ ∃ N : ℤ, ∀ n : ℤ, n < N → lorentzianWeightGrading n = ⊥ :=
\end{leancode}

\begin{dictionary}
\chseventeenL{⊥} is the zero subspace.  The first theorem quantifies over \emph{every} candidate grading that puts
the charged vacua at their energies; it therefore rules out \emph{any} \chseventeenL{GradedVertexOperatorAlgebra}
structure with the physical weights, not only the one constructed.  The second is the instance for the
actual grading \chseventeenL{lorentzianWeightGrading} $=$ \chseventeenL{gradingLatGen lorentzianLattice}
(Definition~\ref{ch17:def:weightgrading}).  The only extra input is \chseventeenL{lorentzianChargeState_ne_zero}
(the coefficient of $1\otimes e^x$ in $e^x$ is $1$).
\end{dictionary}

\begin{proof}[Proof on paper]
Given $N$, take $x=(1,N-1)$; then $\Delta(x)=N-1<N$, so $e^x\in\mathcal V_{N-1}=0$, contradicting $e^x\neq0$.
\end{proof}

\begin{physics}
Lorentzian compactifications have no ground state: the energy $p^+p^-$ is unbounded below.  This is the
mathematical form of the statement that $\chseventeenVII$ is not the chiral algebra of a unitary compact CFT.
\end{physics}

\begin{status}
Proved unconditionally; both printed (\leanpath{VOA/\allowbreak Moonshine/\allowbreak MoonshineCNAudit.lean},
\leanpath{VOA/\allowbreak Moonshine/\allowbreak MoonshineCNAudit2.lean}); standard axioms.
\end{status}

\begin{definition}[The weight grading]\label{ch17:def:weightgrading}
For an integral lattice $L$ and $n\in\Z$, $(V_L)_n=\sum_\alpha(\chseventeenF_r)_{n-\lfloor\langle\alpha,\alpha\rangle/2\rfloor}\otimes e^\alpha$
(the floor is exact when $L$ is even, the only case used later)
(Chapter~\ref{ch:03}, Definition~\ref{ch03:def:gradingLatGen}, the library's \chseventeenL{gradingLatGen}).  For
$\chseventeenII$ this is \chseventeenL{lorentzianWeightGrading}.
\end{definition}

\noindent File \leanpath{VOA/\allowbreak Moonshine/\allowbreak MoonshineCNGrading.lean}.
\begin{leancode}
noncomputable def lorentzianWeightGrading : ℤ → Submodule ℂ LorentzianFockSpace :=
  gradingLatGen lorentzianLattice
\end{leancode}

\begin{dictionary}
Same grading as for positive-definite lattices; it is the \emph{theorems} about it that change.  The weight of
$p\otimes e^\alpha$ is (level of $p$) $+\,\Delta(\alpha)$; that this is the $L_0$-eigenvalue is proved in
Theorem~\ref{ch17:thm:virasoro}.
\end{dictionary}

\begin{physics}
Conformal weight $=$ oscillator level $+$ $\tfrac12$momentum$^2$, with the Lorentzian square.
\end{physics}

\begin{status}
Definition.
\end{status}

\begin{theorem}[Internality without positivity; infinite weight spaces]\label{ch17:thm:grading}\label{ch17:thm:weightzero}
(i) For \emph{every} integral lattice $L$, $V_L=\bigoplus_{n\in\Z}(V_L)_n$ is an internal direct sum.
(ii) For $\chseventeenII$, $\chseventeenL{halfNorm}(\alpha)=\Delta(\alpha)$ and $e^\alpha\in(V)_{\Delta(\alpha)}$.
(iii) The states $e^{(i,0)}$, $i\in\Z$, are linearly independent; hence $(\chseventeenVII)_0$ is
infinite-dimensional.
\end{theorem}

\noindent File \leanpath{VOA/\allowbreak Moonshine/\allowbreak MoonshineCNGrading.lean}.
\begin{leancode}
/-- Internal integer grading, including indefinite lattices and infinite fibers. -/
theorem latticeWeightGrading_isInternal (L : IntegralLattice) :
    DirectSum.IsInternal (gradingLatGen L) :=

theorem lorentzian_halfNorm (α : lorentzianLattice.carrier) :
    halfNorm lorentzianLattice α = lorentzianEnergy α := by

theorem lorentzianChargeState_mem_weight (α : lorentzianLattice.carrier) :
    lorentzianChargeState α ∈ lorentzianWeightGrading (lorentzianEnergy α) := by

/-- Infinitely many linearly independent isotropic charged ground states. -/
theorem lorentzianIsotropicChargeStates_independent :
    LinearIndependent ℂ (fun i : ℤ => lorentzianChargeState ![i, 0]) := by

/-- This is an actual obstruction to the finite-weight-space `VertexAlgebra`. -/
theorem lorentzianWeightZero_not_finiteDimensional :
    ¬ FiniteDimensional ℂ (lorentzianWeightGrading 0) := by
\end{leancode}

\begin{dictionary}
(i) has no hypothesis at all, not even evenness: the library already had internality under positivity
(\chseventeenL{gradingLatGen_isInternal_of_posDef}); the new proof works for infinite fibres.
\chseventeenL{![i, 0]} is the vector $(i,0)$.  \chseventeenL{LinearIndependent ℂ f} is Mathlib's linear independence of a
family.  \chseventeenL{FiniteDimensional ℂ (lorentzianWeightGrading 0)} refers to the subspace as a vector space in
its own right.  \textbf{Fidelity:} (iii) is precisely the obstruction to the library's
\chseventeenL{PreVertexAlgebra}/\chseventeenL{VertexOperatorAlgebra}, whose field \chseventeenL{finrank} requires every weight space
to be finite-dimensional.
\end{dictionary}

\begin{proof}[Proof on paper]
(i) $V_L=\bigoplus_{(w,\alpha)}(\chseventeenF_r)_w\otimes e^\alpha$ is a direct sum of the bigraded pieces; regrouping
the summands of a direct sum along the map $(w,\alpha)\mapsto w+\langle\alpha,\alpha\rangle/2$ is still a direct sum,
whatever the sizes of the fibres.  (ii) $2\Delta(\alpha)=\langle\alpha,\alpha\rangle$ and $1$ has level $0$.  (iii) The
linear functional ``coefficient of $1\otimes e^{(i,0)}$'' is $1$ on $e^{(i,0)}$ and $0$ on $e^{(j,0)}$, $j\ne i$;
each $(i,0)$ is isotropic, so all these states lie in weight $0$.
\end{proof}

\begin{physics}
Infinitely many weight-zero states: every isotropic (lightlike) momentum gives a massless ground state.  The
$e^{(i,0)}$, $i\in\Z$, pairwise pair to $\langle(i,0),(j,0)\rangle=0$, so their vertex operators commute and multiply
like a group algebra $\C[\Z]$ inside weight zero; this is a direct consequence of $\chseventeenII$ containing isotropic
vectors.  In
string language these are the zero-energy light-cone momenta; the infinite degeneracy is why a character
$\operatorname{tr}q^{L_0}$ of $\chseventeenVII$ does not exist, and why the theory must be graded by momentum
\emph{and} weight (Theorem~\ref{ch17:thm:momentum}) to have finite pieces.
\end{physics}

\begin{status}
Proved unconditionally; all printed in \leanpath{VOA/\allowbreak Moonshine/\allowbreak MoonshineCNAudit2.lean}; standard axioms.
\end{status}

%======================================================================
\section{The construction of the fields}\label{ch17:sec:construction}
%======================================================================

The construction follows the positive-definite construction of Chapter~\ref{ch:03}
(\S\S\ref{ch03:sec:vertexops}--\ref{ch03:sec:reconstruction}) step by step, with the real frame replaced by
an arbitrary complex frame $F$ (Definition~\ref{ch17:def:frame}).  Each file states which positive-definite
file it adapts (\chseventeenL{LatticeGenVertex}, \chseventeenL{LatticeGenSameSignLocal}, …); those originals are unchanged.
We present the definitions and the theorems that enter the final structure; the many intermediate
contraction lemmas are described collectively at the end of the section.  Throughout, $L$ is an integral
lattice of rank $r$ with a complex frame $F$, and $\alpha(n)$ denotes the frame oscillator.

\subsection{The exponentials}

\begin{definition}[Creation and annihilation exponentials]\label{ch17:def:exponentials}
For $\alpha\in L$ put $x_m(\alpha)=\sum_a\alpha^\perp_ax_{m,a}$ (the polynomial of $\alpha(-m-1)\chseventeenvac$), and define
polynomials $P_N(\alpha)$ and operators $A_N(\alpha)$ recursively by $P_0=1$, $A_0=\mathrm{id}$ and
\begin{equation*}
  (N+1)\,P_{N+1}(\alpha)=\sum_{m=0}^{N}x_m(\alpha)\,P_{N-m}(\alpha),\qquad
  (N+1)\,A_{N+1}(\alpha)=\sum_{m=0}^{N}\bigl(-\alpha(m+1)\bigr)A_{N-m}(\alpha).
\end{equation*}
Then $E^-(\alpha,z)=\sum_NP_N(\alpha)z^N=\exp\bigl(\sum_{n\ge1}\alpha(-n)z^n/n\bigr)\chseventeenvac$ and
$E^+(\alpha,z)=\sum_NA_N(\alpha)z^{-N}=\exp\bigl(-\sum_{n\ge1}\alpha(n)z^{-n}/n\bigr)$.
\end{definition}

\noindent File \leanpath{VOA/\allowbreak Moonshine/\allowbreak MoonshineCNExponential.lean}.
\begin{leancode}
noncomputable def creationVariable (α : L.carrier) (m : ℕ) : HeisenbergN.FockN (rank L) :=
  ∑ a : Fin (rank L), F.coordinate α a • (MvPolynomial.X ((m, a)) : HeisenbergN.FockN (rank L))

noncomputable def creationCoefficient (α : L.carrier) : ℕ → HeisenbergN.FockN (rank L)
  | 0 => 1
  | (N+1) => ((N+1 : ℂ))⁻¹ • ∑ m ∈ Finset.range (N+1), creationVariable F α m * creationCoefficient α (N - m)

noncomputable def annihilationCoefficient (α : L.carrier) : ℕ → End ℂ (HeisenbergN.FockN (rank L))
  | 0 => LinearMap.id
  | (N+1) => ((N+1 : ℂ))⁻¹ • ∑ m ∈ Finset.range (N+1),
      (-(frameOscillator F α (m+1 : ℤ))).comp (annihilationCoefficient α (N - m))
\end{leancode}

\begin{dictionary}
\chseventeenL{MvPolynomial.X ((m, a))} is the variable $x_{m,a}$, i.e.\ $h_a(-m-1)\chseventeenvac$.  The recursions are
the Newton identities $\partial_zE=E\,\partial_z(\text{exponent})$ read coefficientwise; using them, rather
than exponentials of infinite sums, keeps every coefficient a finite polynomial or a finite composite of
operators.  The frame $F$ and the lattice $L$ are section variables.  \textbf{Fidelity:} these are exactly
the coefficient expansions of FLM's exponentials $E^\mp(-\alpha,z)$ (FLM, \S4.1 and \S8.4; FLM write the vertex operator with $-\alpha$ in the exponentials) in the convention of
Chapter~\ref{ch:03}, Definition~\ref{ch03:def:vSchur}; no convergence or completed tensor product is used.
\end{dictionary}

\begin{physics}
$E^-(\alpha,z)E^+(\alpha,z)$ is the oscillator part ${:}e^{i\alpha\cdot X(z)}{:}$ of the vertex operator, with
the creation half on the left (normal ordering).  On a fixed vector $E^+$ acts by a finite sum, while $E^-$
has infinitely many terms: the familiar fact that a vertex operator creates states at all oscillator levels.
\end{physics}

\begin{status}
Definitions; all lemmas of the file printed in \leanpath{VOA/\allowbreak Moonshine/\allowbreak MoonshineCNAudit2.lean}.
\end{status}

\subsection{The charged fields}

\begin{definition}[Charged field]\label{ch17:def:charged}
Let $\varepsilon:L\times L\to\C$.  The \emph{dressing} field is the normal-ordered product
$\mathcal E(\alpha,z)=E^-(\alpha,z)E^+(\alpha,z)$ on $\chseventeenF_r$.  The \emph{charged field} of $\alpha$ is the
formal distribution on $V_L$ with
\begin{equation}\label{ch17:eq:chargedfield}
  Y_\varepsilon(e^\alpha,z)\,(p\otimes e^\beta) = \varepsilon(\alpha,\beta)\,z^{\langle\alpha,\beta\rangle}\,
  \mathcal E(\alpha,z)p\otimes e^{\alpha+\beta},
\end{equation}
i.e.\ coefficient of $z^t$ equal to $\varepsilon(\alpha,\beta)\,\mathcal E_{t-\langle\alpha,\beta\rangle}(\alpha)p\otimes e^{\alpha+\beta}$.
For $\chseventeenII$, \chseventeenL{lorentzianChargedField} uses $\varepsilon_{\mathrm{std}}$ and
\chseventeenL{lorentzianTargetChargedField} (equal to \chseventeenL{canonicalChargedField lorentzianComplexFrame}) uses the
FLM cocycle.
\end{definition}

\noindent{\raggedright Files \leanpath{VOA/\allowbreak Moonshine/\allowbreak MoonshineCNChargedFields.lean},
\leanpath{VOA/\allowbreak Moonshine/\allowbreak MoonshineCNCurrents.lean}.\par}
\begin{leancode}
noncomputable def dressingField (α : L.carrier) : FormalDist (End ℂ (FockN (rank L))) :=
  nprod (creationField_isField F α) (annihilationField_isField F α)

theorem chargedField_tmul (α β : L.carrier) (t : ℤ) (p : FockN (rank L)) :
    chargedField F ε α t (latStateGen L p β)
      = ε α β • latStateGen L (dressingField F α (t - L.form α β) p) (α + β) := by

/-- The leading product of two charged ground states, with its actual cocycle. -/
theorem chargedField_charge_leading (α β : L.carrier) :
    mode (chargedField F ε α) (-L.form α β - 1) (latChargeGen L β) =
      ε α β • latChargeGen L (α + β) := by

noncomputable abbrev canonicalChargedField (α : L.carrier) :=
  chargedField F (latticeCocycle L) α

/-- Charged fields using the fixed isotropic-coordinate cocycle. -/
noncomputable def lorentzianChargedField (α : lorentzianLattice.carrier) :
    FormalDist (End ℂ LorentzianFockSpace) :=
  chargedField lorentzianComplexFrame standardLorentzianCocycle α

/-- The same construction with the chosen-basis cocycle required by the old target. -/
noncomputable def lorentzianTargetChargedField (α : lorentzianLattice.carrier) :
    FormalDist (End ℂ LorentzianFockSpace) :=
  chargedField lorentzianComplexFrame (latticeCocycle lorentzianLattice) α
\end{leancode}

\begin{dictionary}
\chseventeenL{nprod} is the normal-ordered product of two fields (Chapter~\ref{ch:L03}, \S\ref{l03:sec:vafacts});
\chseventeenL{creationField F α t} is $P_t(\alpha)$ times (multiplication) for $t\ge0$ and $0$ for $t<0$, and
\chseventeenL{annihilationField F α t} is $A_{-t}(\alpha)$ for $t\le0$ and $0$ for $t>0$ (same file).  The definition of
\chseventeenL{chargedField} itself (not shown) is the linear extension of \chseventeenL{chargedField_tmul} over the basis
$p\otimes e^\beta$, built with \chseventeenL{TensorProduct.lift} and \chseventeenL{Finsupp.lsum}; \chseventeenL{chargedField_tmul} is
its defining equation, which is \eqref{ch17:eq:chargedfield} coefficientwise ($t\mapsto t-\langle\alpha,\beta\rangle$ is
the factor $z^{\langle\alpha,\beta\rangle}$, i.e.\ $z^{\alpha(0)}$).  \chseventeenL{chargedField_charge_leading} is
\eqref{ch17:eq:chargeleading}: the mode $-\langle\alpha,\beta\rangle-1$ picks the coefficient $\mathcal E_0=\mathrm{id}$ on the
vacuum.  The companion \chseventeenL{chargedField_charge_higher_modes} gives the vanishing of higher modes.
\textbf{Two cocycles:} the vertex algebra uses \chseventeenL{lorentzianTargetChargedField}.
\end{dictionary}

\begin{physics}
\eqref{ch17:eq:chargedfield} is the FLM/Frenkel--Kac vertex operator (sign convention as in
Definition~\ref{ch17:def:exponentials}): shift momentum by $\alpha$ (the operator $e_\alpha$ with Klein
factor $\varepsilon$), multiply by $z^{\alpha\cdot p}$, and dress with the oscillator exponential.  It is the operator
${:}e^{i\alpha\cdot X(z)}{:}$ of a compact boson, here with a Lorentzian $\alpha\cdot X$.
\end{physics}

\begin{status}
Definitions and proved identities; printed in \leanpath{VOA/\allowbreak Moonshine/\allowbreak MoonshineCNAudit2.lean}; standard axioms.
\end{status}

\begin{theorem}[Locality of charged fields]\label{ch17:thm:chargedlocal}
Let $C:L\times L\to\C$ be bimultiplicative in the second variable with $C(\alpha,\beta)=(-1)^{\langle\alpha,\beta\rangle}C(\beta,\alpha)$.
Then for all $\alpha,\beta\in L$ the charged fields $Y_C(e^\alpha,z)$, $Y_C(e^\beta,z)$ are mutually local.  Both the
FLM cocycle and $\varepsilon_{\mathrm{std}}$ satisfy the hypothesis (both are bimultiplicative, $\pm1$-valued, with
commutator $(-1)^{\langle\alpha,\beta\rangle}$); hence both families of Lorentzian charged
fields are mutually local.
\end{theorem}

\noindent File \leanpath{VOA/\allowbreak Moonshine/\allowbreak MoonshineCNChargedLocality.lean}.
\begin{leancode}
/-- Precisely the scalar identities used to cancel the bosonic reordering sign. -/
structure ChargeCommutationFactors (L : IntegralLattice) where
  value : L.carrier → L.carrier → ℂ
  add_right : ∀ α β γ, value α (β + γ) = value α β * value α γ
  swap : ∀ α β, value α β = (-1 : ℂ) ^ L.form α β * value β α

theorem chargedField_local  (α β : L.carrier) :
    Local (chargedField F C.value α) (chargedField F C.value β) := by

theorem lorentzianChargedField_local (α β : lorentzianLattice.carrier) :
    Local (lorentzianChargedField α) (lorentzianChargedField β) :=

/-- Mutual locality for the cocycle prescribed by `LorentzianVertexTarget`. -/
theorem lorentzianTargetChargedField_local (α β : lorentzianLattice.carrier) :
    Local (lorentzianTargetChargedField α) (lorentzianTargetChargedField β) :=
\end{leancode}

\begin{dictionary}
In \chseventeenL{chargedField_local}, $F$ and \chseventeenL{C : ChargeCommutationFactors L} are section variables, so the
statement is: for every frame and every factor system $C$.  The two fields \chseventeenL{add_right}, \chseventeenL{swap} are
exactly what the computation uses; \chseventeenL{canonicalChargeCommutationFactors L hEven} supplies them for the FLM
cocycle (using evenness and $\varepsilon(\alpha,\alpha)^2=1$), \chseventeenL{lorentzianChargeCommutationFactors} for
$\varepsilon_{\mathrm{std}}$.  No positivity and no sign condition on $\langle\alpha,\beta\rangle$ is assumed.
\end{dictionary}

\begin{proof}[Proof on paper]
Moving $E^+(\alpha,z_1)$ past $E^-(\beta,z_2)$ produces the scalar $(1-z_2/z_1)^{\langle\alpha,\beta\rangle}$, and the
momentum factors give $z_1^{\langle\alpha,\beta\rangle}$; so
$Y(e^\alpha,z_1)Y(e^\beta,z_2)=C(\alpha,\beta)(z_1-z_2)^{\langle\alpha,\beta\rangle}\,{:}\cdots{:}$, where the normal-ordered
product ${:}\cdots{:}$ is symmetric in the two fields, and $(z_1-z_2)^k$ is expanded in $z_2/z_1$.  If
$k=\langle\alpha,\beta\rangle\ge0$ this is a polynomial, and the swap identity $C(\alpha,\beta)(z_1-z_2)^k=C(\beta,\alpha)(z_2-z_1)^k$
shows that the two orderings are \emph{equal}: the fields commute coefficientwise.  If $k<0$ the two
orderings are the expansions of the same rational function in the two regions, and multiplying by
$(z_1-z_2)^{|k|}$ makes both equal.  The Lean proof implements the negative case by finite binomial
differences rather than by rational functions, so that every sum is finite on a fixed vector.
\end{proof}

\begin{physics}
$e^\alpha(z)e^\beta(w)\sim(z-w)^{\langle\alpha,\beta\rangle}e^{\alpha+\beta}(w)$: the OPE has a pole of order $-\langle\alpha,\beta\rangle$ when
the pairing is negative and is regular when it is nonnegative.  In a positive-definite lattice a fixed $\alpha$ has negative pairing with $\beta$ only
within a bounded range; in $\chseventeenII$ the pairing $\langle(1,0),(0,-n)\rangle=-n$ is unbounded below, so pole orders are
unbounded.  The proof handles each pair separately, so this causes no difficulty, but no argument may use a uniform
bound.
\end{physics}

\begin{status}
Proved unconditionally (for every frame and factor system); all three printed in
\leanpath{VOA/\allowbreak Moonshine/\allowbreak MoonshineCNAudit2.lean}; standard axioms.
\end{status}

\subsection{Currents, translation, reconstruction}

\begin{definition}[Currents and translation operator]\label{ch17:def:currents}
For a frame colour $a$, the \emph{current} $h_a(z)$ on $V_L$ is the Fock current $\sum_nh_a(n)z^{-n-1}$
acting on the first factor, plus the \emph{momentum zero mode} $h_a(0)(p\otimes e^\beta)=\beta^\perp_a\,p\otimes e^\beta$.
The derivative currents are $\partial^ih_a(z)/i!$.  The translation operator is
$T=L^{\mathrm{Fock}}_{-1}\otimes1+\sum_ah_a(-1)\,h_a(0)$.
\end{definition}

\noindent{\raggedright Files \leanpath{VOA/\allowbreak Moonshine/\allowbreak MoonshineCNCurrents.lean},
\leanpath{VOA/\allowbreak Moonshine/\allowbreak MoonshineCNTranslation.lean}.\par}
\begin{leancode}
noncomputable def frameMomentum (a : Fin (rank L)) : LatModuleGen L →ₗ[ℂ] LatModuleGen L :=
  LinearMap.lTensor (FockN (rank L)) (groupMomentum F a)

noncomputable def frameCurrent (a : Fin (rank L)) : FormalDist (End ℂ (LatModuleGen L)) :=
  fun j => LinearMap.rTensor (L.carrier →₀ ℂ) (currentN (rank L) 1 a j)
    + (if j = -1 then frameMomentum F a else 0)

noncomputable def frameDerivativeCurrent (a : Fin (rank L)) (i : ℕ) : FormalDist (End ℂ (LatModuleGen L)) :=
  (i.factorial : ℂ)⁻¹ • (VOA.deriv^[i] (frameCurrent F a))

noncomputable def frameTranslation : End ℂ (LatModuleGen L) :=
  LinearMap.rTensor (L.carrier →₀ ℂ) (HeisenbergN.LvirN (rank L) 1 (-1)) + momentumCoupling F
\end{leancode}

\begin{dictionary}
\chseventeenL{groupMomentum F a} multiplies $e^\beta$ by $\beta^\perp_a$ (lemma \chseventeenL{frameMomentum_latStateGen}).  The
coefficient index $j=-1$ is the coefficient of $z^{-1}$, i.e.\ the mode $h_a(0)$: so the momentum enters
exactly as the zero mode, as it must.  \chseventeenL{rTensor}/\chseventeenL{lTensor} act on the first/second tensor factor.
\chseventeenL{LvirN (rank L) 1 (-1)} is the Fock-space $L_{-1}$ (Chapter~\ref{ch:03}); \chseventeenL{momentumCoupling F}
$=\sum_a(h_a(-1)\otimes1)\circ$\chseventeenL{frameMomentum F a}.  \chseventeenL{VOA.deriv^[i]} is the $i$-fold formal derivative.
\end{dictionary}

\begin{physics}
$h_a(z)=i\partial\phi_a(z)$ is the $a$-th free-boson current, whose zero mode is the momentum; $T=L_{-1}$ is
the generator of translations, and the cross term $h_a(-1)p_a$ is the momentum contribution
$\alpha(-1)\alpha(0)$ to $L_{-1}=\frac12\sum{:}\alpha\alpha{:}_{-1}$.
\end{physics}

\begin{status}
Definitions; printed in \leanpath{VOA/\allowbreak Moonshine/\allowbreak MoonshineCNAudit2.lean}.
\end{status}

\begin{theorem}[Reconstruction on the lattice state space]\label{ch17:thm:reconstruct}
Let $L$ be even, with a complex frame $F$.  Define the \emph{word states}
$h_{a_1}(-i_1-1)\cdots h_{a_k}(-i_k-1)e^\alpha$ and the \emph{word fields}
${:}\tfrac{\partial^{i_1}h_{a_1}}{i_1!}\cdots\tfrac{\partial^{i_k}h_{a_k}}{i_k!}Y(e^\alpha,z){:}$ (with the FLM-cocycle charged
field).  There is a linear map $Y:V_L\to\End(V_L)[[z^{\pm1}]]$ sending every word state to its word field,
such that every $Y(u,z)$ is a field, all are mutually local, $Y(u,z)\chseventeenvac=u+O(z)$, and every $Y(u,z)$ is
$T$-covariant.  The library's \chseventeenL{frameStateField F hEven} is such a map, chosen once and for all.
\end{theorem}

\noindent{\raggedright Files \leanpath{VOA/\allowbreak Moonshine/\allowbreak MoonshineCNReconstruction.lean},
\leanpath{VOA/\allowbreak Moonshine/\allowbreak MoonshineCNStateFields.lean}.\par}
\begin{leancode}
noncomputable def frameWordState (L : IntegralLattice) : List (ℕ × Fin (rank L)) → L.carrier → LatModuleGen L
  | [], α => latChargeGen L α
  | (ia :: rest), α =>
      LinearMap.rTensor (L.carrier →₀ ℂ) (currentN (rank L) 1 ia.2 ia.1) (frameWordState L rest α)

theorem frameStateField_exists (hEven : L.Even)  :
    ∃ Yfull : LatModuleGen L →ₗ[ℂ] FormalDist (End ℂ (LatModuleGen L)),
      (∀ (lα : List (ℕ × Fin (rank L)) × L.carrier),
        Yfull (Function.uncurry (frameWordState L) lα)
          = Function.uncurry (fun l α => frameWordField F l α) lα) ∧
      (∀ u : LatModuleGen L, IsField (Yfull u)) ∧
      (∀ u v : LatModuleGen L, Local (Yfull u) (Yfull v)) ∧
      (∀ u : LatModuleGen L, (Yfull u) 0 (latVacuumGen L) = u) ∧
      (∀ u : LatModuleGen L, ∀ q : ℤ, q < 0 → (Yfull u) q (latVacuumGen L) = 0) ∧
      (∀ u : LatModuleGen L, Reconstruction.DCovariant (frameTranslation F) (Yfull u)) :=

noncomputable def frameStateField (hEven : L.Even)  :
    LatModuleGen L →ₗ[ℂ] FormalDist (End ℂ (LatModuleGen L)) :=
  (frameStateField_exists F hEven).choose
\end{leancode}

\begin{dictionary}
A word is a list of pairs $(i,a)$; \chseventeenL{currentN (rank L) 1 a i} applied as the coefficient of $z^i$ is the mode
$h_a(-i-1)$.  \chseventeenL{frameWordField F l α} is defined recursively in the same file as the normal-ordered
product of \chseventeenL{frameDerivativeCurrent F a i} with the word field of the tail, starting from
\chseventeenL{canonicalChargedField F α}.  \chseventeenL{IsField} means lower truncation on each vector;
\chseventeenL{(Yfull u) 0} is the $z^0$ coefficient, so \chseventeenL{(Yfull u) 0 (latVacuumGen L) = u} and the vanishing for
$q<0$ say $Y(u,z)\chseventeenvac=u+O(z)$.  \chseventeenL{DCovariant T A} is $[T,A(z)]=\partial_zA(z)$ (Chapter~\ref{ch:03},
Definition~\ref{ch03:def:DCovariant}).  \chseventeenL{.choose} picks a witness; it is unique by Goddard's theorem
(Theorem~\ref{ch17:thm:goddard}), so the choice is harmless.  The existence is the library's reconstruction
theorem \chseventeenL{Reconstruction.reconstruct} (Chapter~\ref{ch:L03}, Theorem~\ref{l03:thm:reconstruct}), applied
to the generating fields; its hypotheses are discharged by \chseventeenL{frameWordField_local} (locality of all word
fields, by Dong's lemma from Theorem~\ref{ch17:thm:chargedlocal} and current--current/current--charge
locality), the vacuum properties and $T$-covariance.
\end{dictionary}

\begin{proof}[Proof on paper]
The word states span $V_L$ (they are the monomial basis).  The generating fields $h_a(z)$, $Y(e^\alpha,z)$ are
mutually local, $T$-covariant and create their states from the vacuum; by Dong's lemma their normal-ordered
products and derivatives remain mutually local; the reconstruction theorem then extends $Y$ linearly from the
spanning set, the point being that the extension is well defined because a local, $T$-covariant field is
determined by the state it creates.
\end{proof}

\begin{physics}
The state--field correspondence of a free-boson CFT: the field of an oscillator state is the normal-ordered
product of derivatives of currents with the vertex operator, ${:}\partial^{i_1}\phi\cdots e^{i\alpha\cdot\phi}{:}$.
\end{physics}

\begin{status}
Proved for every even lattice and every complex frame; \chseventeenL{frameStateField_exists} and the word-field
lemmas printed in \leanpath{VOA/\allowbreak Moonshine/\allowbreak MoonshineCNAudit2.lean}; standard axioms.
\end{status}

\subsection{The conformal vector and the Virasoro algebra}

\begin{definition}[Conformal vector and Sugawara field]\label{ch17:def:omega}
$\omega_L=\tfrac12\sum_ah_a(-1)^2\chseventeenvac\otimes e^0$ (orthonormal frame colours), and the Sugawara field
$\tfrac12\sum_a{:}h_a(z)h_a(z){:}$ on $V_L$, whose modes are $L_n$.
\end{definition}

\noindent{\raggedright Files \leanpath{VOA/\allowbreak Moonshine/\allowbreak MoonshineCNConformal.lean},
\leanpath{VOA/\allowbreak Moonshine/\allowbreak MoonshineCNVirasoro.lean}.\par}
\begin{leancode}
noncomputable def latticeConformalVector (L : IntegralLattice) : LatModuleGen L := mkEα L 0 (HeisenbergN.ωN (rank L) 1)

noncomputable def frameSugawaraField : FormalDist (End ℂ (LatModuleGen L)) :=
  (1 / 2 : ℂ) • ∑ a : Fin (rank L), nprod (frameCurrent_isField F a) (frameCurrent_isField F a)

noncomputable def frameVirasoroMode (n : ℤ) : End ℂ (LatModuleGen L) := frameSugawaraField F (-n - 2)
\end{leancode}

\begin{dictionary}
\chseventeenL{ωN (rank L) 1} is the Fock-space Sugawara vector $\tfrac12\sum_ah_a(-1)^2\chseventeenvac$ (Chapter~\ref{ch:03},
\eqref{ch03:eq:omegaHeis}), placed in momentum $0$ by \chseventeenL{mkEα L 0}.  The coefficient of $z^{-n-2}$ of the
Sugawara field is $L_n$; hence \chseventeenL{frameVirasoroMode n} $=L_n$.  Note the conformal vector does not depend on
the frame, while the currents' zero modes do (through \chseventeenL{frameMomentum}).  In the isotropic basis the same
vector is $\tfrac12(u(-1)v(-1)+v(-1)u(-1))\chseventeenvac$ (Theorem~\ref{ch17:thm:dualbasis}).
\end{dictionary}

\begin{physics}
$T(z)=\frac12\sum_a{:}\partial\phi_a\partial\phi_a{:}$, the free-boson stress tensor; in light-cone
coordinates $T=-{:}\partial X^+\partial X^-{:}$ up to normalisation.
\end{physics}

\begin{status}
Definitions; printed in \leanpath{VOA/\allowbreak Moonshine/\allowbreak MoonshineCNAudit2.lean}.
\end{status}

\begin{theorem}[Virasoro, $L_0$, translation, grading]\label{ch17:thm:virasoro}
Let $L$ be even with a complex frame.  Then, for $Y=$ \chseventeenL{frameStateField F hEven} and $\omega=\omega_L$:
$[L_m,L_n]=(m-n)L_{m+n}+\tfrac{\operatorname{rank}L}{12}(m^3-m)\delta_{m+n,0}$; $L_0=n$ on $(V_L)_n$;
$Y(L_{-1}v,z)=\partial_zY(v,z)$; and $u_{(j)}(V_L)_q\subseteq(V_L)_{p+q-j-1}$ for $u\in(V_L)_p$.
\end{theorem}

\noindent{\raggedright Files \leanpath{VOA/\allowbreak Moonshine/\allowbreak MoonshineCNVirasoro.lean},
\leanpath{VOA/\allowbreak Moonshine/\allowbreak MoonshineCNConformalTranslation.lean},
\leanpath{VOA/\allowbreak Moonshine/\allowbreak MoonshineCNGradedFields.lean}.\par}
\begin{leancode}
theorem frameStateField_virasoro (hEven : L.Even)  (m n : ℤ) :
    ⁅mode (frameStateField F hEven (latticeConformalVector L)) (m + 1),
      mode (frameStateField F hEven (latticeConformalVector L)) (n + 1)⁆
      = ((m : ℂ) - n) • mode (frameStateField F hEven (latticeConformalVector L)) (m + n + 1)
        + (if m + n = 0 then (((m : ℂ) ^ 3 - m) / 12 * (rank L : ℂ)) • (1 : End ℂ (LatModuleGen L))
           else 0) := by

theorem frameStateField_L0_grading (hEven : L.Even)  {n : ℤ} (v : LatModuleGen L)
    (hv : v ∈ gradingLatGen L n) :
    mode (frameStateField F hEven (latticeConformalVector L)) 1 v = (n : ℂ) • v := by

theorem frameStateField_translation_omega (hEven : L.Even)  (v : LatModuleGen L) :
    frameStateField F hEven (mode (frameStateField F hEven (latticeConformalVector L)) 0 v)
      = deriv (frameStateField F hEven v) := by

theorem frameStateField_grading_mode (hEven : L.Even)  {p q : ℤ} (u : LatModuleGen L)
    (hu : u ∈ gradingLatGen L p) (v : LatModuleGen L) (hv : v ∈ gradingLatGen L q) (j : ℤ) :
    mode (frameStateField F hEven u) j v ∈ gradingLatGen L (p + q - j - 1) := by
\end{leancode}

\begin{dictionary}
These are verbatim the fields \chseventeenL{virasoro}, \chseventeenL{L0_grading}, \chseventeenL{translation}, \chseventeenL{grading_mode} of
\chseventeenL{NearVOA} for the data $(V_L,$ \chseventeenL{frameStateField}, $\omega_L$, $c=\operatorname{rank}L)$.  The central
charge is \chseventeenL{(rank L : ℂ)}, the rank of the lattice, independently of its signature.
\end{dictionary}

\begin{proof}[Proof on paper]
The key computation (in \leanpath{VOA/\allowbreak Moonshine/\allowbreak MoonshineCNSugawara.lean}) is that $Y(\omega_L,z)$ equals the
Sugawara field: the reconstructed field of the word state $\frac12\sum h_a(-1)^2\chseventeenvac$ is
$\frac12\sum{:}h_ah_a{:}$.  On a momentum sector $\chseventeenF_r\otimes e^\beta$ the Sugawara modes are those of $r$
free bosons whose zero modes are the scalars $\beta^\perp_a$; the free-boson Sugawara computation gives the Virasoro
algebra with $c=r$ \emph{for any values of the zero modes}, so the frame and the signature do not enter.
$L_0=$ level $+\frac12\sum_a(\beta^\perp_a)^2=$ level $+\frac12\langle\beta,\beta\rangle$ by \eqref{ch17:eq:framepairing}, which is
the grading.  Translation: $L_{-1}$ coincides with the operator $T$ of Definition~\ref{ch17:def:currents}, and
every reconstructed field is $T$-covariant, which with the creation property gives $Y(Tv,z)=\partial Y(v,z)$.
The grading property is proved for word fields (each oscillator mode and each charged-field coefficient has a
definite weight) and extended by linearity.
\end{proof}

\begin{physics}
The two light-cone bosons have $c=2$; the Lorentzian signature does not change the central charge, which
counts bosons, not signs.  $L_0=N+\frac12p^2$ is the mass-shell operator of string theory.
\end{physics}

\begin{status}
Proved for every even lattice with a complex frame; printed in
\leanpath{VOA/\allowbreak Moonshine/\allowbreak MoonshineCNAudit2.lean}; standard axioms.
\end{status}

\subsection{The technical files of the construction, collectively}

The construction files not singled out above prove the intermediate identities; they are adaptations of the
positive-definite files named in their docstrings, with the real frame replaced by \chseventeenL{ComplexLatticeFrame}
and every contraction using the frame pairing \eqref{ch17:eq:framepairing}.  None introduces a hypothesis.

\begin{center}
\small
\begin{longtable}{p{0.31\textwidth}p{0.62\textwidth}}
\toprule
file (\texttt{MoonshineCN}\dots) & what it proves, and why it is needed \\ \midrule
\endhead
\texttt{HalfCommutation} & reordering $E^+(\alpha,z_1)E^-(\beta,z_2)$ by one factor, the normally ordered four-field product, $\Z$-indexed annihilation coefficients (adapts \texttt{LatticeGenLocality}, \texttt{LatticeGenSameSign}) \\
\texttt{PositiveContraction} & the contraction sum for $\langle\alpha,\beta\rangle\ge0$ and the resulting coefficientwise commutation (adapts \texttt{LatticeGenContraction}, \texttt{LatticeGenEEMaster}) \\
\texttt{NegativeContraction}, \texttt{NegativeExpansion}, \texttt{NegativeCollapse} & the general contraction sum and the finite binomial-difference argument for $\langle\alpha,\beta\rangle<0$, expanded and then collapsed to the locality identity (adapt \texttt{LatticeGenOppSign}, \texttt{LatticeGenOppExpand}, \texttt{LatticeGenOppCollapse}) \\
\texttt{ChargedLocality} & Theorem~\ref{ch17:thm:chargedlocal} \\
\texttt{Currents}, \texttt{CurrentCharge} & currents with momentum zero modes, their mutual locality, and $[h_a(m),Y(e^\alpha,z)]=\alpha^\perp_az^mY(e^\alpha,z)$ (adapt \texttt{LatticeGenBosonFields}, \texttt{LatticeGenBosonCharge}) \\
\texttt{Translation}, \texttt{CreationTranslation}, \texttt{ChargeTranslation} & $T$-covariance of currents, of the creation exponential, and of charged fields (adapt \texttt{LatticeGenTranslation}, \texttt{LatticeGenVirasoroEminus}, \texttt{LatticeGenVirasoroCharge}) \\
\texttt{Reconstruction}, \texttt{StateFields} & word fields, their locality, Theorem~\ref{ch17:thm:reconstruct} and the near-VOA field axioms of \chseventeenL{frameStateField} (adapt \texttt{LatticeGenWordFields}, \texttt{LatticeGenVOAInstance}) \\
\texttt{Conformal}, \texttt{Sugawara}, \texttt{Virasoro} & $Y(\omega_L,z)=$ Sugawara field, its mode commutators with currents and charged fields, and the Virasoro relation (adapt \texttt{LatticeGenConformal}, \texttt{LatticeGenSugawaraTmul}, \texttt{LatticeGenVirasoro}) \\
\texttt{ConformalTranslation}, \texttt{GradedFields} & $L_{-1}=T$, $L_0=$ weight, and the weight shift of every mode (adapt \texttt{LatticeGenL0Translation}, \texttt{LatticeGenGradedField}) \\
\bottomrule
\end{longtable}
\end{center}

%======================================================================
\section{The vertex algebra of the hyperbolic plane}\label{ch17:sec:realization}
%======================================================================

\begin{theorem}[Lattice vertex algebras of indefinite lattices; $\chseventeenVII$]\label{ch17:thm:realization}
Let $L$ be an even integral lattice with nondegenerate form and a complex frame $F$.  Then $V_L$, with the
reconstructed vertex operators of Theorem~\ref{ch17:thm:reconstruct}, vacuum $1\otimes e^0$, conformal vector
$\omega_L$, central charge $\operatorname{rank}L$ and grading $(V_L)_n$, is an even-lattice vertex realisation of
$L$ (Definition~\ref{ch17:def:target}).  In particular, for $\chseventeenII$ with its Lorentzian frame, the target
\chseventeenL{LorentzianVertexTarget} is inhabited, and $\chseventeenVII$ is a near-VOA
\chseventeenL{lorentzianNearVOA} on $\chseventeenF_2\otimes\C[\chseventeenII]$ with central charge $2$, whose field of $e^\alpha$ is the
FLM-cocycle charged field and whose $L_0$ acts on $(V)_n$ by $n$.
\end{theorem}

\noindent File \leanpath{VOA/\allowbreak Moonshine/\allowbreak MoonshineCNRealization.lean}.
\begin{leancode}
/-- An even nondegenerate integral lattice with a complex bilinear frame has
an unbounded conformal vertex algebra on the genuine Fock/group carrier. -/
noncomputable def complexFrameLatticeRealization (hEven : L.Even)
    (hNondegenerate : L.form.Nondegenerate) : EvenLatticeVertexRealization L where
  grading := gradingLatGen L
  isInternal := latticeWeightGrading_isInternal L
  Y := frameStateField F hEven
  vacuum := latVacuumGen L

/-- The required Lorentzian lattice realization, with no unproved input data. -/
noncomputable def lorentzianVertexRealization : EvenLatticeVertexRealization lorentzianLattice :=
  complexFrameLatticeRealization lorentzianComplexFrame lorentzian_even lorentzian_form_nondegenerate

theorem lorentzianVertexTarget_inhabited : LorentzianVertexTarget :=
  ⟨lorentzianVertexRealization⟩

/-- The same constructed object in the repository's existing structure. -/
noncomputable def lorentzianNearVOA : NearVOA LorentzianFockSpace :=
  lorentzianVertexRealization.toUnboundedConformalVertexAlgebra.toNearVOA

theorem lorentzianNearVOA_centralCharge : lorentzianNearVOA.centralCharge = 2 := by

theorem lorentzianNearVOA_charge_field (α : lorentzianLattice.carrier) :
    lorentzianNearVOA.Y (lorentzianChargeState α) = lorentzianTargetChargedField α :=

theorem lorentzianNearVOA_energy (n : ℤ) (v : LorentzianFockSpace)
    (hv : v ∈ lorentzianWeightGrading n) :
    mode (lorentzianNearVOA.Y lorentzianNearVOA.conformalVector) 1 v = (n : ℂ) • v :=
\end{leancode}

\begin{dictionary}
\textbf{The construction is explicit.}  Every data field of \chseventeenL{complexFrameLatticeRealization} is a named
definition: \chseventeenL{grading := gradingLatGen L}, \chseventeenL{Y := frameStateField F hEven}, \chseventeenL{vacuum := latVacuumGen L},
\chseventeenL{conformalVector := latticeConformalVector L}, \chseventeenL{centralCharge := (rank L : ℂ)} (the omitted lines;
every axiom field is one of the theorems of \S\ref{ch17:sec:construction}).  So
\chseventeenL{lorentzianNearVOA} is not an existential witness but a specific structure; its fields can be computed
(\chseventeenL{lorentzianNearVOA_charge_field} is \chseventeenL{frameStateField_charge}, an instance of the defining
property of reconstruction).  \textbf{Hypotheses.} \chseventeenL{hEven} and \chseventeenL{hNondegenerate} are discharged for
$\chseventeenII$ by \chseventeenL{lorentzian_even}, \chseventeenL{lorentzian_form_nondegenerate}; the frame by
\chseventeenL{lorentzianComplexFrame}.  Nondegeneracy is a field of the target and is not used in the construction of
the fields.  \textbf{What the statement is not.} It is not a \chseventeenL{VertexOperatorAlgebra} (impossible by
Theorem~\ref{ch17:thm:grading}); simplicity of $\chseventeenVII$ is not stated; and the charged fields use the FLM
cocycle, not $\varepsilon_{\mathrm{std}}$.  \textbf{Uniqueness.} The vertex map is the unique local, creative,
$T$-covariant extension of the generating fields (Goddard), so the choice made by \chseventeenL{.choose} in
\chseventeenL{frameStateField} is immaterial.
\end{dictionary}

\begin{proof}[Proof on paper]
Assemble: internality (Theorem~\ref{ch17:thm:grading}), the field axioms (Theorem~\ref{ch17:thm:reconstruct}),
Virasoro, $L_0$, translation and grading (Theorem~\ref{ch17:thm:virasoro}); the charge fields are the
reconstructed fields of the word $[\,]$, whose leading products are \eqref{ch17:eq:chargeleading}; the
charged vacuum has weight $\langle\alpha,\alpha\rangle/2$ (Theorem~\ref{ch17:thm:grading}(ii)).  Central charge $\operatorname{rank}\chseventeenII=2$.
\end{proof}

\begin{physics}
$\chseventeenVII$ is the chiral algebra of the two light-cone bosons $X^\pm$ compactified on the even
self-dual Lorentzian torus, with the standard Klein factors.  Its central charge $2$ is the contribution of
the two longitudinal directions of the string; it is the ``$2$'' in $24+2=26$.
\end{physics}

\begin{status}
Proved with no hypotheses: \chseventeenL{lorentzianVertexRealization}, \chseventeenL{lorentzianVertexTarget_inhabited},
\chseventeenL{lorentzianNearVOA}, \chseventeenL{lorentzianNearVOA_centralCharge}, \chseventeenL{lorentzianNearVOA_charge_field},
\chseventeenL{lorentzianNearVOA_energy}; \chseventeenL{complexFrameLatticeRealization} conditional only on its displayed
hypotheses (evenness, nondegeneracy, a frame).  All printed in
\leanpath{VOA/\allowbreak Moonshine/\allowbreak MoonshineCNAudit2.lean}; that audit has 383 \texttt{\#print axioms} lines (including declarations of other files it imports); according to
the development record (\S\ref{ch17:sec:notes}) every list is a subset of
\chseventeenL{[propext, Classical.choice, Quot.sound]}, with no \chseventeenL{sorryAx}.
\end{status}

\begin{theorem}[The conformal vector in light-cone form]\label{ch17:thm:dualbasis}
Let $u=(1,0)$, $v=(0,1)$ and write $u(-1)$, $v(-1)$ for the frame oscillators acting on the first factor.
Then $\omega=\tfrac12\bigl(u(-1)v(-1)+v(-1)u(-1)\bigr)\chseventeenvac$, and $u(-1)$ is the $(-1)$-mode of the Heisenberg
current $Y(u(-1)\chseventeenvac,z)=\sum_au^\perp_ah_a(z)$ (zero mode included).  Moreover
$\langle u,u\rangle=\langle v,v\rangle=0$, $\langle u,v\rangle=1$.
\end{theorem}

\noindent File \leanpath{VOA/\allowbreak Moonshine/\allowbreak MoonshineCNConformalBasis.lean}.
\begin{leancode}
/-- Negative-one Heisenberg action, lifted to the full lattice carrier. -/
noncomputable def lorentzianCreation (α : lorentzianLattice.carrier) :
    Module.End ℂ LorentzianFockSpace :=
  LinearMap.rTensor (lorentzianLattice.carrier →₀ ℂ)
    (frameOscillator lorentzianComplexFrame α (-1))

/-- The Heisenberg current, including its momentum zero mode. -/
noncomputable def lorentzianHeisenbergField (α : lorentzianLattice.carrier) :
    FormalDist (Module.End ℂ LorentzianFockSpace) :=
  ∑ a, frameCoordinate lorentzianComplexFrame α a • frameCurrent lorentzianComplexFrame a

theorem lorentzianCreation_eq_current_mode (α : lorentzianLattice.carrier) :
    mode (lorentzianHeisenbergField α) (-1) = lorentzianCreation α := by

/-- The prescribed dual-basis formula for the actual reconstructed conformal vector. -/
theorem lorentzianNearVOA_conformal_dualBasis :
    lorentzianNearVOA.conformalVector = (1 / 2 : ℂ) •
      (lorentzianCreation ![1, 0] (lorentzianCreation ![0, 1] lorentzianVacuum) +
       lorentzianCreation ![0, 1] (lorentzianCreation ![1, 0] lorentzianVacuum)) := by
\end{leancode}

\begin{dictionary}
This theorem is the referee's check that the conformal vector is the right one in lattice terms.  The
Sugawara vector of a lattice is $\frac12\sum_i\beta_i(-1)\beta_i^\vee(-1)\chseventeenvac$ for a basis $\beta_i$ and its dual
basis $\beta_i^\vee$; for $(u,v)$ the dual basis is $(v,u)$, giving the displayed formula.  The proof is a
polynomial identity: $u^\perp=(1,i)$, $v^\perp=(\frac12,-\frac i2)$ and $u^\perp_1v^\perp_1x_1^2+\dots=\frac12(x_1^2+x_2^2)$.
The companion \chseventeenL{lorentzianIsotropicBasis_pairing} records the four pairings.
\end{dictionary}

\begin{proof}[Proof on paper]
$u(-1)v(-1)\chseventeenvac=(x_1+ix_2)(\tfrac12x_1-\tfrac i2x_2)=\tfrac12(x_1^2+x_2^2)=\omega$, and $u(-1),v(-1)$ commute.
\end{proof}

\begin{physics}
In light-cone variables the stress tensor of two bosons is $T=-{:}\partial X^+\partial X^-{:}$ (sign by
convention): it pairs the two lightlike directions.  This is the stress tensor whose modes appear in the
light-cone decomposition of \S\ref{ch17:sec:hyp}.
\end{physics}

\begin{status}
Proved unconditionally; printed in \leanpath{VOA/\allowbreak Moonshine/\allowbreak MoonshineCNAudit2.lean}; standard axioms.
\end{status}

\subsection{The momentum grading}

\begin{definition}[Momentum sectors]\label{ch17:def:momentum}
For $\alpha\in L$, the momentum-$\alpha$ sector is $\chseventeenF_r\otimes e^\alpha=\sum_w(\chseventeenF_r)_w\otimes e^\alpha$.  A
formal distribution $A$ \emph{has momentum $\alpha$} if every coefficient maps momentum $\beta$ into momentum $\alpha+\beta$.
\end{definition}

\noindent File \leanpath{VOA/\allowbreak Moonshine/\allowbreak MoonshineCNMomentum.lean}.
\begin{leancode}
noncomputable def latticeMomentumGrading (L : IntegralLattice) (α : L.carrier) :
    Submodule ℂ (LatModuleGen L) := ⨆ w : ℤ, bigradα L (w, α)

/-- A field of momentum `α` adds `α` to the momentum of its argument. -/
def HasLatticeMomentum {L : IntegralLattice}
    (A : FormalDist (End ℂ (LatModuleGen L))) (α : L.carrier) : Prop :=
  ∀ (t : ℤ) (β : L.carrier) (v : LatModuleGen L), v ∈ latticeMomentumGrading L β →
    A t v ∈ latticeMomentumGrading L (α + β)
\end{leancode}

\begin{dictionary}
\chseventeenL{bigradα L (w, α)} $=(\chseventeenF_r)_w\otimes e^\alpha$ (Chapter~\ref{ch:03}, Definition~\ref{ch03:def:gradingLatGen});
the supremum over $w$ is the whole sector, equal to the range of $p\mapsto p\otimes e^\alpha$
(\chseventeenL{latticeMomentumGrading_eq_range}).  Each sector is a full copy of the Fock space, hence
infinite-dimensional; finite-dimensional pieces are the bigraded ones (momentum \emph{and} weight).
\end{dictionary}

\begin{physics}
$U(1)^2$ charge conservation: momentum is the eigenvalue of the zero modes of the two currents.
\end{physics}

\begin{status}
Definitions.
\end{status}

\begin{theorem}[The momentum grading is internal and respected by all modes]\label{ch17:thm:momentum}
For every integral lattice $L$, $V_L=\bigoplus_\alpha(\chseventeenF_r\otimes e^\alpha)$.  For even $L$ with a complex frame,
the field of every state of momentum $\alpha$ has momentum $\alpha$; in particular in $\chseventeenVII$,
$u_{(n)}v$ has momentum $\alpha+\beta$ when $u,v$ have momenta $\alpha,\beta$.
\end{theorem}

\noindent File \leanpath{VOA/\allowbreak Moonshine/\allowbreak MoonshineCNMomentum.lean}.
\begin{leancode}
theorem latticeMomentumGrading_isInternal (L : IntegralLattice) :
    DirectSum.IsInternal (latticeMomentumGrading L) := by

theorem frameStateField_momentum (hEven : L.Even) (α : L.carrier)
    (u : LatModuleGen L) (hu : u ∈ latticeMomentumGrading L α) :
    HasLatticeMomentum (frameStateField F hEven u) α := by

theorem lorentzianNearVOA_momentum (α β : lorentzianLattice.carrier)
    (u v : LorentzianFockSpace) (hu : u ∈ latticeMomentumGrading lorentzianLattice α)
    (hv : v ∈ latticeMomentumGrading lorentzianLattice β) (n : ℤ) :
    mode (lorentzianNearVOA.Y u) n v ∈ latticeMomentumGrading lorentzianLattice (α + β) :=
\end{leancode}

\begin{dictionary}
Exactly as stated; no positivity.  The proof goes through word states, which span each sector, and the
fact that currents have momentum $0$ and charged fields momentum $\alpha$; normal products add momenta.
\end{dictionary}

\begin{proof}[Proof on paper]
$V_L=\bigoplus_{(w,\alpha)}(\chseventeenF_r)_w\otimes e^\alpha$, regroup by $\alpha$.  Each generating field has a definite momentum
(the charged field shifts $e^\beta\mapsto e^{\alpha+\beta}$, currents preserve momentum), normal-ordered products and
derivatives add momenta, and $Y$ is linear.
\end{proof}

\begin{physics}
Momentum conservation in every OPE.  The Monster Lie algebra inherits this as its $\chseventeenII$-grading
(Chapter~\ref{ch:18}).
\end{physics}

\begin{status}
Proved; printed in \leanpath{VOA/\allowbreak Moonshine/\allowbreak MoonshineCNAudit2.lean}; standard axioms.
\end{status}

%======================================================================
\section{Comparison with positive-definite lattices}\label{ch17:sec:compare}
%======================================================================

The following table compares the library's lattice vertex operator algebra \chseventeenL{latticeVOAGen}
(Chapter~\ref{ch:03}, Theorem~\ref{ch03:thm:latticeVOA}; Chapter~\ref{ch:L05},
Theorem~\ref{l05:thm:latticeVOAGen}) with \chseventeenL{complexFrameLatticeRealization}.

\begin{center}
\small
\begin{longtable}{p{0.24\textwidth}p{0.34\textwidth}p{0.34\textwidth}}
\toprule
& positive-definite $L$ (\chseventeenL{latticeVOAGen}) & indefinite $L$ (this chapter) \\ \midrule
\endhead
hypotheses & \chseventeenL{hEven}, \chseventeenL{hPos} & \chseventeenL{hEven}, \chseventeenL{hNondegenerate}, a complex frame \\
frame & real, $C=\sqrt G$ & complex, any $C$ with $C^TC=G$ over $\C$ \\
state space & $\chseventeenF_r\otimes\C[L]$ & the same type \chseventeenL{LatModuleGen L} \\
cocycle & \chseventeenL{latticeCocycle} & the same \\
weight spaces & finite-dimensional & infinite-dimensional (Thm~\ref{ch17:thm:weightzero}) \\
lower bound & $V_n=0$ for $n<0$, $V_0=\C\chseventeenvac$ & none (Thm~\ref{ch17:thm:notbounded}) \\
locality & both signs of $\langle\alpha,\beta\rangle$ & both signs, proofs re-done for complex frames \\
structure & \chseventeenL{VertexOperatorAlgebra} & \chseventeenL{NearVOA} via \chseventeenL{EvenLatticeVertexRealization} \\
central charge & $\operatorname{rank}L$ & $\operatorname{rank}L$ \\
general theory & all of Chapter~\ref{ch:L03} & re-proved in \S\ref{ch17:sec:near} \\
invariant form & FHL, via \chseventeenL{IsFullyInvariant} & FHL with charge-sensitive cutoff (\S\ref{ch17:sec:fhl}) \\
\bottomrule
\end{longtable}
\end{center}

\begin{remark}\label{ch17:rem:compare}
The positive-definite construction could not simply be instantiated: its frame lemma
\chseventeenL{orthCoord_inner} carries the hypothesis \chseventeenL{hPos}, and \chseventeenL{lorentzian_not_posDef} shows that no
\chseventeenL{PosDef} witness exists for $\chseventeenII$.  The library therefore duplicated the construction with a
frame parameter rather than introducing a false positivity witness.  Conversely, the new construction applies to positive-definite
lattices too, with the real frame viewed as complex; no comparison theorem between the two constructions is
stated.
\end{remark}

\begin{physics}
Nothing in the free-boson CFT depends on the signature except unitarity and the spectrum's lower bound:
OPEs, the stress tensor and the central charge are signature-blind.  The comparison table is the
mathematical form of that remark.
\end{physics}

%======================================================================
\section{The isotropic Fock model of a momentum sector}\label{ch17:sec:hyp}
%======================================================================

The no-ghost theorem (Chapter~\ref{ch:18}) is a statement about one momentum sector at a time.  For that
purpose the library uses a second description of a sector of $\chseventeenVII$, in which the two
light-cone oscillator families are separate polynomial variables, rather than the orthonormal complex
frame of \S\ref{ch17:sec:frame}.

\subsection{Why this is a momentum sector: the derivation by hand}\label{ch17:sec:hypwhy}

Fix a momentum $\beta=(x,y)\in\chseventeenII$ and consider the sector $\chseventeenF_2\otimes e^\beta$.  Let $u=(1,0)$, $v=(0,1)$ and
write $a(n)=v(n)$, $b(n)=u(n)$ for the corresponding frame oscillators (Definition~\ref{ch17:def:frame}), with zero
modes given by the momentum.  Then:
\begin{enumerate}[label=(\arabic*)]
\item \emph{Heisenberg relations in null coordinates.}  Since $\langle u,u\rangle=\langle v,v\rangle=0$ and $\langle u,v\rangle=1$,
  Theorem~\ref{ch17:thm:frame} gives $[a(m),a(n)]=[b(m),b(n)]=0$ and $[a(m),b(n)]=m\delta_{m+n,0}$.
\item \emph{Zero modes.}  On $e^\beta$, $a(0)=\langle v,\beta\rangle=x$ and $b(0)=\langle u,\beta\rangle=y$.  We write
  $\kappa=x$, $\lambda=y$.
\item \emph{The stress tensor.}  By Theorem~\ref{ch17:thm:dualbasis}, $\omega=\frac12(u(-1)v(-1)+v(-1)u(-1))\chseventeenvac$, so
  $T(z)=\frac12\bigl({:}b(z)a(z){:}+{:}a(z)b(z){:}\bigr)={:}a(z)b(z){:}$ (the two currents commute up to a central term
  that normal ordering removes), i.e.
  \begin{equation*}
    L_n=\sum_{k\in\Z}{:}a(n-k)\,b(k){:} .
  \end{equation*}
\item \emph{Reading off the terms.}  For $n\ge1$ the summands of $L_{-n}$ are: $k=-n$, giving $a(0)b(-n)=\kappa\,b(-n)$;
  $k=0$, giving $a(-n)b(0)=\lambda\,a(-n)$; $-n<k<0$, giving $a(-n-k)b(k)$, a product of two creation operators;
  and $k>0$ or $k<-n$, an annihilator to the right of a creator.  For $n=0$ the term $a(0)b(0)=\kappa\lambda$ appears
  once and every other term counts level, so $L_0=xy+\text{level}$, matching the weight
  $\Delta(\beta)+\text{level}$.
\end{enumerate}
Now rename the creation operators as variables: $a(-j)\chseventeenvac\mapsto a_j$, $b(-j)\chseventeenvac\mapsto b_j$.  The Fock space of
the two commuting families is the polynomial ring $\C[a_j,b_j:j\ge1]$, on which $a(-j)$, $b(-j)$ multiply and, by (1),
$a(j)=j\,\partial/\partial b_j$, $b(j)=j\,\partial/\partial a_j$.  This is $\chseventeenF^{\mathrm{iso}}$ below; as a Lean type it is literally
\chseventeenL{FockN 2}, i.e.\ $\chseventeenF_2$ \emph{in null coordinates}, only the meaning of the variables changes.  The
explicit change of variables from the orthonormal frame is the algebra isomorphism of Theorem~\ref{ch17:thm:sectorequiv};
that it intertwines the Virasoro actions is proved in Chapter~\ref{ch:18}.  The terms of Definition~\ref{ch17:def:hypsugawara}
below are exactly the four kinds listed in (4).

\paragraph{Where $c=26$ enters.}  Nothing in this section or in \S\ref{ch17:sec:tensor} uses $c=26$: the light-cone
splitting (Theorems~\ref{ch17:thm:hyppbw} and~\ref{ch17:thm:tensorlightcone}) holds for every central charge and for an
arbitrary matter operator family.  The value $26$ is used in Chapter~\ref{ch:18}, in the abstract ``$c=26$ lemma'': on the relevant
weight spaces $[L_2+\tfrac32L_1^2,\,L_{-2}+\tfrac32L_{-1}^2]$ acts as the scalar $\tfrac c2-13$, so for $c=26$ (and only then)
$L_1$ and $L_2$ map the spurious states into spurious states.  This makes the spurious part of a weight-one physical
state again physical, hence null, so that physical states modulo null states are the transverse ones.  This section presents that model (the 20 files
\texttt{MoonshineCNHyperbolic*}); its identification with the actual sectors of $\chseventeenVII$ is
Theorem~\ref{ch17:thm:sectorequiv} below and the comparison of Virasoro actions is in Chapter~\ref{ch:18}.

\subsection{Oscillators and level}

\begin{definition}[The isotropic Fock space]\label{ch17:def:hypfock}
$\chseventeenF^{\mathrm{iso}}=\C[a_j,b_j:j\ge1]$, graded by level ($a_j,b_j$ have level $j$).  For $\kappa,\lambda\in\C$ the
oscillators are, for $j\ge1$,
\begin{equation*}
  a(-j)=a_j\cdot,\quad a(j)=j\,\frac{\partial}{\partial b_j},\quad a(0)=\kappa;\qquad
  b(-j)=b_j\cdot,\quad b(j)=j\,\frac{\partial}{\partial a_j},\quad b(0)=\lambda .
\end{equation*}
\end{definition}

\noindent File \leanpath{VOA/\allowbreak Moonshine/\allowbreak MoonshineCNHyperbolicFock.lean}.
\begin{leancode}
abbrev hypFock := FockN 2

/-- One isotropic oscillator family with arbitrary scalar zero mode. -/
def hypOscillator (a : Fin 2) (q : ℂ) (n : ℤ) : Module.End ℂ hypFock :=
  if n < 0 then oscN 2 1 a n
  else if 0 < n then oscN 2 1 a.rev n
  else q • LinearMap.id

/-- Integer level grading, independent of momentum. -/
def hypLevel (N : ℤ) : Submodule ℂ hypFock := gradingZN 2 N
\end{leancode}

\begin{dictionary}
\chseventeenL{FockN 2} $=\C[x_{j,i}:j\in\N,i\in\{0,1\}]$; the variable \chseventeenL{X (j,0)} is $a_{j+1}$ and \chseventeenL{X (j,1)} is
$b_{j+1}$.  \chseventeenL{oscN 2 1 c n} for $n<0$ multiplies by \chseventeenL{X (-n-1, c)}, for $n>0$ is $n\,\partial/\partial$\chseventeenL{X (n-1, c)}.
The twist is \chseventeenL{a.rev}: \chseventeenL{Fin.rev} exchanges $0\leftrightarrow1$, so the \emph{positive} modes of one family
differentiate the \emph{other} family.  The family index \chseventeenL{a : Fin 2} names $a$ (\chseventeenL{0}) or $b$ (\chseventeenL{1});
\chseventeenL{hypOscA κ} $=$ \chseventeenL{hypOscillator 0 κ}, \chseventeenL{hypOscB ell} $=$ \chseventeenL{hypOscillator 1 ell}.  The zero mode
\chseventeenL{q} is a free complex parameter (the momentum), so the model describes \emph{all} sectors at once.
\chseventeenL{gradingZN 2 N} is the level-$N$ subspace (zero for $N<0$).  \textbf{Convention.} Under the
identification with the actual sectors used in Chapter~\ref{ch:18} (\chseventeenL{hypLorentzianSector_zeroCharges}),
the $a$-family is the current of $(0,1)$ and the $b$-family that of $(1,0)$, so at momentum $(x,y)$ one has
$\kappa=x$, $\lambda=y$ and $\kappa\lambda=xy$, as derived in \S\ref{ch17:sec:hypwhy}.  Some of the authors' working notes use the opposite
labelling ($\kappa=y$); since everything below is symmetric under the swap (Theorem~\ref{ch17:thm:hyplightcone}(iii)), the
difference is harmless, but see \S\ref{ch17:sec:traps}.
\end{dictionary}

\begin{physics}
$a_j,b_j$ are the light-cone creation operators $\alpha^+_{-j},\alpha^-_{-j}$; the two families commute among
themselves and are conjugate to each other, as lightlike directions are.  $\kappa,\lambda$ are the light-cone
momenta $p^+,p^-$.
\end{physics}

\begin{status}
Definitions; printed in \leanpath{VOA/\allowbreak Moonshine/\allowbreak MoonshineCNAudit5.lean}.
\end{status}

\begin{theorem}[Hyperbolic Heisenberg relations; level]\label{ch17:thm:hypheis}
$[\varphi(m),\varphi'(n)]=m\,\delta_{\varphi'\ne\varphi}\,\delta_{m+n,0}$ for families $\varphi,\varphi'\in\{a,b\}$ and any zero modes; the level
grading is internal with finite-dimensional pieces, $\chseventeenF^{\mathrm{iso}}_N=0$ for $N<0$,
$\chseventeenF^{\mathrm{iso}}_0=\C$; and $\varphi(n)$ lowers level by $n$.
\end{theorem}

\noindent File \leanpath{VOA/\allowbreak Moonshine/\allowbreak MoonshineCNHyperbolicFock.lean}.
\begin{leancode}
/-- The full hyperbolic Heisenberg relation, including both scalar zero modes. -/
theorem hypOscillator_comm (a b : Fin 2) (q r : ℂ) (m n : ℤ) :
    hypOscillator a q m * hypOscillator b r n -
      hypOscillator b r n * hypOscillator a q m =
    (if a = b.rev ∧ m + n = 0 then (m : ℂ) else 0) • LinearMap.id := by

theorem hypLevel_isInternal : DirectSum.IsInternal hypLevel := gradingZN_isInternal 2

theorem hypLevel_finiteDimensional (N : ℤ) : FiniteDimensional ℂ (hypLevel N) :=

theorem hypOscillator_mem_level (a : Fin 2) (q : ℂ) (n N : ℤ) (p : hypFock)
    (hp : p ∈ hypLevel N) : hypOscillator a q n p ∈ hypLevel (N - n) := by
\end{leancode}

\begin{dictionary}
\chseventeenL{a = b.rev} says the two families differ.  So $[a(m),a(n)]=[b(m),b(n)]=0$ and $[a(m),b(n)]=m\delta_{m+n,0}$:
the Heisenberg algebra of $\chseventeenII\otimes\C$ in the isotropic basis.
\end{dictionary}

\begin{proof}[Proof on paper]
$[j\partial_{b_j},b_j\cdot]=j$; all other pairs commute; zero modes are scalars.
\end{proof}

\begin{physics}
$[\alpha^+_m,\alpha^-_n]=m\delta_{m+n,0}$ and $[\alpha^\pm_m,\alpha^\pm_n]=0$: the light-cone oscillator algebra.
\end{physics}

\begin{status}
Proved; printed in \leanpath{VOA/\allowbreak Moonshine/\allowbreak MoonshineCNAudit5.lean}; standard axioms.
\end{status}

\subsection{The Sugawara operators}

\begin{definition}[Explicit Sugawara operators]\label{ch17:def:hypsugawara}
For $\kappa,\lambda\in\C$ and $n\ge1$,
\begin{align*}
L_{-n}&=\kappa b_n+\lambda a_n+\sum_{i=1}^{n-1}a_{n-i}b_i+\sum_{k\ge1}k\Bigl(a_{n+k}\frac{\partial}{\partial a_k}+b_{n+k}\frac{\partial}{\partial b_k}\Bigr),\\
L_{n}&=n\kappa\frac{\partial}{\partial a_n}+n\lambda\frac{\partial}{\partial b_n}+\sum_{k>n}k\Bigl(a_{k-n}\frac{\partial}{\partial a_k}+b_{k-n}\frac{\partial}{\partial b_k}\Bigr)
  +\sum_{k=1}^{n-1}k(n-k)\frac{\partial^2}{\partial a_k\partial b_{n-k}},\\
L_0&=\kappa\lambda+\sum_{k\ge1}k\Bigl(a_k\frac{\partial}{\partial a_k}+b_k\frac{\partial}{\partial b_k}\Bigr).
\end{align*}
These are the normal-ordered sums $L_n=\sum_k{:}a(n-k)b(k){:}$.
\end{definition}

\noindent File \leanpath{VOA/\allowbreak Moonshine/\allowbreak MoonshineCNHyperbolicSugawara.lean}.
\begin{leancode}
/-- The first-order part of `L_(-n)`: each variable of level `j+1` shifts up by `n`. -/
def hypRaisingDerivation (n : ℕ) : Derivation ℂ hypFock hypFock :=
  mkDerivation ℂ fun v => (v.1 + 1 : ℂ) • X (v.1 + n, v.2)

/-- The first-order part of `L_n`, including the two momentum derivatives. -/
def hypLoweringDerivation (κ ell : ℂ) (n : ℕ) : Derivation ℂ hypFock hypFock :=
  mkDerivation ℂ fun v =>
    if v.1 + 1 = n then (v.1 + 1 : ℂ) • MvPolynomial.C (if v.2 = 0 then κ else ell)
    else if n < v.1 + 1 then (v.1 + 1 : ℂ) • X (v.1 - n, v.2)
    else 0

/-- The multiplication part of `L_(-(j+1))`, with exactly `j` quadratic terms. -/
def hypNegativePotential (κ ell : ℂ) (j : ℕ) : hypFock :=
  κ • X (j, 1) + ell • X (j, 0) +
    ∑ k ∈ Finset.range j, X (j - 1 - k, 0) * X (k, 1)

/-- The finite second-order part of `L_(j+1)`. -/
def hypPositiveSecondOrder (j : ℕ) : Module.End ℂ hypFock :=
  ∑ k ∈ Finset.range j, ((k + 1 : ℂ) * (j - k : ℂ)) •
    ((pderiv (k, 0)).toLinearMap * (pderiv (j - 1 - k, 1)).toLinearMap)

/-- The explicit Sugawara operator at arbitrary momentum. -/
def hypSugawara (κ ell : ℂ) (n : ℤ) : Module.End ℂ hypFock :=
  if n < 0 then
    LinearMap.mul ℂ hypFock (hypNegativePotential κ ell (-n - 1).toNat) +
      (hypRaisingDerivation (-n).toNat).toLinearMap
  else if 0 < n then
    (hypLoweringDerivation κ ell n.toNat).toLinearMap +
      hypPositiveSecondOrder (n - 1).toNat
  else (κ * ell) • LinearMap.id + hypEnergyDerivation.toLinearMap
\end{leancode}

\begin{dictionary}
\chseventeenL{mkDerivation ℂ f} is the unique derivation of the polynomial ring with prescribed values $f(v)$ on
the variables; this encodes the infinite first-order sums without any infinite sum of operators.  With
$v=(k-1,i)$ (the variable of level $k$ in family $i\in\{0,1\}$): \chseventeenL{hypRaisingDerivation n} sends it to $k$ times the
variable of level $k+n$; \chseventeenL{hypLoweringDerivation κ ell n} sends it to $n\kappa$ or $n\lambda$ if $k=n$, to $k$ times
the variable of level $k-n$ if $k>n$, and to $0$ otherwise.  So the term $n\kappa\,\partial/\partial a_n$ is the
$k=n$ case for family $a$.  \chseventeenL{hypNegativePotential κ ell j} is $\kappa b_{j+1}+\lambda a_{j+1}+\sum_{k<j}a_{j-k}b_{k+1}$;
\chseventeenL{pderiv (k, 0)} is $\partial/\partial a_{k+1}$; \chseventeenL{hypEnergyDerivation} $=$ \chseventeenL{hypRaisingDerivation 0} is the
level operator.  \chseventeenL{(-n - 1).toNat} etc.\ convert $\Z$ to $\N$ on the relevant branch.  \textbf{Fidelity.} The
definition is the explicit formula of Definition~\ref{ch17:def:hypsugawara} term by term; that it \emph{is} the
Sugawara construction is then proved (Theorem~\ref{ch17:thm:hypvirasoro}), not assumed.
\end{dictionary}

\begin{physics}
$L_n=\sum_k{:}\alpha^+_{n-k}\alpha^-_k{:}$, the light-cone stress tensor.  The terms $\kappa b_n$ and $\lambda a_n$ are
$p^+\alpha^-_{-n}$ and $p^-\alpha^+_{-n}$: in light-cone gauge the constraint $L_{-n}$ is dominated by
$p^+\alpha^-_{-n}$ when $p^+\neq0$, which is why one can ``solve for'' $\alpha^-$.  This is exactly the leading term
exploited in Theorem~\ref{ch17:thm:hyppbw}.
\end{physics}

\begin{status}
Definitions; printed in \leanpath{VOA/\allowbreak Moonshine/\allowbreak MoonshineCNAudit5.lean}.
\end{status}

\begin{theorem}[Commutant]\label{ch17:thm:hypcommutant}
An endomorphism of $\chseventeenF^{\mathrm{iso}}$ that commutes with all nonzero oscillator modes of both families is a
scalar, namely the constant term of its value on $1$.
\end{theorem}

\noindent File \leanpath{VOA/\allowbreak Moonshine/\allowbreak MoonshineCNHyperbolicCommutant.lean}.
\begin{leancode}
/-- Every endomorphism commuting with the nonzero oscillators is scalar. -/
theorem hypCommutant_scalar (D : Module.End ℂ hypFock) (charges : Fin 2 → ℂ)
    (hD : ∀ a n, n ≠ 0 → D * hypOscillator a (charges a) n =
      hypOscillator a (charges a) n * D) :
    D = coeff 0 (D 1) • LinearMap.id := by
\end{leancode}

\begin{dictionary}
\chseventeenL{charges} are the zero modes (irrelevant, as only $n\ne0$ is used).  \chseventeenL{coeff 0 (D 1)} is the constant term of
$D(1)$.  No finite-dimensionality or continuity is assumed: the operator is arbitrary.
\end{dictionary}

\begin{proof}[Proof on paper]
Commuting with all multiplications by variables gives $D(p)=p\,D(1)$; commuting with all $\partial/\partial a_j$,
$\partial/\partial b_j$ gives that every partial derivative of $D(1)$ vanishes, so $D(1)$ is constant (characteristic $0$).
\end{proof}

\begin{physics}
The Fock module of the Heisenberg algebra is irreducible (Schur's lemma): any operator built from oscillators
whose commutators with all oscillators vanish is a c-number.  This is how central charges are computed: the
anomaly in $[L_m,L_n]$ commutes with everything, so it is a number, read off on the vacuum.
\end{physics}

\begin{status}
Proved; printed in \leanpath{VOA/\allowbreak Moonshine/\allowbreak MoonshineCNAudit5.lean}; standard axioms.
\end{status}

\begin{theorem}[Virasoro with $c=2$; oscillator commutators; degree]\label{ch17:thm:hypvirasoro}
For all $\kappa,\lambda$ and $m,n\in\Z$:
$[L_m,L_n]=(m-n)L_{m+n}+\tfrac{2}{12}(m^3-m)\delta_{m+n,0}$; $[L_m,\varphi(n)]=-n\,\varphi(m+n)$ for both families (with
$\varphi(0)$ the scalar zero mode); $L_0=\kappa\lambda+N$ on level $N$, and the level-$N$ space is exactly the
$(\kappa\lambda+N)$-eigenspace of $L_0$; and $L_n$ maps level $N$ to level $N-n$.
\end{theorem}

\noindent{\raggedright Files \leanpath{VOA/\allowbreak Moonshine/\allowbreak MoonshineCNHyperbolicVirasoro.lean},
\leanpath{VOA/\allowbreak Moonshine/\allowbreak MoonshineCNHyperbolicMixedModes.lean},
\leanpath{VOA/\allowbreak Moonshine/\allowbreak MoonshineCNHyperbolicModeGrading.lean}.\par}
\begin{leancode}
/-- The concrete operators satisfy Virasoro at central charge exactly two. -/
theorem hypSugawara_virasoro (κ ell : ℂ) (m n : ℤ) :
    hypSugawara κ ell m * hypSugawara κ ell n -
      hypSugawara κ ell n * hypSugawara κ ell m =
      ((m : ℂ) - n) • hypSugawara κ ell (m + n) +
        if m + n = 0 then (((m : ℂ) ^ 3 - m) / 12 * 2) • 1 else 0 := by

/-- The complete mixed relation, including zero oscillator index and all momenta. -/
theorem hypSugawara_oscillator_comm (κ ell : ℂ) (m : ℤ) (a : Fin 2) (n : ℤ) :
    ⁅hypSugawara κ ell m, hypOscillator a (if a = 0 then κ else ell) n⁆ =
      -(n : ℂ) • hypOscillator a (if a = 0 then κ else ell) (m + n) := by

theorem hypLevel_eq_sugawara_eigenspace (κ ell : ℂ) (N : ℤ) :
    hypLevel N = Module.End.eigenspace (hypSugawara κ ell 0) (κ * ell + (N : ℂ)) := by

/-- Every signed mode has level degree exactly minus its index. -/
theorem hypSugawara_mem_level (κ ell : ℂ) (n N : ℤ) (p : hypFock)
    (hp : p ∈ hypLevel N) : hypSugawara κ ell n p ∈ hypLevel (N - n) := by
\end{leancode}

\begin{dictionary}
The zero mode in the oscillator commutator is tied to the momentum by \chseventeenL{(if a = 0 then κ else ell)}: the
family-$a$ oscillator has zero mode $\kappa$, family $b$ has $\lambda$ (this is necessary: $[L_m,a(-m)]=m\,a(0)$ must
produce the momentum).  \chseventeenL{Module.End.eigenspace f μ} is $\ker(f-\mu)$.  The central term is exactly
$\frac{m^3-m}{12}\cdot2$.  No hypothesis on $\kappa,\lambda$.
\end{dictionary}

\begin{proof}[Proof on paper]
The mixed relation is checked on the generators of the polynomial ring and
on $1$ (both sides are derivations up to the multiplication terms, which are handled by the explicit
creation commutators \chseventeenL{hypSugawara_negative_mul_variable}, \chseventeenL{hypSugawara_positive_creation}).  Then
$D=[L_m,L_n]-(m-n)L_{m+n}$ commutes with every oscillator by the Jacobi identity
($[[L_m,L_n],\varphi(k)]=(m-n)[L_{m+n},\varphi(k)]$), so it is a scalar by Theorem~\ref{ch17:thm:hypcommutant}; it has level
degree $-(m+n)$, so it vanishes unless $m+n=0$; for $n=-m$, $m\ge1$, its value is the constant term of
$L_mL_{-m}1-2mL_01=(2m\kappa\lambda+\sum_{k=1}^{m-1}k(m-k))-2m\kappa\lambda=\frac{m^3-m}6$
(\chseventeenL{hypSugawara_positive_negative_constant}, \chseventeenL{hypSugawara_central_sum}).  The level statements follow from
the explicit $L_0$.
\end{proof}

\begin{physics}
Two bosons give $c=2$, at every momentum, as in Theorem~\ref{ch17:thm:virasoro}.  The mixed relation says
that the light-cone oscillators are primary fields of weight one.
\end{physics}

\begin{status}
Proved with no hypotheses; printed in \leanpath{VOA/\allowbreak Moonshine/\allowbreak MoonshineCNAudit7.lean}; standard axioms.
\end{status}

\subsection{The contravariant form}

\begin{definition}[Isotropic monomial form]\label{ch17:def:hypform}
For an exponent vector $d$ (a monomial $\prod a_j^{\alpha_j}b_j^{\beta_j}$) let $z(d)=\prod_{\text{variables}}(\text{level})^{\text{exp}}\cdot\text{exp}!$
and let $\bar d$ be $d$ with the two families exchanged.  The form $B^{\mathrm{iso}}$ pairs monomial $d$ with
monomial $\bar d$ with value $z(d)$ and is zero otherwise: $B^{\mathrm{iso}}(a^\alpha b^\beta,a^{\alpha'}b^{\beta'})=\delta_{\alpha\beta'}\delta_{\beta\alpha'}z(\alpha)z(\beta)$.
\end{definition}

\noindent File \leanpath{VOA/\allowbreak Moonshine/\allowbreak MoonshineCNHyperbolicForm.lean}.
\begin{leancode}
def hypMonomialNorm (d : (ℕ × Fin 2) →₀ ℕ) : ℂ :=
  d.prod fun v k => (v.1 + 1 : ℂ) ^ k * (k.factorial : ℂ)

/-- The explicit isotropic monomial pairing. -/
def hypForm : LinearMap.BilinForm ℂ hypFock :=
  Finsupp.lsum ℂ fun d =>
    LinearMap.smulRight LinearMap.id (hypMonomialNorm d • lcoeff ℂ (hypSwapExponent d))
\end{leancode}

\begin{dictionary}
\chseventeenL{(ℕ × Fin 2) →₀ ℕ} is an exponent vector (finitely supported).  \chseventeenL{d.prod f} is $\prod_vf(v,d_v)$.
\chseventeenL{hypSwapExponent} applies \chseventeenL{Fin.rev} to the family index.  \chseventeenL{hypForm p q} $=\sum_dp_d\,z(d)\,q_{\bar d}$,
where $p_d$ is the coefficient of $p$; this is the formula of the definition.
\end{dictionary}

\begin{physics}
A contravariant pairing of the light-cone Fock space.  Each oscillator is adjoint to its own opposite mode,
$a(-j)^\dagger=a(j)=j\,\partial/\partial b_j$ (Theorem~\ref{ch17:thm:hypform}); because $a(j)$ differentiates the \emph{other}
family, the form pairs an $a$-monomial with the corresponding $b$-monomial, reflecting $\langle u,v\rangle=1$,
$\langle u,u\rangle=\langle v,v\rangle=0$.  The form takes real values on polynomials with real coefficients, and there it is
not positive: $a_1+b_1$ has norm $2$ and $a_1-b_1$ has norm $-2$.  This real form is the one that carries the
``ghost'' story: the light-cone sector contains states of negative norm.
\end{physics}

\begin{status}
Definition; printed in \leanpath{VOA/\allowbreak Moonshine/\allowbreak MoonshineCNAudit5.lean}.
\end{status}

\begin{theorem}[Properties of the isotropic form]\label{ch17:thm:hypform}
$B^{\mathrm{iso}}$ is symmetric, $B^{\mathrm{iso}}(1,q)$ is the constant term of $q$, it is nondegenerate, different
levels are orthogonal and it is nondegenerate on each level; every oscillator satisfies
$B^{\mathrm{iso}}(\varphi(n)p,q)=B^{\mathrm{iso}}(p,\varphi(-n)q)$; and every Sugawara mode satisfies
$B^{\mathrm{iso}}(L_np,q)=B^{\mathrm{iso}}(p,L_{-n}q)$.
\end{theorem}

\noindent{\raggedright Files \leanpath{VOA/\allowbreak Moonshine/\allowbreak MoonshineCNHyperbolicForm.lean},
\leanpath{VOA/\allowbreak Moonshine/\allowbreak MoonshineCNHyperbolicSugawaraAdjoint.lean}.\par}
\begin{leancode}
theorem hypForm_symm (p q : hypFock) : hypForm p q = hypForm q p := by

theorem hypForm_one (q : hypFock) : hypForm 1 q = coeff 0 q := by

theorem hypForm_nondegenerate (p : hypFock) (hp : p ≠ 0) : ∃ q, hypForm p q ≠ 0 := by

theorem hypForm_weight_nondegenerate (N : ℤ) (p : hypFock) (hp : p ∈ hypLevel N)
    (hB : ∀ q ∈ hypLevel N, hypForm p q = 0) : p = 0 := by

/-- Both isotropic oscillator families are adjoint under reversal of the mode. -/
theorem hypForm_oscillator_adjoint (a : Fin 2) (charge : ℂ) (n : ℤ) (p q : hypFock) :
    hypForm (hypOscillator a charge n p) q =
      hypForm p (hypOscillator a charge (-n) q) := by

/-- All signed modes are contravariant for the concrete hyperbolic form. -/
theorem hypForm_sugawara_adjoint (κ ell : ℂ) (n : ℤ) (p q : hypFock) :
    hypForm (hypSugawara κ ell n p) q = hypForm p (hypSugawara κ ell (-n) q) := by
\end{leancode}

\begin{dictionary}
\textbf{A sign to notice.} Here oscillators are \emph{self}-adjoint, $\varphi(n)^\dagger=\varphi(-n)$, whereas for the lattice form
of Chapter~\ref{ch:04} (Definition~\ref{ch04:def:uform}) one has $h(n)^\dagger=-h(-n)$.  The two forms are therefore
different and are not identified by the identity map; this is an
explicit convention difference (and the module docstring says ``This is the sector contravariant form, not the
opposite-momentum FHL form of the lattice vertex algebra'').  Both make $L_n$ adjoint to $L_{-n}$, which is all
the no-ghost theorem needs.  \chseventeenL{hypForm_oscillator_adjoint} holds for any zero mode \chseventeenL{charge}.
\end{dictionary}

\begin{proof}[Proof on paper]
Symmetry and level orthogonality are visible from the definition.  On level $N$ the Gram matrix in the
monomial basis is a permutation matrix ($d\leftrightarrow\bar d$) times nonzero diagonal entries $z(d)$, hence
invertible.  Adjointness of $a(-j)=a_j\cdot$ and $a(j)=j\partial/\partial b_j$ amounts to $z(d+e_v)=j(d_v+1)z(d)$ for a
variable $v$ of level $j$ (\chseventeenL{hypMonomialNorm_add_single}).  For $L_n$: the negative modes are adjoint to the
positive ones by induction on polynomials using the commutators with creation operators; the remaining cases
follow from symmetry.
\end{proof}

\begin{physics}
$L_n^\dagger=L_{-n}$ is the reality of the stress tensor, the property that makes ``null'' (zero-norm) spurious
states decouple.  Level-wise nondegeneracy is what guarantees that the radical of the form on physical states
is exactly the null states.
\end{physics}

\begin{status}
Proved; the form statements printed in \leanpath{VOA/\allowbreak Moonshine/\allowbreak MoonshineCNAudit5.lean},
\chseventeenL{hypForm_sugawara_adjoint} in \leanpath{VOA/\allowbreak Moonshine/\allowbreak MoonshineCNAudit7.lean}; standard axioms.
\end{status}

\subsection{One-family vacua}

\begin{definition}[One-family subspace]\label{ch17:def:onefamily}
For a family $\varphi\in\{a,b\}$, $\chseventeenF^{\mathrm{iso}}(\varphi)$ is the subalgebra of polynomials in the variables of family
$c$ only; its augmentation subspace is the set of $p$ with zero constant term.
\end{definition}

\noindent File \leanpath{VOA/\allowbreak Moonshine/\allowbreak MoonshineCNHyperbolicOneFamily.lean}.
\begin{leancode}
/-- Polynomials using just the indicated isotropic family. -/
def hypOneFamily (a : Fin 2) : Submodule ℂ hypFock :=
  (supported ℂ {v : ℕ × Fin 2 | v.2 = a}).toSubmodule
\end{leancode}

\begin{dictionary}
\chseventeenL{supported ℂ s} is the subalgebra of polynomials whose variables lie in $s$.
\end{dictionary}

\begin{physics}
Polynomials in the $\alpha^+$ oscillators only: the states with no $\alpha^-$ excitation.  In the Goddard--Thorn
argument the role of the transverse/longitudinal space $K$ is played by the common kernel of the positive modes
$K_n=k\cdot\alpha_n$ of a lightlike vector $k$; here $K_n=a(n)=n\partial/\partial b_n$, whose common kernel is exactly the $b$-free
polynomials.
\end{physics}

\begin{status}
Definition.
\end{status}

\begin{theorem}[One-family subspace: kernel, form, radical]\label{ch17:thm:onefamily}
(i) $p\in\chseventeenF^{\mathrm{iso}}(\varphi)$ iff $\varphi(j)p=0$ for all $j\ge1$.  (ii) For $p,q\in\chseventeenF^{\mathrm{iso}}(\varphi)$,
$B^{\mathrm{iso}}(p,q)=p(0)q(0)$.  (iii) Hence the radical of $B^{\mathrm{iso}}$ restricted to $\chseventeenF^{\mathrm{iso}}(\varphi)$ is the augmentation
subspace, and the quotient is one-dimensional, identified with $\C$ by evaluation at $0$.
\end{theorem}

\noindent File \leanpath{VOA/\allowbreak Moonshine/\allowbreak MoonshineCNHyperbolicOneFamily.lean}.
\begin{leancode}
/-- Exact common-kernel description for either isotropic family. -/
theorem hypOneFamily_iff_oscillator (a : Fin 2) (charge : ℂ) (p : hypFock) :
    p ∈ hypOneFamily a ↔ ∀ j : ℕ, hypOscillator a charge (j + 1 : ℤ) p = 0 := by

/-- The form restricted to one isotropic family is the product of evaluations. -/
theorem hypOneFamily_form (a : Fin 2) (p q : hypFock)
    (hp : p ∈ hypOneFamily a) (hq : q ∈ hypOneFamily a) :
    hypForm p q = coeff 0 p * coeff 0 q := by

theorem hypOneFamily_radical_iff (a : Fin 2) (p : hypOneFamily a) :
    (∀ q : hypOneFamily a, hypForm p q = 0) ↔ p ∈ hypOneFamilyAugmentation a := by

/-- The concrete one-family radical quotient is evaluation at zero. -/
def hypOneFamilyQuotientEquiv (a : Fin 2) :
    (hypOneFamily a ⧸ hypOneFamilyAugmentation a) ≃ₗ[ℂ] ℂ :=
\end{leancode}

\begin{dictionary}
(i): the positive modes of family $\varphi$ differentiate the \emph{other} family, so they kill exactly the polynomials
not involving it.  \chseventeenL{⧸} is the quotient module and \chseventeenL{≃ₗ[ℂ]} a linear isomorphism;
\chseventeenL{hypOneFamilyQuotientEquiv_mk} (same file) shows it is evaluation at $0$.
\end{dictionary}

\begin{proof}[Proof on paper]
(ii) Exchanging families sends a nonconstant $c$-only monomial to a monomial of the other family, which pairs to
zero with every $c$-only monomial; only constants survive.  (iii) follows.
\end{proof}

\begin{physics}
Within the ``longitudinal'' states, everything except the ground state is null: the $\alpha^+$ excitations have zero
norm against each other.  In the sector string module of Chapter~\ref{ch:18} (matter $\otimes$ this space) this is what
makes the physical quotient equal to the matter transverse space.
\end{physics}

\begin{status}
Proved; printed in \leanpath{VOA/\allowbreak Moonshine/\allowbreak MoonshineCNAudit6.lean}; standard axioms.
\end{status}

\subsection{The light-cone splitting}

Order monomials by their \emph{$b$-count} (total degree in the $b$'s) and then by their \emph{$b$-level} (total
level of the $b$'s).

\begin{theorem}[Leading term, PBW map, light-cone splitting]\label{ch17:thm:hyppbw}\label{ch17:thm:hyplightcone}
Fix $\kappa,\lambda$.  (i) For $p$ homogeneous of $b$-count $d$ and $b$-level $w$,
$L_{-(j+1)}p-\kappa\,b_{j+1}p$ is strictly lower than $(d+1,w+j+1)$ in the lexicographic order.
(ii) Let $\Phi$ send a monomial $a^\alpha b_{\lambda_1}\cdots b_{\lambda_k}$ ($\lambda_1\ge\dots\ge\lambda_k$) to
$L_{-\lambda_1}\cdots L_{-\lambda_k}a^\alpha$.  Then $\Phi(m)-\kappa^km$ is strictly lower than $m$; hence if $\kappa\ne0$,
$\Phi$ is a linear automorphism of $\chseventeenF^{\mathrm{iso}}$.
(iii) If $\kappa\ne0$, then $\chseventeenF^{\mathrm{iso}}=\Bigl(\sum_{n\ge1}\operatorname{im}L_{-n}\Bigr)\oplus\chseventeenF^{\mathrm{iso}}(a)$; if $\lambda\ne0$,
the same with $\chseventeenF^{\mathrm{iso}}(b)$; so for every nonzero momentum one of the two common kernels
$\bigcap_j\ker \varphi(j)$ is a complement of the spurious space.
\end{theorem}

\noindent{\raggedright Files \leanpath{VOA/\allowbreak Moonshine/\allowbreak MoonshineCNHyperbolicFiltration.lean},
\leanpath{VOA/\allowbreak Moonshine/\allowbreak MoonshineCNHyperbolicPBWMap.lean},
\leanpath{VOA/\allowbreak Moonshine/\allowbreak MoonshineCNHyperbolicPBWEquiv.lean},
\leanpath{VOA/\allowbreak Moonshine/\allowbreak MoonshineCNHyperbolicLightCone.lean},
\leanpath{VOA/\allowbreak Moonshine/\allowbreak MoonshineCNHyperbolicSwap.lean}.\par}
\begin{leancode}
theorem hypSugawara_negative_leading_filtered (κ ell : ℂ) (j d w : ℕ)
    (p : hypFock) (hp : p ∈ hypBHomogeneous d w) :
    hypSugawara κ ell (-(j + 1 : ℤ)) p - κ • (X (j, 1) * p) ∈
      hypBFiltered (d + 1) (w + j + 1) := by

def hypPBWMap (κ ell : ℂ) : Module.End ℂ hypFock :=
  Finsupp.lsum ℂ fun e => LinearMap.smulRight LinearMap.id
    (stringNegativeWord (fun j => hypSugawara κ ell (-(j + 1 : ℤ)))
      (monomial (hypAExponent e) 1) (hypCanonicalBWord e))

/-- The sorted-word map is triangular with diagonal `κ` to the b-count. -/
theorem hypPBWMap_leading (κ ell : ℂ) (e : (ℕ × Fin 2) →₀ ℕ) :
    hypPBWMap κ ell (monomial e 1) -
      κ ^ (Finsupp.weight hypBCountWeight e) • monomial e 1 ∈
      hypBFiltered (Finsupp.weight hypBCountWeight e) (Finsupp.weight hypBLevelWeight e) := by

/-- A linear automorphism of the polynomial space supplied by negative Sugawara words. -/
def hypPBWEquiv (κ ell : ℂ) (hκ : κ ≠ 0) : hypFock ≃ₗ[ℂ] hypFock :=

/-- The actual sum of negative Sugawara ranges. -/
def hypNegativeSpurious (κ ell : ℂ) : Submodule ℂ hypFock :=
  ⨆ j : ℕ, LinearMap.range (hypSugawara κ ell (-(j + 1 : ℤ)))

/-- Concrete Fock-space light-cone splitting with only the nonzero-momentum hypothesis. -/
theorem hypFock_lightConeSplit (κ ell : ℂ) (hκ : κ ≠ 0) :
    IsCompl (hypNegativeSpurious κ ell) (hypOneFamily 0) := by

/-- The complementary light-cone choice when the other momentum is nonzero. -/
theorem hypFock_lightConeSplit_swapped (κ ell : ℂ) (hell : ell ≠ 0) :
    IsCompl (hypNegativeSpurious κ ell) (hypOneFamily 1) := by

/-- Every nonzero momentum has a concrete choice of longitudinal common kernel. -/
theorem hypFock_lightConeSplit_nonzero (κ ell : ℂ) (h : κ ≠ 0 ∨ ell ≠ 0) :
    ∃ a : Fin 2, IsCompl (hypNegativeSpurious κ ell)
      (⨅ j : ℕ, LinearMap.ker (hypOscillator a (if a = 0 then κ else ell) (j + 1 : ℤ))) := by
\end{leancode}

\begin{dictionary}
\chseventeenL{hypBHomogeneous d w}: polynomials homogeneous of $b$-count $d$ and $b$-level $w$ (weights
\chseventeenL{hypBCountWeight}, \chseventeenL{hypBLevelWeight}: $b_{j+1}$ has count $1$ and level $j+1$, $a$'s have $0$).
\chseventeenL{hypBFiltered d w}: the span of monomials with count $<d$, or count $=d$ and level $<w$.
\chseventeenL{stringNegativeWord D x [i₁, …, i_k]} $=D_{i_1}(\cdots D_{i_k}(x))$ (defined in the Chapter~\ref{ch:18} file
\leanpath{VOA/\allowbreak Moonshine/\allowbreak MoonshineCNStringStraightening.lean}); \chseventeenL{hypCanonicalBWord e} is the $b$-exponent
multiset sorted decreasingly, \chseventeenL{hypAExponent e} the $a$-part.  So \chseventeenL{hypPBWMap} is $\Phi$.  \chseventeenL{IsCompl U W}
means $U\cap W=0$ and $U+W$ is everything.  \chseventeenL{hypPBWEquiv} packages \chseventeenL{hypPBWMap_injective} and
\chseventeenL{hypPBWMap_surjective}; no claim that it is multiplicative.  \textbf{Hypotheses:} only $\kappa\ne0$ (resp.\
$\lambda\ne0$, or one of them); this is the nonzero-momentum condition.  At zero momentum the theorem is false
(for $\kappa=\lambda=0$ the level-one part of the spurious space is $L_{-1}\C=0$, so $b_1$ lies in neither summand), and
Chapter~\ref{ch:18} treats that case separately.
\end{dictionary}

\begin{proof}[Proof on paper]
Order pairs $(d,w)$ (count, level) lexicographically: $(d,w)<(d',w')$ iff $d<d'$, or $d=d'$ and $w<w'$; a polynomial is
``strictly lower than $(d,w)$'' if all its monomials have (count, level) $<(d,w)$.  (i) In $L_{-(j+1)}$, the term $\kappa b_{j+1}$
raises count by one and level by $j+1$.  The term $\lambda a_{j+1}$ and the $a$-part of the derivation preserve count and
level; the $b$-part of the derivation, $b_{k+j+1}\partial/\partial b_k$, preserves count and raises level by $j+1$, which
is still strictly lower because the count is $d<d+1$.  The quadratic terms $a_{j-k}b_{k+1}$ raise count by one but level
only by $k+1\le j$.
(ii) Iterating (i) gives $\Phi(m)=\kappa^km+\text{lower}$ (\chseventeenL{hypSugawara_negative_word_leading}).  Injectivity: if
$\Phi(p)=0$, look at a maximal monomial of $p$ in the order; its coefficient in $\Phi(p)$ is $\kappa^k$ times its
coefficient in $p$.  Surjectivity: by double well-founded induction on count and level, each monomial is
$\kappa^{-k}\Phi(m)$ plus strictly lower terms.  No finite-dimensionality is used.
(iii) The negative modes satisfy the Witt relations $[L_{-(i+1)},L_{-(j+1)}]=(j-i)L_{-(i+j+2)}$
(\chseventeenL{hypSugawara_negative_comm}), so every word straightens into sorted words of the same energy
(\chseventeenL{hypSugawara_negative_straightening}).  Thus $\sum\operatorname{im}L_{-n}$ is the span of nonempty sorted words applied
to $b$-free polynomials, i.e.\ $\Phi$(polynomials with $b$-count $\ge1$); since $\Phi$ is bijective and is the identity on
$b$-free polynomials, the complement is $\chseventeenF^{\mathrm{iso}}(a)$.  The swap $a\leftrightarrow b$ (\chseventeenL{hypSwapPolynomial}) conjugates
$L_n(\kappa,\lambda)$ to $L_n(\lambda,\kappa)$ and gives the case $\lambda\ne0$.
\end{proof}

\begin{physics}
This is the light-cone decomposition at the heart of the no-ghost theorem.  At momentum with $p^+\neq0$ every
state is uniquely a ``longitudinal'' state (no $\alpha^-$ excitations) plus Virasoro descendants: the Virasoro
constraints can be solved for the $\alpha^-$ oscillators.  Brower and Goddard--Thorn use the same triangularity
($L_{-n}\approx p^+\alpha^-_{-n}$).  After tensoring with a $c=24$ matter theory (Theorem~\ref{ch17:thm:tensorlightcone})
and imposing $c=26$, Chapter~\ref{ch:18} deduces that physical states modulo null states are the transverse
states.
\end{physics}

\begin{status}
Proved; hypotheses only $\kappa\ne0$, $\lambda\ne0$ or $\kappa\ne0\vee\lambda\ne0$ as displayed.  Printed in
\leanpath{VOA/\allowbreak Moonshine/\allowbreak MoonshineCNAudit6.lean}; standard axioms.
\end{status}

\paragraph{The remaining \texttt{Hyperbolic*} files.}  \texttt{HyperbolicNegativeModes} (differential identities
for the negative modes: commutators with $\partial/\partial a_t$, with $X$, and with all oscillators),
\texttt{HyperbolicNegativeVirasoro} (the negative Witt relations and straightening),
\texttt{HyperbolicPositiveCreation} (positive modes against creation operators),
\texttt{HyperbolicMixedModes} and \texttt{HyperbolicSugawaraAdjoint} (all signed mode--oscillator commutators;
adjointness), \texttt{HyperbolicCentralScalar} (the finite sums $\sum_{k<j}(k+1)(j-k)=j(j+1)(j+2)/6$ and the vacuum
coefficient), \texttt{HyperbolicFilteredWords} (the iterated leading-term estimate),
\texttt{HyperbolicPBWCoordinates} (the PBW map agrees with every sorted word on every $b$-free base) and
\texttt{HyperbolicSwap} (the family exchange) are the supporting lemmas of
Theorems~\ref{ch17:thm:hypvirasoro}--\ref{ch17:thm:hyppbw}; all are printed in \leanpath{VOA/\allowbreak Moonshine/\allowbreak MoonshineCNAudit5.lean},
\texttt{Audit6} or \texttt{Audit7}.

%======================================================================
\section{The tensor theory \texorpdfstring{$V^\natural\otimes V_{\mathrm{II}_{1,1}}$}{V-natural x V(II(1,1))}}\label{ch17:sec:tensor}
%======================================================================

\subsection{Tensor products of near-VOAs}

\begin{definition}[Tensor product of near-VOAs]\label{ch17:def:tensor}
Let $V$ on $A$ and $W$ on $B$ be near-VOAs.  On $A\otimes B$ put
\begin{equation}\label{ch17:eq:tensorfield}
  Y(a\otimes b,z)(x\otimes y)=\sum_{m\in\Z}\bigl(\text{coefficient of }z^m\text{ in }Y(a,z)x\bigr)\otimes
  \bigl(\text{coefficient of }z^{s-m}\text{ in }Y(b,z)y\bigr)\ \text{ at }z^s,
\end{equation}
i.e.\ $Y(a\otimes b,z)=Y(a,z)\otimes Y(b,z)$; the vacuum $\chseventeenvac\otimes\chseventeenvac$; the conformal vector
$\omega_V\otimes\chseventeenvac+\chseventeenvac\otimes\omega_W$; the central charge $c_V+c_W$; and the grading
$(A\otimes B)_n=\operatorname{span}\{a\otimes b:a\in A_i,\ b\in B_j,\ i+j=n\}$.
\end{definition}

\noindent{\raggedright Files \leanpath{VOA/\allowbreak Moonshine/\allowbreak MoonshineCNTensorFields.lean},
\leanpath{VOA/\allowbreak Moonshine/\allowbreak MoonshineCNTensorConformal.lean},
\leanpath{VOA/\allowbreak Moonshine/\allowbreak MoonshineCNTensorGrading.lean},
\leanpath{VOA/\allowbreak Moonshine/\allowbreak MoonshineCNTensorRealization.lean}.\par}
\begin{leancode}
/-- The state–field map on the algebraic tensor product. -/
noncomputable def nearTensorField (V : NearVOA A) (W : NearVOA B) :
    A ⊗[ℂ] B →ₗ[ℂ] FormalDist (End ℂ (A ⊗[ℂ] B)) :=
  TensorProduct.lift (nearTensorFieldBilinear V W)

theorem nearTensorField_tmul_apply (V : NearVOA A) (W : NearVOA B)
    (a x : A) (b y : B) (s : ℤ) :
    nearTensorField V W (a ⊗ₜ b) s (x ⊗ₜ y) =
      ∑ᶠ m : ℤ, (V.Y a m x) ⊗ₜ[ℂ] (W.Y b (s - m) y) := by

/-- The conformal vector on the tensor carrier. -/
noncomputable def nearTensorConformalVector (V : NearVOA A) (W : NearVOA B) : A ⊗[ℂ] B :=
  V.conformalVector ⊗ₜ W.vacuum + V.vacuum ⊗ₜ W.conformalVector

/-- The span of pure tensors whose two integer weights add to `n`. -/
def nearTensorGrading (V : NearVOA A) (W : NearVOA B) (n : ℤ) :
    Submodule ℂ (A ⊗[ℂ] B) :=
  Submodule.span ℂ {t | ∃ (i j : ℤ) (a : A) (b : B),
    a ∈ V.grading i ∧ b ∈ W.grading j ∧ i + j = n ∧ a ⊗ₜ[ℂ] b = t}

/-- Tensor products of near-VOAs require neither finite weight spaces nor a lower bound. -/
noncomputable def nearTensorVOA (V : NearVOA A) (W : NearVOA B) : NearVOA (A ⊗[ℂ] B) where
  grading := nearTensorGrading V W
  isInternal := nearTensorGrading_isInternal V W
  Y := nearTensorField V W
  vacuum := V.vacuum ⊗ₜ W.vacuum
\end{leancode}

\begin{dictionary}
\chseventeenL{⊗[ℂ]} is the algebraic tensor product; \chseventeenL{a ⊗ₜ b} a pure tensor.  \chseventeenL{nearTensorFieldBilinear} is the
bilinear map $(a,b)\mapsto$ \chseventeenL{nprod} of the first-factor lift of $Y(a,z)$ and the second-factor lift of
$Y(b,z)$; \chseventeenL{TensorProduct.lift} extends it linearly.  \chseventeenL{nearTensorField_tmul_apply} shows what it is:
\eqref{ch17:eq:tensorfield}, with \chseventeenL{∑ᶠ} a finite sum (Mathlib's \chseventeenL{finsum}, which is the sum when the
support is finite --- here it is, by truncation of $Y(a,z)x$ and $Y(b,z)y$).  The coefficient index $s$ is the
power $z^s$; \chseventeenL{V.Y a m x} is the $z^m$-coefficient.  The omitted fields of \chseventeenL{nearTensorVOA} are
\chseventeenL{conformalVector := nearTensorConformalVector V W}, \chseventeenL{centralCharge := V.centralCharge + W.centralCharge}
and the proofs (Theorem~\ref{ch17:thm:tensor}).  \textbf{No hypotheses}: neither factor needs finite weight spaces or
a lower bound.
\end{dictionary}

\begin{physics}
Two decoupled chiral CFTs side by side: fields multiply, stress tensors add, central charges add.  The tensor
product of $\Vnat$ with the light-cone bosons is the chiral matter algebra of Borcherds' string.
\end{physics}

\begin{status}
Definitions; printed in \leanpath{VOA/\allowbreak Moonshine/\allowbreak MoonshineCNAudit3.lean}.
\end{status}

\begin{theorem}[The tensor product is a near-VOA]\label{ch17:thm:tensor}
For near-VOAs $V,W$, the data of Definition~\ref{ch17:def:tensor} form a near-VOA; in particular its weight grading is
internal, its modes have the right weight, its fields are mutually local, translation covariance holds, and
\begin{equation*}
  [L_m,L_n]=(m-n)L_{m+n}+\tfrac{c_V+c_W}{12}(m^3-m)\delta_{m+n,0}.
\end{equation*}
\end{theorem}

\noindent{\raggedright Files \leanpath{VOA/\allowbreak Moonshine/\allowbreak MoonshineCNTensorConformal.lean},
\leanpath{VOA/\allowbreak Moonshine/\allowbreak MoonshineCNTensorGrading.lean}.\par}
\begin{leancode}
/-- The central term is the sum of the two factor central charges. -/
theorem nearTensorField_virasoro (V : NearVOA A) (W : NearVOA B) (m n : ℤ) :
    ⁅mode (nearTensorField V W (nearTensorConformalVector V W)) (m + 1),
      mode (nearTensorField V W (nearTensorConformalVector V W)) (n + 1)⁆ =
      ((m : ℂ) - n) •
        mode (nearTensorField V W (nearTensorConformalVector V W)) (m + n + 1) +
      (if m + n = 0 then
        (((m : ℂ) ^ 3 - m) / 12 * (V.centralCharge + W.centralCharge)) •
          (1 : End ℂ (A ⊗[ℂ] B)) else 0) := by

/-- Eigenvalues of the total `L₀` prove independence without a finite-fiber assumption. -/
theorem nearTensorGrading_isInternal (V : NearVOA A) (W : NearVOA B) :
    DirectSum.IsInternal (nearTensorGrading V W) := by

/-- Every tensor coefficient has the required total-weight shift. -/
theorem nearTensorGrading_mode (V : NearVOA A) (W : NearVOA B)
    {k l : ℤ} (u : A ⊗[ℂ] B) (hu : u ∈ nearTensorGrading V W k)
    (v : A ⊗[ℂ] B) (hv : v ∈ nearTensorGrading V W l) (n : ℤ) :
    mode (nearTensorField V W u) n v ∈ nearTensorGrading V W (k + l - n - 1) := by
\end{leancode}

\begin{dictionary}
Since these are fields of \chseventeenL{nearTensorVOA}, the structure's existence already certifies all near-VOA axioms;
we print the three that carry content.  The central charge is literally \chseventeenL{V.centralCharge + W.centralCharge}.
\end{dictionary}

\begin{proof}[Proof on paper]
$Y(\omega,z)=T_V(z)\otimes1+1\otimes T_W(z)$, so $L_n=L^V_n\otimes1+1\otimes L^W_n$, and the two summands commute; their
Virasoro relations add, including the central terms.  $L_0$ acts on $a\otimes b$ ($a\in A_i$, $b\in B_j$) by $i+j$, so the
pieces $(A\otimes B)_n$ are contained in distinct eigenspaces of one operator, hence independent; they span since
pure tensors of homogeneous elements span.  Locality of $Y(a\otimes b)$ and $Y(a'\otimes b')$: the first-factor and
second-factor lifts commute, and each pair of same-factor lifts is local, so Dong's lemma applies.  Translation:
$L_{-1}=L^V_{-1}\otimes1+1\otimes L^W_{-1}$ and the derivative of a product is given by the Leibniz rule.
\end{proof}

\begin{physics}
$T=T_V+T_W$; central charges of decoupled sectors add.  The internality argument (eigenvalues of $L_0$, no
finiteness) is the mathematical version of ``energy eigenstates of the combined system''.
\end{physics}

\begin{status}
Proved for arbitrary near-VOAs, no hypotheses; printed in \leanpath{VOA/\allowbreak Moonshine/\allowbreak MoonshineCNAudit3.lean}
(\chseventeenL{nearTensorVOA} included); standard axioms.
\end{status}

\subsection{The matter algebra of central charge 26}

\begin{definition}[The matter near-VOA]\label{ch17:def:matter}
$\chseventeenM=\Vnat\otimes\chseventeenVII$, the tensor product of the near-VOA of $\Vnat$ (Chapter~\ref{ch:14},
\chseventeenL{VnatVOA}, forgetting finiteness and CFT type) with $\chseventeenVII$.
\end{definition}

\noindent File \leanpath{VOA/\allowbreak Moonshine/\allowbreak MoonshineCNTensorRealization.lean}.
\begin{leancode}
/-- The matter space in the no-ghost construction. -/
abbrev moonshineMatterSpace := Vnat ⊗[ℂ] LorentzianFockSpace

/-- `V♮ ⊗ V_{II_{1,1}}`, with its full tensor state–field map and total integer weight. -/
noncomputable def moonshineMatterNearVOA : NearVOA moonshineMatterSpace :=
  nearTensorVOA VnatVOA.toNearVOA lorentzianNearVOA
\end{leancode}

\begin{dictionary}
\chseventeenL{Vnat} is the carrier of the Moonshine module (Chapter~\ref{ch:14}); \chseventeenL{VnatVOA :
VertexOperatorAlgebra Vnat} is the named vertex operator algebra of \leanpath{VOA/\allowbreak Moonshine/\allowbreak AutVnat.lean},
defined as \chseventeenL{vnatVOA vnat_twistedLocality}, whose vertex map, vacuum, conformal vector and grading are the
explicit ones (\chseventeenL{VnatVOA_Y}, \chseventeenL{VnatVOA_vacuum}, \chseventeenL{VnatVOA_conformalVector}, \chseventeenL{VnatVOA_grading}, all
\chseventeenL{rfl}) of the theorem \chseventeenL{moonshine}.  So $\chseventeenM$ is built from the actual $\Vnat$, not from an
abstract $c=24$ VOA.
\end{dictionary}

\begin{physics}
The full matter chiral algebra of Borcherds' string: Monster CFT $\otimes$ light-cone bosons.
\end{physics}

\begin{status}
Definitions; constructed with no hypotheses (\leanpath{VOA/\allowbreak Moonshine/\allowbreak MoonshineCNAudit3.lean}).
\end{status}

\begin{theorem}[Central charge 26]\label{ch17:thm:c26}
$\chseventeenM$ has central charge $26$; its conformal vector is $\omega_\natural\otimes\chseventeenvac+\chseventeenvac\otimes\omega_{\chseventeenII}$;
and $Y(a\otimes b,z)(x\otimes y)=Y(a,z)x\otimes Y(b,z)y$ coefficientwise.
\end{theorem}

\noindent File \leanpath{VOA/\allowbreak Moonshine/\allowbreak MoonshineCNTensorRealization.lean}.
\begin{leancode}
/-- The central charge is exactly `24 + 2 = 26`, without hypotheses. -/
theorem moonshineMatterNearVOA_centralCharge : moonshineMatterNearVOA.centralCharge = 26 := by

theorem moonshineMatterNearVOA_conformalVector :
    moonshineMatterNearVOA.conformalVector =
      VnatVOA.conformalVector ⊗ₜ lorentzianNearVOA.vacuum +
        VnatVOA.vacuum ⊗ₜ lorentzianNearVOA.conformalVector := rfl

theorem moonshineMatterNearVOA_tmul_field (a x : Vnat) (b y : LorentzianFockSpace) (s : ℤ) :
    moonshineMatterNearVOA.Y (a ⊗ₜ b) s (x ⊗ₜ y) =
      ∑ᶠ m : ℤ, (VnatVOA.Y a m x) ⊗ₜ[ℂ] (lorentzianNearVOA.Y b (s - m) y) :=
\end{leancode}

\begin{dictionary}
\textbf{How the $26$ arises in Lean.} The proof first rewrites the central charge as \chseventeenL{(24 : ℂ) +
lorentzianNearVOA.centralCharge} by \chseventeenL{change}, i.e.\ \emph{definitionally}: the central charge of
\chseventeenL{VnatVOA} is the literal $24$ of the Moonshine construction; then Theorem~\ref{ch17:thm:realization} gives $2$.
\textbf{Fidelity.} The statement is about the actual structure, not about an operator family with the right
commutators; the conformal vector equation is \chseventeenL{rfl}, i.e.\ true by the definition.
\end{dictionary}

\begin{proof}[Proof on paper]
$c=24+2$ by Theorem~\ref{ch17:thm:tensor}.
\end{proof}

\begin{physics}
$26=24+2$: the Monster CFT replaces the $24$ transverse bosons of the critical bosonic string, and the two
light-cone bosons complete the critical dimension.  This is the input $c=26$ of the no-ghost theorem.
\end{physics}

\begin{status}
Proved with no hypotheses; printed in \leanpath{VOA/\allowbreak Moonshine/\allowbreak MoonshineCNAudit3.lean}; standard axioms.
\end{status}

\begin{theorem}[The $\chseventeenII$-momentum grading of $\chseventeenM$]\label{ch17:thm:mattermomentum}
Let $\chseventeenM^\alpha=\Vnat\otimes(\chseventeenF_2\otimes e^\alpha)$.  Then $\chseventeenM=\bigoplus_{\alpha\in\chseventeenII}\chseventeenM^\alpha$
(internal), and $u_{(n)}v\in\chseventeenM^{\alpha+\beta}$ for $u\in\chseventeenM^\alpha$, $v\in\chseventeenM^\beta$, $n\in\Z$.
\end{theorem}

\noindent File \leanpath{VOA/\allowbreak Moonshine/\allowbreak MoonshineCNTensorMomentum.lean}.
\begin{leancode}
/-- Momentum comes from the lattice factor; all of `V♮` has zero momentum. -/
noncomputable def moonshineMatterMomentum (α : lorentzianLattice.carrier) : Submodule ℂ moonshineMatterSpace :=
  TensorVA.tenSub (⊤ : Submodule ℂ Vnat) (latticeMomentumGrading lorentzianLattice α)

theorem moonshineMatterMomentum_isInternal : DirectSum.IsInternal moonshineMatterMomentum := by

/-- Every actual tensor vertex mode adds its two lattice momenta. -/
theorem moonshineMatterMomentum_mode (α β : lorentzianLattice.carrier)
    (u v : moonshineMatterSpace) (hu : u ∈ moonshineMatterMomentum α)
    (hv : v ∈ moonshineMatterMomentum β) (n : ℤ) :
    mode (moonshineMatterNearVOA.Y u) n v ∈ moonshineMatterMomentum (α + β) := by
\end{leancode}

\begin{dictionary}
\chseventeenL{TensorVA.tenSub P Q} is the span of $p\otimes q$ with $p\in P$, $q\in Q$; with $P=\top=\Vnat$ this is
$\Vnat\otimes\chseventeenF_2e^\alpha$.  Internality is proved via the eigenspaces of the operator multiplying $e^\alpha$ by the
injective label $\alpha_0+i\alpha_1\in\C$ (\chseventeenL{lorentzianChargeLabel}); this is a proof device, not an extra datum.
\end{dictionary}

\begin{proof}[Proof on paper]
Tensor the decomposition of Theorem~\ref{ch17:thm:momentum} with $\Vnat$; for modes, use
\eqref{ch17:eq:tensorfield} and Theorem~\ref{ch17:thm:momentum} on the second factor.
\end{proof}

\begin{physics}
Spacetime momentum conservation.  This grading descends to the $\chseventeenII$-grading of the Monster Lie algebra
(\chseventeenL{monsterLieMomentum}, Chapter~\ref{ch:18}), whose pieces are the root spaces.
\end{physics}

\begin{status}
Proved; printed in \leanpath{VOA/\allowbreak Moonshine/\allowbreak MoonshineCNAudit3.lean}; standard axioms.
\end{status}

\subsection{Bilinear forms}

The no-ghost theorem needs a bilinear form on $\chseventeenM$ for which $L_n$ is adjoint to $L_{-n}$ and which is
nondegenerate; the Lie algebra of Chapter~\ref{ch:18} needs, in addition, full FHL invariance (treated in
\S\ref{ch17:sec:fhl}).

Four bilinear forms occur in this chapter; the following table says how they differ.
\begin{center}
\small
\begin{tabular}{p{0.19\textwidth}p{0.15\textwidth}p{0.17\textwidth}p{0.17\textwidth}p{0.15\textwidth}}
\toprule
form & space & adjoint of an oscillator & pairs momenta & used for \\ \midrule
$B_{\mathcal F}$ (Fock form, Ch.~\ref{ch:04}) & $\chseventeenF_r$ & $h(n)^\dagger=-h(-n)$ & --- & building $B_L$ \\
$B_L$ (lattice form) & $V_L$ & $h(n)^\dagger=-h(-n)$ & $\alpha$ with $-\alpha$ & $B$, FHL invariance \\
$B=B_\natural\otimes B_L$ & $\chseventeenM$ & as $B_L$ & $\alpha$ with $-\alpha$ & Monster Lie form \\
$B^{\mathrm{iso}}$ (isotropic form) & $\chseventeenF^{\mathrm{iso}}$ (one sector) & $\varphi(n)^\dagger=\varphi(-n)$ & within the sector & no-ghost theorem \\
\bottomrule
\end{tabular}
\end{center}
All four make $L_n$ adjoint to $L_{-n}$; only $B$ is shown to be FHL-invariant (\S\ref{ch17:sec:fhl}).  $B_L$ and
$B^{\mathrm{iso}}$ are related through momentum reversal (Theorem~\ref{ch17:thm:momentumform}) and a change of variables.

\begin{definition}[Virasoro-adjoint form of a near-VOA]\label{ch17:def:nearform}
A \emph{near-invariant form} on a near-VOA $V$ is a bilinear form $B$ with $B(L_na,b)=B(a,L_{-n}b)$ for all $n,a,b$.
The tensor product of two near-invariant forms is a near-invariant form on the tensor near-VOA.
\end{definition}

\noindent{\raggedright Files \leanpath{VOA/\allowbreak Moonshine/\allowbreak MoonshineCNLatticeForm.lean},
\leanpath{VOA/\allowbreak Moonshine/\allowbreak MoonshineCNTensorForm.lean}.\par}
\begin{leancode}
/-- The existing invariant-form convention, extended to a near-VOA. -/
structure NearInvariantForm {M : Type*} [AddCommGroup M] [Module ℂ M] (V : NearVOA M) where
  form : LinearMap.BilinForm ℂ M
  virasoro_adj : ∀ (n : ℤ) (a b : M),
    form (nearVirasoroMode V n a) b = form a (nearVirasoroMode V (-n) b)

/-- Tensoring Virasoro-adjoint forms preserves every adjointness relation. -/
noncomputable def nearTensorInvariantForm {V : NearVOA A} {W : NearVOA B}
    (F : NearInvariantForm V) (G : NearInvariantForm W) : NearInvariantForm (nearTensorVOA V W) where
  form := F.form.tmul G.form
  virasoro_adj := nearTensorForm_virasoro_adjoint F G
\end{leancode}

\begin{dictionary}
\chseventeenL{nearVirasoroMode V n} $=$ \chseventeenL{mode (V.Y V.conformalVector) (n + 1)} $=L_n$.  \textbf{The name is weaker than it
sounds.} Exactly like the library's \chseventeenL{VertexOperatorAlgebra.InvariantForm} (Chapter~\ref{ch:L04},
Definition~\ref{l04:def:invform}), \chseventeenL{NearInvariantForm} records \emph{only} Virasoro adjointness, the $u=\omega$ slice
of FHL invariance; the module docstring says so (``Full FHL invariance is a separate property'').  A referee
must not read \chseventeenL{NearInvariantForm} as FHL invariance.  \chseventeenL{F.form.tmul G.form} is the form
$(a\otimes b,a'\otimes b')\mapsto F(a,a')G(b,b')$.
\end{dictionary}

\begin{physics}
$L_n^\dagger=L_{-n}$: the BPZ conjugation of the stress tensor.  This is the only property of the inner product
that the Virasoro-algebraic part of the no-ghost theorem uses.
\end{physics}

\begin{status}
Definitions; printed in \leanpath{VOA/\allowbreak Moonshine/\allowbreak MoonshineCNAudit3.lean}.
\end{status}

\begin{definition}[The lattice and matter forms]\label{ch17:def:matterform}
On $\chseventeenVII$ the form is the library's lattice form $B_L$ of Chapter~\ref{ch:04}
(Definition~\ref{ch04:def:latform}): $B_L(p\otimes e^\alpha,q\otimes e^\beta)=\delta_{\alpha+\beta,0}\,\chi(\alpha)\,B_{\mathcal F}(p,q)$, with $B_{\mathcal F}$ the
Fock form (oscillators skew-adjoint, $h(n)^\dagger=-h(-n)$) and $\chi(\alpha)=(-1)^{\langle\alpha,\alpha\rangle/2}\varepsilon(\alpha,\alpha)$.  On
$\chseventeenM$ it is $B_\natural\otimes B_L$, with $B_\natural$ the normalised invariant form \chseventeenL{natForm} of $\Vnat$
(Chapter~\ref{ch:14}).
\end{definition}

\noindent{\raggedright Files \leanpath{VOA/\allowbreak Moonshine/\allowbreak MoonshineCNLatticeForm.lean},
\leanpath{VOA/\allowbreak Moonshine/\allowbreak MoonshineCNTensorForm.lean}, \leanpath{VOA/\allowbreak Moonshine/\allowbreak NondegLattice.lean}.\par}
\begin{leancode}
noncomputable def latForm : LinearMap.BilinForm ℂ (LatModuleGen L) :=
  (uForm (rank L)).tmul (chForm L)

/-- The explicit Fock form paired across opposite lattice charges. -/
noncomputable def lorentzianBilinearForm : LinearMap.BilinForm ℂ LorentzianFockSpace :=
  latForm lorentzianLattice

noncomputable def lorentzianInvariantForm : NearInvariantForm lorentzianNearVOA where
  form := lorentzianBilinearForm
  virasoro_adj := lorentzianBilinearForm_virasoro_adjoint

/-- The product of the established moonshine form with the explicit Lorentzian form. -/
noncomputable def moonshineMatterForm : LinearMap.BilinForm ℂ moonshineMatterSpace :=
  natForm.tmul lorentzianBilinearForm

noncomputable def moonshineMatterInvariantForm : NearInvariantForm moonshineMatterNearVOA where
  form := moonshineMatterForm
  virasoro_adj := moonshineMatterForm_virasoro_adjoint
\end{leancode}

\begin{dictionary}
\chseventeenL{latForm} lives in \leanpath{VOA/\allowbreak Moonshine/\allowbreak NondegLattice.lean} (Chapter~\ref{ch:04}); it is defined for any integral
lattice and its definition never used positivity.  \chseventeenL{chForm L} pairs $e^\alpha$ with $e^{-\alpha}$ with the sign
$\chi(\alpha)=$ \chseventeenL{chSign L α} (Proposition~\ref{ch04:prop:chsign}).  The form pairs \emph{opposite} momenta; it vanishes on
two states of the same nonzero momentum (Theorem~\ref{ch17:thm:momentumform}).
\end{dictionary}

\begin{physics}
The BPZ pairing of a charged state with its charge conjugate: $\langle e^{-\alpha}|e^\alpha\rangle\neq0$.  For the matter theory
this is the product of the Monster CFT's Zamolodchikov metric and the light-cone pairing.
\end{physics}

\begin{status}
Definitions; printed in \leanpath{VOA/\allowbreak Moonshine/\allowbreak MoonshineCNAudit3.lean}.
\end{status}

\begin{theorem}[Properties of the matter form]\label{ch17:thm:matterform}
$B_L$ and $B_\natural\otimes B_L$ are symmetric, normalised ($B(\chseventeenvac,\chseventeenvac)=1$), nondegenerate (every nonzero $x$ pairs
nontrivially with some $y$), and satisfy $B(L_nx,y)=B(x,L_{-n}y)$ for all $n\in\Z$.
\end{theorem}

\noindent{\raggedright Files \leanpath{VOA/\allowbreak Moonshine/\allowbreak MoonshineCNLatticeForm.lean},
\leanpath{VOA/\allowbreak Moonshine/\allowbreak MoonshineCNTensorForm.lean},
\leanpath{VOA/\allowbreak Moonshine/\allowbreak MoonshineCNTensorNondegenerate.lean}.\par}
\begin{leancode}
theorem lorentzianBilinearForm_symmetric (x y : LorentzianFockSpace) :
    lorentzianBilinearForm x y = lorentzianBilinearForm y x := latForm_symm _ x y

theorem lorentzianBilinearForm_nondegenerate (x : LorentzianFockSpace) (hx : x ≠ 0) :
    ∃ y, lorentzianBilinearForm x y ≠ 0 := latForm_nondeg _ x hx

theorem lorentzianBilinearForm_normalized :
    lorentzianBilinearForm lorentzianNearVOA.vacuum lorentzianNearVOA.vacuum = 1 := by

theorem lorentzianBilinearForm_virasoro_adjoint (n : ℤ) (x y : LorentzianFockSpace) :
    lorentzianBilinearForm (nearVirasoroMode lorentzianNearVOA n x) y =
      lorentzianBilinearForm x (nearVirasoroMode lorentzianNearVOA (-n) y) := by

theorem moonshineMatterForm_symmetric (x y : moonshineMatterSpace) :
    moonshineMatterForm x y = moonshineMatterForm y x := by

theorem moonshineMatterForm_normalized :
    moonshineMatterForm moonshineMatterNearVOA.vacuum moonshineMatterNearVOA.vacuum = 1 := by

theorem moonshineMatterForm_virasoro_adjoint (n : ℤ) (x y : moonshineMatterSpace) :
    moonshineMatterForm (nearVirasoroMode moonshineMatterNearVOA n x) y =
      moonshineMatterForm x (nearVirasoroMode moonshineMatterNearVOA (-n) y) :=

theorem moonshineForm_nondegenerate (x : Vnat) (hx : x ≠ 0) : ∃ y, natForm x y ≠ 0 := by

theorem moonshineMatterForm_nondegenerate (x : moonshineMatterSpace) (hx : x ≠ 0) :
    ∃ y, moonshineMatterForm x y ≠ 0 := by
\end{leancode}

\begin{dictionary}
Nondegeneracy is stated in the strong ``global'' form on the infinite-dimensional space: for every $x\ne0$ there is a
$y$ (for $x$ homogeneous of momentum $\alpha$ and weight $N$, one can take $y$ of momentum $-\alpha$ and weight $N$) with $B(x,y)\ne0$.  \chseventeenL{moonshineForm_nondegenerate} comes
from the simplicity of $\Vnat$ (the radical of \chseventeenL{natInvForm} is zero, \chseventeenL{natInvForm_leftRadical}).  Weight-wise
nondegeneracy on each finite-dimensional bigraded piece is not stated here (Chapter~\ref{ch:18} uses the isotropic
model's \chseventeenL{hypForm_weight_nondegenerate} instead).  \textbf{Not claimed here}: positivity (meaningless as it stands, the form being
indefinite) and full FHL invariance (\S\ref{ch17:sec:fhl}).
\end{dictionary}

\begin{proof}[Proof on paper]
Symmetry and nondegeneracy of $B_L$: the monomial--charge basis $x^d\otimes e^\alpha$ is paired with $x^d\otimes e^{-\alpha}$ by the
nonzero number $u(d)\chi(\alpha)$ and with nothing else, so the Gram matrix is a signed permutation matrix
(\chseventeenL{lorentzianStateBasis_pair}).  For the tensor product, expand $x=\sum_{(d,\alpha)}x_{d,\alpha}\otimes(x^d e^\alpha)$ with
$x_{d,\alpha}\in\Vnat$; if $x_{d,\alpha}\ne0$, pick $a$ with $B_\natural(x_{d,\alpha},a)\ne0$ and pair $x$ with $a\otimes x^de^{-\alpha}$
(\chseventeenL{moonshineMatterForm_coefficient}).  Virasoro adjointness of $B_L$: on each sector the Fock form makes the Fock
$L_n$ adjoint to $L_{-n}$ (\chseventeenL{uForm_virasoro_adjoint}, \S\ref{ch17:sec:fhl}), and the momentum terms match because the
paired sectors have momenta $\alpha,-\alpha$ and the frame oscillator of $-\alpha$ is minus that of $\alpha$.  For the tensor form,
$L_n=L^\natural_n\otimes1+1\otimes L^{\chseventeenII}_n$ and both factors are adjoint.
\end{proof}

\begin{physics}
A nondegenerate BPZ pairing with $L_n^\dagger=L_{-n}$ on the matter theory: exactly the input of the no-ghost theorem.
It is a complex bilinear form that pairs momentum $\alpha$ with $-\alpha$ and vanishes on two states of the same nonzero
momentum, so ``positive'' or ``indefinite'' have no meaning for it directly.  The ghost story is told by the
\emph{sector} form, obtained by composing with momentum reversal (Theorem~\ref{ch17:thm:momentumform}) and restricted to
real combinations; in light-cone variables it is the isotropic form $B^{\mathrm{iso}}$ of \S\ref{ch17:sec:hyp}, which is
not positive.  Those negative-norm states are the ghosts that the no-ghost theorem shows to decouple at $c=26$.
\end{physics}

\begin{status}
Proved with no hypotheses; all printed in \leanpath{VOA/\allowbreak Moonshine/\allowbreak MoonshineCNAudit3.lean}; standard axioms.
\end{status}

\subsection{The light-cone splitting of a matter sector}

The sector of $\chseventeenM$ of momentum $(x,y)$ will be identified (Theorem~\ref{ch17:thm:sectorequiv}) with
$\Vnat\otimes\chseventeenF^{\mathrm{iso}}$.  The splitting of \S\ref{ch17:sec:hyp} extends to $V\otimes\chseventeenF^{\mathrm{iso}}$ for an
arbitrary vector space $V$ with operators $D_j$ (later $D_j=L^V_{-(j+1)}$).

\begin{definition}[Tensor negative modes]\label{ch17:def:tensornegative}
For $D:\N\to\End(V)$ and $\kappa,\lambda\in\C$, $\mathcal L_{-(j+1)}=D_j\otimes1+1\otimes L_{-(j+1)}$ on $V\otimes\chseventeenF^{\mathrm{iso}}$; the
tensor one-family space is $V\otimes\chseventeenF^{\mathrm{iso}}(\varphi)$; the spurious space is $\sum_j\operatorname{im}\mathcal L_{-(j+1)}$.
\end{definition}

\noindent{\raggedright Files \leanpath{VOA/\allowbreak Moonshine/\allowbreak MoonshineCNTensorNegative.lean},
\leanpath{VOA/\allowbreak Moonshine/\allowbreak MoonshineCNTensorLightCone.lean}.\par}
\begin{leancode}
/-- Index `j` denotes the genuine tensor mode with integer index `-(j+1)`. -/
def hypTensorNegative (D : ℕ → Module.End ℂ V) (κ ell : ℂ) (j : ℕ) :
    Module.End ℂ (V ⊗[ℂ] hypFock) :=
  (D j).rTensor hypFock + (hypSugawara κ ell (-(j + 1 : ℤ))).lTensor V

def hypTensorOneFamily (a : Fin 2) : Submodule ℂ (V ⊗[ℂ] hypFock) :=
  hypTensorSupported {e | monomial e 1 ∈ hypOneFamily a}

def hypTensorNegativeSpurious (D : ℕ → Module.End ℂ V) (κ ell : ℂ) :
    Submodule ℂ (V ⊗[ℂ] hypFock) := ⨆ j : ℕ, LinearMap.range (hypTensorNegative D κ ell j)
\end{leancode}

\begin{dictionary}
\chseventeenL{hypTensorSupported s}: elements $\sum_ev_e\otimes x^e$ whose polynomial exponents $e$ all lie in $s$ (coordinates via
the monomial basis, \chseventeenL{hypTensorCoordinates}); so \chseventeenL{hypTensorOneFamily a} $=V\otimes\chseventeenF^{\mathrm{iso}}(\varphi)$.  $V$ is an
arbitrary complex vector space (section variable); no grading, form or basis is assumed.
\end{dictionary}

\begin{physics}
$\mathcal L_{-n}=L^{\mathrm{matter}}_{-n}+L^{\mathrm{lc}}_{-n}$: the total Virasoro lowering operators of the string, matter plus
light-cone.
\end{physics}

\begin{status}
Definitions; printed in \leanpath{VOA/\allowbreak Moonshine/\allowbreak MoonshineCNAudit7.lean}.
\end{status}

\begin{theorem}[Tensor PBW map and light-cone splitting]\label{ch17:thm:tensorlightcone}
(i) For any $D$, if $\kappa\ne0$ the sorted-word map $v\otimes a^\alpha b_\lambda\mapsto\mathcal L_{-\lambda_1}\cdots\mathcal L_{-\lambda_k}(v\otimes a^\alpha)$ is a
linear automorphism of $V\otimes\chseventeenF^{\mathrm{iso}}$.  (ii) If moreover $[D_i,D_j]=(j-i)D_{i+j+1}$, then
$V\otimes\chseventeenF^{\mathrm{iso}}=\bigl(\sum_j\operatorname{im}\mathcal L_{-(j+1)}\bigr)\oplus V\otimes\chseventeenF^{\mathrm{iso}}(a)$; similarly with $b$ when $\lambda\ne0$; and for
$(\kappa,\lambda)\ne0$ one of the two common kernels $\bigcap_j\ker(1\otimes \varphi(j))$ complements the spurious space.
\end{theorem}

\noindent{\raggedright Files \leanpath{VOA/\allowbreak Moonshine/\allowbreak MoonshineCNTensorPBWEquiv.lean},
\leanpath{VOA/\allowbreak Moonshine/\allowbreak MoonshineCNTensorLightCone.lean},
\leanpath{VOA/\allowbreak Moonshine/\allowbreak MoonshineCNTensorLightConeKernels.lean}.\par}
\begin{leancode}
/-- No negative Virasoro relation is needed for injectivity. -/
theorem hypTensorPBWMap_injective (D : ℕ → Module.End ℂ V) (κ ell : ℂ) (hκ : κ ≠ 0) :
    Function.Injective (hypTensorPBWMap D κ ell) := by

/-- Surjectivity holds for arbitrary matter operators. -/
theorem hypTensorPBWMap_surjective (D : ℕ → Module.End ℂ V) (κ ell : ℂ) (hκ : κ ≠ 0) :
    Function.Surjective (hypTensorPBWMap D κ ell) := by

/-- Concrete tensor-sector splitting, assuming only the matter negative relation. -/
theorem hypTensor_lightConeSplit (D : ℕ → Module.End ℂ V)
    (hD : ∀ i j, D i * D j - D j * D i = ((j : ℂ) - i) • D (i + j + 1))
    (κ ell : ℂ) (hκ : κ ≠ 0) :
    IsCompl (hypTensorNegativeSpurious D κ ell) (hypTensorOneFamily 0) := by

theorem hypTensorOneFamily_eq_commonKernel (a : Fin 2) (charge : ℂ) :
    hypTensorOneFamily (V := V) a =
      ⨅ j : ℕ, LinearMap.ker ((hypOscillator a charge (j + 1 : ℤ)).lTensor V) := by

/-- For every nonzero momentum, the genuine negative tensor ranges complement
one of the actual longitudinal common kernels. -/
theorem hypTensor_lightConeSplit_nonzero (D : ℕ → Module.End ℂ V)
    (hD : ∀ i j, D i * D j - D j * D i = ((j : ℂ) - i) • D (i + j + 1))
    (κ ell : ℂ) (h : κ ≠ 0 ∨ ell ≠ 0) :
    ∃ a : Fin 2, IsCompl (hypTensorNegativeSpurious D κ ell)
      (⨅ j : ℕ, LinearMap.ker
        ((hypOscillator a (if a = 0 then κ else ell) (j + 1 : ℤ)).lTensor V)) := by
\end{leancode}

\begin{dictionary}
\textbf{The hypothesis \chseventeenL{hD}} is the Witt relation for the negative matter modes: with
$D_j=L^V_{-(j+1)}$, $[L_{-(i+1)},L_{-(j+1)}]=(j-i)L_{-(i+j+2)}$, which is the Virasoro relation for negative modes (no central
term since $m+n\ne0$).  It is a genuine hypothesis here because $V$ and $D$ are arbitrary; for $\Vnat$ it is discharged in
Chapter~\ref{ch:18} (\chseventeenL{vnatHypTensor_lightConeSplit} assumes nothing).  Injectivity and surjectivity of the PBW map need
\emph{no} relation on $D$ (the matter term preserves both filtration indices); only the identification of its image
with the spurious space uses \chseventeenL{hD}, through straightening.  \chseventeenL{hypTensorPBWEquiv} packages (i).
\end{dictionary}

\begin{proof}[Proof on paper]
Write elements as polynomials with coefficients in $V$.  $\mathcal L_{-(j+1)}(v\otimes p)=\kappa\,v\otimes b_{j+1}p+$(terms of lower
order), where $D_jv\otimes p$ has the \emph{same} count and level as $v\otimes p$, hence is lower than the leading term.  The
proofs of Theorem~\ref{ch17:thm:hyppbw} go through coefficientwise with values in $V$; straightening uses \chseventeenL{hD}.
\end{proof}

\begin{physics}
At nonzero momentum, every state of the string Fock space ($\text{matter}\otimes\text{light-cone}$) is uniquely
``longitudinal-free'' plus total Virasoro descendants.  This is the decomposition $\mathcal H=\mathcal S\oplus\mathcal K$ of
Goddard--Thorn; the matter theory enters only through its negative Virasoro modes.
\end{physics}

\begin{status}
Proved with the displayed hypotheses only; printed in \leanpath{VOA/\allowbreak Moonshine/\allowbreak MoonshineCNAudit7.lean}; standard axioms.
The remaining tensor files (\texttt{TensorNegative} coordinates and filtration, \texttt{TensorPBWMap} leading terms,
\texttt{TensorLightConeKernels} swap) are the supporting lemmas.
\end{status}

%======================================================================
\section{Vertex algebra theory for near-VOAs}\label{ch17:sec:near}
%======================================================================

The library's general structure theory (Goddard uniqueness, translation covariance, the $n$-th product and
commutator formulas, associativity, the Borcherds identity, skew-symmetry: Chapter~\ref{ch:L03}) is stated for
\chseventeenL{VertexOperatorAlgebra}, whose weight spaces are finite-dimensional and bounded below.  The proofs, however, use
only locality, the vacuum axioms, regularity and translation covariance.  The seven files
\texttt{MoonshineCNNear\{Goddard, Translation, NthProduct, Iterate, Borcherds, Skew, Consequences\}} (plus the
specialisations in \texttt{NearMatterIdentities}) port them to \chseventeenL{NearVOA}, ``without finiteness assumptions''.
All declarations below are in namespace \chseventeenL{VOA.NearVOA} and take a near-VOA \chseventeenL{V : NearVOA 𝒱} (or
\chseventeenL{H}) as their first explicit argument; no other hypothesis is added unless displayed.

\begin{definition}[Products, Virasoro modes, spurious states]\label{ch17:def:nearops}
For a near-VOA: $u_{(n)}v=$ \chseventeenL{V.smul u n v}; $L_n=\omega_{(n+1)}$; $T=L_{-1}$; the \emph{spurious subspace} is
$\sum_{n\ge1}\operatorname{im}L_{-n}$.
\end{definition}

\noindent{\raggedright Files \leanpath{VOA/\allowbreak Moonshine/\allowbreak MoonshineCNNearGoddard.lean},
\leanpath{VOA/\allowbreak Moonshine/\allowbreak MoonshineCNNearConsequences.lean}.\par}
\begin{leancode}
noncomputable def smul (V : NearVOA 𝒱) (u : 𝒱) (n : ℤ) (v : 𝒱) : 𝒱 := mode (V.Y u) n v

noncomputable def T (V : NearVOA 𝒱) : End ℂ 𝒱 := V.L (-1)

/-- The subspace generated by negative Virasoro modes. -/
noncomputable def spurious (V : NearVOA 𝒱) : Submodule ℂ 𝒱 :=
  ⨆ n : ℕ, LinearMap.range (V.L (-(n + 1 : ℤ)))
\end{leancode}

\begin{dictionary}
\chseventeenL{V.L n} is defined in the same file as \chseventeenL{mode (V.Y V.conformalVector) (n + 1)} (lemma \chseventeenL{L_def}); it coincides
definitionally with \chseventeenL{nearVirasoroMode V n} of \S\ref{ch17:sec:tensor}.  \chseventeenL{⨆ n : ℕ} over $n\ge0$ of the range of
$L_{-(n+1)}$ is $\sum_{m\ge1}\operatorname{im}L_{-m}$.
\end{dictionary}

\begin{physics}
The spurious states are the Virasoro descendants: in string theory, the states of the form $L_{-1}\chi_1+L_{-2}\chi_2+\cdots$
that decouple from physical amplitudes when they are also physical.
\end{physics}

\begin{status}
Definitions; printed in \leanpath{VOA/\allowbreak Moonshine/\allowbreak MoonshineCNAudit9.lean}.
\end{status}

\begin{theorem}[Goddard uniqueness for near-VOAs]\label{ch17:thm:goddard}
If a formal distribution $A$ is local with every $Y(b,z)$ and $A(z)\chseventeenvac=Y(a,z)\chseventeenvac$, then $A=Y(a,z)$.
\end{theorem}

\noindent File \leanpath{VOA/\allowbreak Moonshine/\allowbreak MoonshineCNNearGoddard.lean}.
\begin{leancode}
/-- Ordinary locality and the vacuum tower determine a vertex operator. -/
theorem goddard_uniqueness (V : NearVOA 𝒱) (A : FormalDist (End ℂ 𝒱)) (a : 𝒱)
    (hAloc : ∀ b : 𝒱, Local A (V.Y b))
    (hAvac : ∀ m : ℤ, A m V.vacuum = (V.Y a) m V.vacuum) : A = V.Y a := by
\end{leancode}

\begin{dictionary}
\chseventeenL{A m V.vacuum} is the $z^m$-coefficient of $A(z)\chseventeenvac$.  The file also proves the sign-generic pointwise version
\chseventeenL{goddard_pointwiseε} (for $\varepsilon$-locality) and Dong's lemma for near-VOA fields (\chseventeenL{dong_fields}).
Same statement as Chapter~\ref{ch:L03}, Theorem~\ref{l03:thm:goddard}, for the larger class of near-VOAs.
\end{dictionary}

\begin{proof}[Proof on paper]
For each $b$, $(z_1-z_2)^N[A(z_1)-Y(a,z_1)]Y(b,z_2)\chseventeenvac=(z_1-z_2)^NY(b,z_2)[A(z_1)-Y(a,z_1)]\chseventeenvac=0$ by locality; setting
$z_2=0$ (allowed since $Y(b,z_2)\chseventeenvac$ has no negative powers) gives $z_1^N(A(z_1)-Y(a,z_1))b=0$, so
$A(z_1)b=Y(a,z_1)b$.
\end{proof}

\begin{physics}
A field of a CFT is determined by the state it creates and by its OPEs: two local fields creating the same state are
equal.
\end{physics}

\begin{status}
Proved; printed in \leanpath{VOA/\allowbreak Moonshine/\allowbreak MoonshineCNAudit9.lean}; standard axioms.
\end{status}

\begin{theorem}[Translation, products, commutator, iterate]\label{ch17:thm:nearbasic}
For a near-VOA and all $u,v,x$, $m,n,\ell\in\Z$, $k\in\N$:
(i) $[L_{-1},u_{(n)}]=-n\,u_{(n-1)}$;
(ii) the $k$-th product field of $Y(u,z)$ and $Y(v,z)$ is $Y(u_{(k)}v,z)$;
(iii) $[u_{(m)},v_{(n)}]=\sum_{j\ge0}\binom mj(u_{(j)}v)_{(m+n-j)}$, a finite sum;
(iv) $(u_{(\ell)}v)_{(N)}x=\sum_{i\ge0}(-1)^i\binom\ell i\bigl(u_{(\ell-i)}v_{(N+i)}x-(-1)^\ell v_{(\ell+N-i)}u_{(i)}x\bigr)$.
\end{theorem}

\noindent{\raggedright Files \leanpath{VOA/\allowbreak Moonshine/\allowbreak MoonshineCNNearTranslation.lean},
\leanpath{VOA/\allowbreak Moonshine/\allowbreak MoonshineCNNearNthProduct.lean},
\leanpath{VOA/\allowbreak Moonshine/\allowbreak MoonshineCNNearConsequences.lean},
\leanpath{VOA/\allowbreak Moonshine/\allowbreak MoonshineCNNearIterate.lean}.\par}
\begin{leancode}
theorem L_comm_mode (V : NearVOA 𝒱) (u : 𝒱) (n : ℤ) :
    ⁅V.L (-1), mode (V.Y u) n⁆ = (-n) • mode (V.Y u) (n - 1) := by

theorem nth_product (V : NearVOA 𝒱) (u v : 𝒱) (k : ℕ) :
    cCoeff k (bracketDistε 1 (V.Y u) (V.Y v)) = V.Y (V.smul u k v) := by

theorem mode_commutator_lie (V : NearVOA 𝒱) (u v : 𝒱) (m n : ℤ) :
    ∃ N : ℕ, ⁅mode (V.Y u) m, mode (V.Y v) n⁆ =
      ∑ j : Fin N, gBinom m (j : ℕ) •
        mode (V.Y (V.smul u (j : ℕ) v)) (m + n - (j : ℕ)) := by

theorem mode_iterate (V : NearVOA 𝒱) (u v : 𝒱) (ℓ N : ℤ) (x : 𝒱) :
    V.smul (V.smul u ℓ v) N x
      = ∑ᶠ i : ℕ, ((-1 : ℂ) ^ i * gBinom ℓ i) •
          (V.smul u (ℓ - i) (V.smul v (N + i) x)
            - ((1 : ℂ) * (-1 : ℂ) ^ ℓ) • V.smul v (ℓ + N - i) (V.smul u i x)) := by
\end{leancode}

\begin{dictionary}
\chseventeenL{gBinom m j} is the generalised binomial coefficient $\binom mj=m(m-1)\cdots(m-j+1)/j!$ for $m\in\Z$.
\chseventeenL{bracketDistε 1 A B} is the commutator distribution $[A(z_1),B(z_2)]$ and \chseventeenL{cCoeff k} extracts the coefficient of
$\partial^{(k)}\delta(z_1-z_2)$ in its delta-expansion; (ii) says that this coefficient is the field of $u_{(k)}v$ (the OPE pole
coefficient).  In (iii), \chseventeenL{∃ N} together with \chseventeenL{∑ j : Fin N} is the finite sum over $0\le j<N$ ($N$ depends on
$u,v$ via regularity).  In (iv), \chseventeenL{∑ᶠ} is a finite sum on the fixed vector $x$.  These are the statements of
Chapter~\ref{ch:L03}, Theorems~\ref{l03:thm:transcov}, \ref{l03:thm:nthproduct}, \ref{l03:thm:comm} and
\ref{l03:thm:iterate}, verbatim in content.
\end{dictionary}

\begin{proof}[Proof on paper]
As in Chapter~\ref{ch:L03}: translation covariance from (VOA3) and Goddard; the $n$-th product theorem from Dong's lemma
and Goddard (the product field is local with everything and creates $u_{(k)}v$); the commutator formula from the
delta-function expansion of the commutator of local fields; associativity for negative indices from the normal-ordered
product formula of the field of $u_{(-k-1)}v$ (\chseventeenL{assoc_neg}).  No lower bound on weights is ever used.
\end{proof}

\begin{physics}
The OPE $u(z)v(w)\sim\sum_{k\ge0}(u_{(k)}v)(w)/(z-w)^{k+1}$ and its consequences: commutators of modes are
determined by the singular part of the OPE, the regular part by normal ordering.
\end{physics}

\begin{status}
Proved; printed in \leanpath{VOA/\allowbreak Moonshine/\allowbreak MoonshineCNAudit9.lean}; standard axioms.
\end{status}

\begin{theorem}[Borcherds identity and skew-symmetry for near-VOAs]\label{ch17:thm:nearborcherds}
For all $u,v,x$ and $\ell,m,n\in\Z$,
\begin{multline*}
\sum_{i\ge0}(-1)^i\binom\ell i\Bigl(u_{(\ell+m-i)}v_{(n+i)}x-(-1)^\ell v_{(\ell+n-i)}u_{(m+i)}x\Bigr)
=\sum_{i\ge0}\binom mi\,(u_{(\ell+i)}v)_{(m+n-i)}x,
\end{multline*}
and $v_{(t)}u=\sum_{j\ge0}\frac{(-1)^{t+j+1}}{j!}L_{-1}^j\bigl(u_{(t+j)}v\bigr)$.
\end{theorem}

\noindent{\raggedright Files \leanpath{VOA/\allowbreak Moonshine/\allowbreak MoonshineCNNearBorcherds.lean},
\leanpath{VOA/\allowbreak Moonshine/\allowbreak MoonshineCNNearSkew.lean}.\par}
\begin{leancode}
theorem jacobi (V : NearVOA 𝒱) (u v : 𝒱) (x : 𝒱) (ℓ m n : ℤ) :
    (∑ᶠ i : ℤ, (negOnePow i * cbz ℓ i) •
        (V.smul u (ℓ + m - i) (V.smul v (n + i) x)
          - ((1 : ℂ) * negOnePow ℓ) • V.smul v (ℓ + n - i) (V.smul u (m + i) x)))
      = ∑ᶠ i : ℤ, cbz m i • V.smul (V.smul u (ℓ + i) v) (m + n - i) x := by

theorem smul_eq_skew (V : NearVOA 𝒱) (a b : 𝒱) (t : ℤ) :
    V.smul b t a
      = ∑ᶠ j : ℕ, ((j.factorial : ℂ)⁻¹) •
          (sgn (-t - 1 - (j : ℤ))) •
            ((V.L (-1) ^ j) (V.smul a (t + (j : ℤ)) b)) := by
\end{leancode}

\begin{dictionary}
\chseventeenL{negOnePow i} $=(-1)^i$; \chseventeenL{cbz ℓ i} is $\binom\ell i$ for $i\ge0$ and $0$ for $i<0$ (so the sums over $\Z$ are sums over
$i\ge0$); \chseventeenL{sgn k} $=(-1)^k$.  The factor \chseventeenL{(1 : ℂ)} is the Koszul sign of a superalgebra, specialised to $1$
(the files are ports of the super versions at sign one).  The Jacobi identity is
the Borcherds identity of Chapter~\ref{ch:L03}, \eqref{l03:eq:borcherds}, in exactly the same index convention ($\ell$ the
product index, $m$ the first mode, $n$ the second); skew-symmetry is Theorem~\ref{l03:thm:skew}, written mode by mode.
\end{dictionary}

\begin{proof}[Proof on paper]
Both sides of the Jacobi identity satisfy the same Pascal recursion in $(\ell,m)$
(\chseventeenL{borcherdsProductSum_rec}, \chseventeenL{borcherdsIterateSum_rec}); they agree for $m=0$ (the iterate formula,
Theorem~\ref{ch17:thm:nearbasic}(iv)) and for $\ell=0$ (the commutator formula); the recursion then propagates agreement
to all $(\ell,m)$, every sum being finite on the fixed vector.  Skew-symmetry: apply Jacobi with $x=\chseventeenvac$ and use
$Y(w,z)\chseventeenvac=e^{zL_{-1}}w$.
\end{proof}

\begin{physics}
The Borcherds identity is the full content of the OPE: associativity and commutativity of operator products as
meromorphic functions.  Skew-symmetry $Y(u,z)v=e^{zL_{-1}}Y(v,-z)u$ is the exchange of the two insertion points.
\end{physics}

\begin{status}
Proved; printed in \leanpath{VOA/\allowbreak Moonshine/\allowbreak MoonshineCNAudit9.lean}; standard axioms.
\end{status}

\begin{theorem}[Consequences used by the physical-state Lie bracket]\label{ch17:thm:nearconseq}
For a near-VOA: (i) $[u_{(0)},v_{(n)}]=(u_{(0)}v)_{(n)}$; (ii) if $L_0u=u$ and $L_ju=0$ for $j>0$ (a weight-one primary),
then $[L_m,u_{(n)}]=-n\,u_{(m+n)}$ and $u_{(0)}$ preserves the spurious subspace; (iii) $u_{(0)}v+v_{(0)}u\in\operatorname{im}L_{-1}$
for all $u,v$.
\end{theorem}

\noindent File \leanpath{VOA/\allowbreak Moonshine/\allowbreak MoonshineCNNearConsequences.lean}.
\begin{leancode}
/-- Every zero mode is a derivation, with no homogeneity assumption. -/
theorem zeroMode_derivation (V : NearVOA 𝒱) (u v : 𝒱) (n : ℤ) :
    ⁅mode (V.Y u) 0, mode (V.Y v) n⁆ = mode (V.Y (V.smul u 0 v)) n := by

/-- Virasoro commutation for a primary state of conformal weight one. -/
theorem primary_comm (V : NearVOA 𝒱) (u : 𝒱) (hu : V.L 0 u = u)
    (hpos : ∀ j : ℤ, 0 < j → V.L j u = 0) (m n : ℤ) :
    ⁅V.L m, mode (V.Y u) n⁆ = (-n) • mode (V.Y u) (m + n) := by

/-- Primary weight-one zero modes preserve spurious states. -/
theorem zeroMode_preserves_spurious (V : NearVOA 𝒱) (u : 𝒱) (hu : V.L 0 u = u)
    (hpos : ∀ j : ℤ, 0 < j → V.L j u = 0) (x : 𝒱) (hx : x ∈ V.spurious) :
    V.smul u 0 x ∈ V.spurious := by

/-- The symmetric zero-product is a translation, for arbitrary states. -/
theorem skew_zero (V : NearVOA 𝒱) (u v : 𝒱) :
    V.smul u 0 v + V.smul v 0 u ∈ LinearMap.range (V.L (-1)) := by
\end{leancode}

\begin{dictionary}
\chseventeenL{hu}, \chseventeenL{hpos} say exactly ``$u$ is a weight-one Virasoro primary''.  These are the hypotheses of the
physical-state condition \eqref{ch17:eq:physical}; the subspace of such $u$ is \chseventeenL{primaryWeightOne}
(Chapter~\ref{ch:18}).  \textbf{Concrete instances.} \leanpath{VOA/\allowbreak Moonshine/\allowbreak MoonshineCNNearMatterIdentities.lean}
restates each theorem of this section for \chseventeenL{moonshineMatterNearVOA} and for \chseventeenL{lorentzianNearVOA}
(\chseventeenL{moonshineMatter_jacobi}, \chseventeenL{lorentzian_skew}, \dots), with no added hypothesis; e.g.
\end{dictionary}
\begin{leancode}
theorem moonshineMatter_skew_zero (u v : moonshineMatterSpace) :
    moonshineMatterNearVOA.smul u 0 v + moonshineMatterNearVOA.smul v 0 u ∈
      LinearMap.range (moonshineMatterNearVOA.L (-1)) :=
\end{leancode}
\noindent (file \leanpath{VOA/\allowbreak Moonshine/\allowbreak MoonshineCNNearMatterIdentities.lean}).

\begin{proof}[Proof on paper]
(i) is the commutator formula at $m=0$ (only $j=0$ survives).  (ii) The commutator formula with $\omega$ and $u$ and the
skew-symmetry $\omega_{(j)}u$ from $u_{(j)}\omega$; for a weight-one primary only $L_{-1}u$ and $L_0u=u$ contribute.  If
$x=L_{-k}y$ then $u_{(0)}x=L_{-k}u_{(0)}y$ since $[L_{-k},u_{(0)}]=0$.  (iii) Skew-symmetry at $t=0$:
$v_{(0)}u=-u_{(0)}v+L_{-1}(\cdots)$.
\end{proof}

\begin{physics}
Borcherds' bracket on physical states is $[u,v]=u_{(0)}v$: (i) is the Jacobi identity of this bracket, (ii) says the
bracket maps physical states to physical states and spurious to spurious (the gauge invariance of string amplitudes),
and (iii) says it is antisymmetric modulo $L_{-1}$-exact states, which are null.
\end{physics}

\begin{status}
Proved; printed in \leanpath{VOA/\allowbreak Moonshine/\allowbreak MoonshineCNAudit9.lean}; standard axioms.
\end{status}

%======================================================================
\section{Contragredient (FHL) adjointness without a lower bound}\label{ch17:sec:fhl}
%======================================================================

The Monster Lie algebra carries an invariant bilinear form, induced from the matter form $B$ of
Theorem~\ref{ch17:thm:matterform}.  Its invariance $B([u,v],w)=B(u,[v,w])$ requires that $B$ be \emph{FHL-invariant}:
for $u$ of weight $N$,
\begin{equation}\label{ch17:eq:fhl}
  B(u_{(n)}x,y)=(-1)^N\sum_{j\ge0}\frac1{j!}\,B\bigl(x,(L_1^ju)_{(2N-n-2-j)}y\bigr)
\end{equation}
(Chapter~\ref{ch:L04}, Definition~\ref{l04:def:fullinv}).  For a VOA of CFT type the sum stops at $j=N$ because
$L_1^ju$ has weight $N-j<0$ otherwise.  In a near-VOA this argument fails; the cutoff must be chosen differently.

\begin{warning}
\textbf{The total-weight cutoff is wrong (a kernel-checked counterexample).}  In $\chseventeenVII$ the state
$w=L_{-1}e^{(1,-1)}$ has total weight $0$ (the charged vacuum $e^{(1,-1)}$ has weight $\Delta(1,-1)=-1$), yet
$L_1w=-2e^{(1,-1)}\neq0$.  So a formula \eqref{ch17:eq:fhl} with the CFT-type cutoff $j\le N=0$ would drop a nonzero term.
This was first noticed on paper and is now a Lean theorem
(Theorem~\ref{ch17:thm:cutoffwitness}).  The correct cutoff is the \emph{oscillator level above the sector's ground
energy}, $N-\Delta(\alpha)$ (Theorem~\ref{ch17:thm:nilpotence}).
\end{warning}

\begin{theorem}[Cutoff witness]\label{ch17:thm:cutoffwitness}
In any near-VOA, if $u\in\mathcal V_N$ and $L_1u=0$ then $L_1L_{-1}u=2N u$.  In $\chseventeenVII$, $w=L_{-1}e^{(1,-1)}$ lies in weight $0$
and $L_1w=-2e^{(1,-1)}$; hence some weight-zero state is not killed by $L_1$.
\end{theorem}

\noindent File \leanpath{VOA/\allowbreak Moonshine/\allowbreak MoonshineCNChargeCutoffWitness.lean}.
\begin{leancode}
noncomputable def lorentzianNegativeChargeDescendant : LorentzianFockSpace :=
  lorentzianNearVOA.L (-1) (lorentzianChargeState ![1, -1])

theorem lorentzianNegativeChargeDescendant_weight :
    lorentzianNegativeChargeDescendant ∈ lorentzianNearVOA.grading 0 := by

/-- At total weight zero the first omitted FHL term is explicitly nonzero. -/
theorem lorentzianNegativeChargeDescendant_virasoroOne :
    lorentzianNearVOA.L 1 lorentzianNegativeChargeDescendant =
      (-2 : ℂ) • lorentzianChargeState ![1, -1] := by

theorem lorentzianVirasoroOne_weightZero_not_zero :
    ∃ u ∈ lorentzianNearVOA.grading 0, lorentzianNearVOA.L 1 u ≠ 0 := by
\end{leancode}

\begin{dictionary}
The general first sentence is the lemma \chseventeenL{virasoroOne_translation_primary} (same file, not printed here; it is
the computation in the proof below).  The charged vacuum is primary
(\chseventeenL{frameChargeState_virasoroOne}, \leanpath{VOA/\allowbreak Moonshine/\allowbreak MoonshineCNGeneratorPrimary.lean}).  Nonvanishing of
$e^{(1,-1)}$ is \chseventeenL{lorentzianChargeState_ne_zero}.
\end{dictionary}

\begin{proof}[Proof on paper]
$[L_1,L_{-1}]=2L_0$, so $L_1L_{-1}u=L_{-1}L_1u+2L_0u=2Nu$; with $N=\Delta(1,-1)=-1$ this is $-2e^{(1,-1)}$.
\end{proof}

\begin{physics}
In a theory with negative-energy ground states, quasi-primary descendants can have weight zero without being
annihilated by $L_1$; the BPZ conjugation $e^{zL_1}(-z^{-2})^{L_0}$ must then keep more terms than the weight suggests.
\end{physics}

\begin{status}
Proved; printed in \leanpath{VOA/\allowbreak Moonshine/\allowbreak MoonshineCNAudit13.lean}; standard axioms.
\end{status}

\subsection{Local nilpotence of \texorpdfstring{$L_1$}{L1}}

\begin{theorem}[Charge-sensitive nilpotence]\label{ch17:thm:nilpotence}
(i) In any near-VOA, let $S$ be an $L_1$-stable subspace with $S\cap\mathcal V_N=0$ for $N<E_0$ (Lean \chseventeenL{E}).  Then $L_1^ku=0$ for
$u\in S\cap\mathcal V_N$ and $k>N-E_0$.  If $\mathcal V$ is spanned by such pieces $S_\alpha\cap\mathcal V_N$, then $L_1$ is locally nilpotent.
(ii) In $V_L$ (any even $L$, complex frame), the momentum-$\alpha$ part of $(V_L)_N$ is $(\chseventeenF_r)_{N-\Delta(\alpha)}\otimes e^\alpha$; hence
$L_1^ku=0$ for $u$ of weight $N$ and momentum $\alpha$ and $k>N-\Delta(\alpha)$.
(iii) The same holds in $\chseventeenM$; in particular $L_1^{(N-\Delta(\alpha))+1}u=0$ and $L_1$ is locally nilpotent on $\chseventeenVII$ and on
$\chseventeenM$.
\end{theorem}

\noindent{\raggedright Files \leanpath{VOA/\allowbreak Moonshine/\allowbreak MoonshineCNNearNilpotence.lean},
\leanpath{VOA/\allowbreak Moonshine/\allowbreak MoonshineCNChargeNilpotence.lean},
\leanpath{VOA/\allowbreak Moonshine/\allowbreak MoonshineCNMatterNilpotence.lean}.\par}
\begin{leancode}
/-- The cutoff counts level above the sector ground energy, including negative energies. -/
theorem virasoroOne_pow_zero_of_lower_bound (S : Submodule ℂ H) (E : ℤ)
    (hS : ∀ u ∈ S, V.L 1 u ∈ S)
    (hE : ∀ N : ℤ, N < E → S ⊓ V.grading N = ⊥)
    {N : ℤ} {u : H} (hu : u ∈ V.grading N) (hSu : u ∈ S)
    (k : ℕ) (hk : (N - E).toNat < k) : (V.L 1 ^ k) u = 0 := by

/-- Fixing total weight and charge recovers exactly one oscillator weight. -/
theorem latticeMomentum_weight_intersection (L : IntegralLattice) (α : L.carrier) (N : ℤ) :
    latticeMomentumGrading L α ⊓ gradingLatGen L N =
      bigradα L (N - halfNorm L α, α) := by

/-- Local nilpotence on arbitrary finite sums of charges and oscillator weights. -/
theorem lorentzianVirasoroOne_locallyNilpotent (u : LorentzianFockSpace) :
    ∃ k : ℕ, (lorentzianNearVOA.L 1 ^ k) u = 0 := by

/-- The requested explicit finite cutoff on the actual reconstructed matter algebra. -/
theorem moonshineMatterVirasoroOne_cutoff (α : lorentzianLattice.carrier)
    {N : ℤ} {u : moonshineMatterSpace} (hu : u ∈ moonshineMatterNearVOA.grading N)
    (hα : u ∈ moonshineMatterMomentum α) :
    (moonshineMatterNearVOA.L 1 ^ ((N - LatGen.halfNorm lorentzianLattice α).toNat + 1)) u = 0 :=

theorem moonshineMatterVirasoroOne_locallyNilpotent (u : moonshineMatterSpace) :
    ∃ k : ℕ, (moonshineMatterNearVOA.L 1 ^ k) u = 0 :=
\end{leancode}

\begin{dictionary}
\chseventeenL{(N - E).toNat} is $\max(N-E_0,0)$.  \chseventeenL{S ⊓ V.grading N} is the intersection.  \chseventeenL{halfNorm L α} $=\Delta(\alpha)$
(Theorem~\ref{ch17:thm:grading}(ii)).  ``Locally nilpotent'' means: for each $u$ some power of $L_1$ kills $u$ (the power
may depend on $u$); this is the precise hypothesis \chseventeenL{hNil} of the next theorems.  For $\chseventeenM$ the lower bound
uses $\Vnat_k=0$ for $k<0$ (CFT type of $\Vnat$): a state of momentum $\alpha$ has weight $\ge \Delta(\alpha)$.
\end{dictionary}

\begin{proof}[Proof on paper]
$L_1$ lowers weight by one and preserves momentum (it is a zero-momentum operator).  Within momentum $\alpha$ the weights
are $\ge \Delta(\alpha)$ because oscillator levels (and $\Vnat$ weights) are $\ge0$.  So after more than $N-\Delta(\alpha)$ steps the weight
drops below the sector's floor and the vector is $0$.  For a general vector, take the maximum of the finitely many
cutoffs of its homogeneous components.
\end{proof}

\begin{physics}
Each momentum sector of the Lorentzian theory \emph{is} bounded below (by its ground energy $p^+p^-$); only the union is
not.  The BPZ conjugate of a state therefore involves finitely many $L_1$-descendants, counted from the sector's ground
state rather than from zero.
\end{physics}

\begin{status}
Proved; printed in \leanpath{VOA/\allowbreak Moonshine/\allowbreak MoonshineCNAudit13.lean}; standard axioms.
\end{status}

\subsection{Opposite modes}

\begin{definition}[Opposite modes]\label{ch17:def:opposite}
Let $V$ be a near-VOA in which $L_1$ is locally nilpotent.  For $u\in\mathcal V_N$ define
\begin{equation*}
  u^{\mathrm{op}}_{(n)}=(-1)^N\sum_{j\ge0}\frac1{j!}\,(L_1^ju)_{(2N-n-2-j)}
\end{equation*}
(a finite sum by nilpotence), and extend linearly over the weight decomposition.
\end{definition}

\noindent{\raggedright Files \leanpath{VOA/\allowbreak Moonshine/\allowbreak MoonshineCNNearOppositeModes.lean},
\leanpath{VOA/\allowbreak Moonshine/\allowbreak MoonshineCNMatterOpposite.lean}.\par}
\begin{leancode}
def oppositeModeTerm (N n : ℤ) (j : ℕ) : H →ₗ[ℂ] End ℂ H where
  toFun u := ((j.factorial : ℂ)⁻¹) • mode (V.Y ((V.L 1 ^ j) u))
    (2 * N - n - 2 - (j : ℤ))

/-- The FHL opposite mode at a specified homogeneous weight. -/
def oppositeModeAtWeight (N n : ℤ) (u : H) : End ℂ H :=
  (-1 : ℂ) ^ N • ∑ᶠ j : ℕ, V.oppositeModeTerm N n j u

/-- Assemble homogeneous opposite modes using the given internal weight decomposition. -/
def oppositeModes (hNil : ∀ u : H, ∃ k : ℕ, (V.L 1 ^ k) u = 0)
    (n : ℤ) : H →ₗ[ℂ] End ℂ H :=

/-- Any nilpotence cutoff computes the same opposite mode. -/
theorem oppositeModeAtWeight_eq_sum (N n : ℤ) {u : H} (k : ℕ)
    (hk : (V.L 1 ^ k) u = 0) :
    V.oppositeModeAtWeight N n u = (-1 : ℂ) ^ N •
      ∑ j ∈ Finset.range k, ((j.factorial : ℂ)⁻¹) •
        mode (V.Y ((V.L 1 ^ j) u)) (2 * N - n - 2 - (j : ℤ)) := by

/-- The opposite zero mode of a weight-one primary is minus its ordinary zero mode. -/
theorem oppositeModes_primary (hNil : ∀ u : H, ∃ k : ℕ, (V.L 1 ^ k) u = 0)
    {u : H} (hu : u ∈ V.grading 1) (hL : V.L 1 u = 0) :
    V.oppositeModes hNil 0 u = -mode (V.Y u) 0 := by

def moonshineMatterOppositeModes (n : ℤ) :
    moonshineMatterSpace →ₗ[ℂ] Module.End ℂ moonshineMatterSpace :=
  moonshineMatterNearVOA.oppositeModes moonshineMatterVirasoroOne_locallyNilpotent n
\end{leancode}

\begin{dictionary}
\chseventeenL{oppositeModeAtWeight N n u} is the displayed formula with $N$ a parameter (it is meaningful for $u\in\mathcal V_N$);
\chseventeenL{∑ᶠ} is finite because $L_1^ju=0$ for large $j$.  \chseventeenL{oppositeModes} (body, using
\chseventeenL{V.isInternal.chooseDecomposition} and \chseventeenL{DirectSum.toModule}) decomposes $u$ into homogeneous components and
applies \chseventeenL{oppositeModeAtWeight} at the right weight on each; \chseventeenL{oppositeModes_homogeneous} confirms that on
$u\in\mathcal V_N$ it equals \chseventeenL{oppositeModeAtWeight N n u}.  The choice of decomposition instance is immaterial (the
decomposition is unique).  \chseventeenL{oppositeModeAtWeight_eq_sum} says \emph{any} killing exponent gives the same value, so the
correct cutoff need not be known; for $\chseventeenM$, \chseventeenL{moonshineMatterOppositeModes_charge_cutoff} uses
$(N-\Delta(\alpha))+1$.  Comparing with Chapter~\ref{ch:L04}, \eqref{l04:eq:fhl}: same coefficients, same mode index
$2N-n-2-j$, different cutoff.
\end{dictionary}

\begin{physics}
$u^{\mathrm{op}}_{(n)}$ are the modes of the BPZ-conjugate field $Y(e^{zL_1}(-z^{-2})^{L_0}u,z^{-1})$.  For a weight-one primary
current $J$ this is $J^{\mathrm{op}}_{(0)}=-J_{(0)}$: charges are anti-self-adjoint for the bilinear pairing.
\end{physics}

\begin{status}
Definitions and proved properties; printed in \leanpath{VOA/\allowbreak Moonshine/\allowbreak MoonshineCNAudit13.lean}; standard axioms.
\end{status}

\subsection{From generators to all fields}

\begin{definition}[Opposite iterate identity]\label{ch17:def:oppiterate}
A near-VOA with locally nilpotent $L_1$ satisfies the \emph{opposite iterate identity} if for all $u,v,y$ and $\ell,n\in\Z$
\begin{equation*}
  (u_{(\ell)}v)^{\mathrm{op}}_{(n)}y=\sum_{i\ge0}(-1)^i\binom\ell i\Bigl(v^{\mathrm{op}}_{(n+i)}u^{\mathrm{op}}_{(\ell-i)}y-(-1)^\ell u^{\mathrm{op}}_{(i)}v^{\mathrm{op}}_{(\ell+n-i)}y\Bigr).
\end{equation*}
The \emph{intertwined states} for a form $B$ are those $u$ with $B(u_{(n)}x,y)=B(x,u^{\mathrm{op}}_{(n)}y)$ for all $n,x,y$.
\end{definition}

\noindent File \leanpath{VOA/\allowbreak Moonshine/\allowbreak MoonshineCNNearAdjointClosure.lean}.
\begin{leancode}
/-- States whose actual modes are adjoint to the constructed opposite modes. -/
def oppositeIntertwinedStates (hNil : ∀ u : H, ∃ k : ℕ, (V.L 1 ^ k) u = 0)
    (B : LinearMap.BilinForm ℂ H) : Submodule ℂ H where
  carrier := {u | ∀ (n : ℤ) (x y : H),
    B (mode (V.Y u) n x) y = B x (V.oppositeModes hNil n u y)}

/-- The precise remaining opposite-field identity needed for adjoint closure. -/
def OppositeIterateIdentity (hNil : ∀ u : H, ∃ k : ℕ, (V.L 1 ^ k) u = 0) : Prop :=
  ∀ (u v y : H) (ℓ n : ℤ),
    V.oppositeModes hNil n (V.smul u ℓ v) y =
      ∑ᶠ i : ℕ, ((-1 : ℂ) ^ i * gBinom ℓ i) •
        (V.oppositeModes hNil (n + i) v (V.oppositeModes hNil (ℓ - i) u y) -
          (-1 : ℂ) ^ ℓ • V.oppositeModes hNil i u
            (V.oppositeModes hNil (ℓ + n - i) v y))
\end{leancode}

\begin{dictionary}
The identity is the iterate formula (Theorem~\ref{ch17:thm:nearbasic}(iv)) with every mode replaced by its opposite and the
order of the two factors reversed, which is what adjunction does to a product.  The docstring's word ``remaining'' is
historical: the identity is proved for every near-VOA with locally nilpotent $L_1$ in
Theorem~\ref{ch17:thm:oppborcherds}.
\end{dictionary}

\begin{physics}
Adjunction reverses the order of operators; the opposite iterate identity says that BPZ conjugation is compatible with
the OPE, i.e.\ that the conjugate fields again satisfy associativity.
\end{physics}

\begin{status}
Definitions; printed in \leanpath{VOA/\allowbreak Moonshine/\allowbreak MoonshineCNAudit13.lean}.
\end{status}

\begin{theorem}[Opposite Borcherds identity; adjoint closure]\label{ch17:thm:oppborcherds}
Let $V$ be a near-VOA with $L_1$ locally nilpotent.  (i) For $u\in\mathcal V_p$, any $v$, any $q,\ell,m,n$:
\begin{multline*}
\sum_i(-1)^i\binom\ell i\Bigl(v^{\mathrm{op},q}_{(n+i)}u^{\mathrm{op},p}_{(\ell+m-i)}y-(-1)^\ell u^{\mathrm{op},p}_{(m+i)}v^{\mathrm{op},q}_{(\ell+n-i)}y\Bigr)
=\sum_i\binom mi\,(u_{(\ell+i)}v)^{\mathrm{op},\,p+q-\ell-i-1}_{(m+n-i)}y,
\end{multline*}
where $w^{\mathrm{op},N}$ denotes \chseventeenL{oppositeModeAtWeight N}.  (ii) $V$ satisfies the opposite iterate identity.
(iii) Let $B$ be any bilinear form and $G$ a set of states generating $V$ under all products $u_{(\ell)}v$.  If every $u\in G$
satisfies $B(u_{(n)}x,y)=B(x,u^{\mathrm{op}}_{(n)}y)$, then every $u\in\mathcal V$ does.
\end{theorem}

\noindent{\raggedright Files \leanpath{VOA/\allowbreak Moonshine/\allowbreak MoonshineCNNearOppositeIterate.lean},
\leanpath{VOA/\allowbreak Moonshine/\allowbreak MoonshineCNNearAdjointClosure.lean}.\par}
\begin{leancode}
theorem oppositeModeAtWeight_borcherds (hNil : ∀ w : H, ∃ k : ℕ, (V.L 1 ^ k) w = 0)
    {p q : ℤ} {u : H} (hu : u ∈ V.grading p) (v : H) (ℓ m n : ℤ) (y : H) :
    (∑ᶠ i : ℤ, (negOnePow i * cbz ℓ i) •
        (V.oppositeModeAtWeight q (n + i) v (V.oppositeModeAtWeight p (ℓ + m - i) u y)
          - negOnePow ℓ • V.oppositeModeAtWeight p (m + i) u
              (V.oppositeModeAtWeight q (ℓ + n - i) v y)))
      = ∑ᶠ i : ℤ, cbz m i •
          V.oppositeModeAtWeight (p + q - (ℓ + i) - 1) (m + n - i) (V.smul u (ℓ + i) v) y := by

/-- **The opposite iterate identity holds** for every near-VOA with locally nilpotent `L₁`. -/
theorem oppositeIterateIdentity_of_locallyNilpotent
    (hNil : ∀ u : H, ∃ k : ℕ, (V.L 1 ^ k) u = 0) : V.OppositeIterateIdentity hNil := by

/-- Adjointness of all fields from adjointness of generators, for every near-VOA with
locally nilpotent `L₁`. The opposite iterate law is now proved, not assumed. -/
theorem adjoint_of_locallyNilpotent_generators
    (hNil : ∀ u : H, ∃ k : ℕ, (V.L 1 ^ k) u = 0) (B : LinearMap.BilinForm ℂ H) (G : Set H)
    (hGen : ∀ S : Submodule ℂ H, G ⊆ S →
      (∀ u ∈ S, ∀ v ∈ S, ∀ ℓ : ℤ, V.smul u ℓ v ∈ S) → S = ⊤)
    (hAdj : ∀ u ∈ G, ∀ (n : ℤ) (x y : H),
      B (mode (V.Y u) n x) y = B x (V.oppositeModes hNil n u y))
    (u : H) (n : ℤ) (x y : H) :
    B (mode (V.Y u) n x) y = B x (V.oppositeModes hNil n u y) :=
\end{leancode}

\begin{dictionary}
\textbf{Hypotheses of (iii)}: \chseventeenL{hNil} (local nilpotence; discharged for $\chseventeenM$ by
\chseventeenL{moonshineMatterVirasoroOne_locallyNilpotent}); \chseventeenL{hGen} says $G$ generates: every subspace containing $G$ and
closed under all products is everything; \chseventeenL{hAdj} is adjointness on generators.  The conclusion is FHL invariance
\eqref{ch17:eq:fhl} for every state, in opposite-mode form.  The earlier conditional version
\chseventeenL{adjoint_of_oppositeIterate_generators} (\leanpath{VOA/\allowbreak Moonshine/\allowbreak MoonshineCNNearAdjointClosure.lean}) has the extra
hypothesis \chseventeenL{hIter : V.OppositeIterateIdentity hNil}, which (ii) discharges.  \textbf{Fidelity.} No lower bound, no
finite-dimensionality, no symmetry or nondegeneracy of $B$ is assumed; this is stronger than the VOA version of
Chapter~\ref{ch:L04} in scope.
\end{dictionary}

\begin{proof}[Proof on paper]
Both sides of (i) are reduced to one block sum
\begin{equation*}
(-1)^{p+q}\sum_{b<B_0}\sum_{a<A_0}\frac1{a!\,b!}\sum_i\binom{2q-\ell-n-2-b}{i}\bigl((L_1^bv)_{(\ell+i)}(L_1^au)\bigr)_{(\cdots)}y,
\end{equation*}
where $A_0,B_0$ are \emph{arbitrary} exponents with $L_1^{A_0}u=0$, $L_1^{B_0}v=0$.  The opposite-pair side: expand both opposite
modes with these cutoffs and apply the Jacobi identity (Theorem~\ref{ch17:thm:nearborcherds}) to each pair
$(L_1^bv,L_1^au)$.  The product side: the opposite mode of $u_{(s)}v$ involves $L_1^t(u_{(s)}v)$, which is expanded by the
``comultiplication'' formula $L_1(u_{(s)}x)=(2p-s-2)u_{(s+1)}x+(L_1u)_{(s)}x+u_{(s)}L_1x$ (\chseventeenL{virasoroOne_smul_of_mem_weight},
from the commutator formula with $\omega$ and $L_0u=pu$), iterated with explicit coefficients (\chseventeenL{virasoroOne_pow_smul_eq_sum});
the triangle of indices $a+b\le t$ is moved into the box $a<A_0,b<B_0$ (the summand vanishes outside), and both sides
collapse by Vandermonde's convolution to the same coefficients.  (ii) is (i) at $m=0$, extended from homogeneous $u$ to
all $u$ by additivity.  (iii): the intertwined states form a subspace containing $G$; by the ordinary iterate formula, (ii)
and adjunction, it is closed under all products; so it is everything.
\end{proof}

\begin{physics}
The statement ``if the generating fields are BPZ-self-conjugate, all fields are'' --- the standard way to establish an
invariant form on a CFT (check it on the currents and vertex operators).  The point of the theorem is that it survives
the absence of a bounded-below spectrum.
\end{physics}

\begin{status}
(i), (ii), (iii) proved with the displayed hypotheses only; printed in \leanpath{VOA/\allowbreak Moonshine/\allowbreak MoonshineCNAudit16.lean}
(43 \texttt{\#print axioms} lines, each reported as exactly \chseventeenL{[propext, Classical.choice, Quot.sound]} in the development record).  The
supporting files \texttt{NearOppositeComul}, \texttt{NearOppositeLeft}, \texttt{NearOppositeRight} contain the two
block-sum reductions.
\end{status}

\subsection{The generators of \texorpdfstring{$\chseventeenM$}{M} are adjoint}

\begin{theorem}[Adjointness of the generating fields]\label{ch17:thm:genadjoint}
(i) On $\chseventeenF_r$ with the Fock form $B_{\mathcal F}$: $B_{\mathcal F}(L_np,q)=B_{\mathcal F}(p,L_{-n}q)$ and $B_{\mathcal F}(\alpha(n)p,q)=-B_{\mathcal F}(p,\alpha(-n)q)$ (frame
oscillators are skew).  (ii) The dressing satisfies $B_{\mathcal F}(\mathcal E_t(\gamma)y,y')=B_{\mathcal F}(y,\mathcal E_{-t}(\gamma)y')$.  (iii) In $V_L$
($L$ even), with $B_L$ and $N=\Delta(\gamma)$: $B_L(e^\gamma_{(n)}x,y)=(-1)^NB_L(x,e^\gamma_{(2N-n-2)}y)$ and
$B_L(h_a(n)x,y)=-B_L(x,h_a(-n)y)$ for the reconstructed fields; and the charged vacua and the currents generate $V_L$.
(iv) In $\chseventeenM$ the same holds for $\chseventeenvac\otimes e^\gamma$, $\chseventeenvac\otimes h_a(-1)\chseventeenvac$, and \eqref{ch17:eq:fhl} holds for every
$a\otimes\chseventeenvac$, $a\in\Vnat_N$.
\end{theorem}

\noindent{\raggedright Files \leanpath{VOA/\allowbreak Moonshine/\allowbreak MoonshineCNFockAdjoint.lean},
\leanpath{VOA/\allowbreak Moonshine/\allowbreak MoonshineCNDressingAdjoint.lean},
\leanpath{VOA/\allowbreak Moonshine/\allowbreak MoonshineCNChargedAdjoint.lean},
\leanpath{VOA/\allowbreak Moonshine/\allowbreak MoonshineCNCurrentAdjoint.lean},
\leanpath{VOA/\allowbreak Moonshine/\allowbreak MoonshineCNTensorGeneratorAdjoint.lean}.\par}
\begin{leancode}
/-- The Fock bilinear form makes `Lₙ` adjoint to `L₋ₙ` in every rank. -/
theorem uForm_virasoro_adjoint (r : ℕ) (n : ℤ) (p q : FockN r) :
    uForm r (LvirN r 1 n p) q = uForm r p (LvirN r 1 (-n) q) := by

theorem uForm_frameOscillator_adjoint {L : IntegralLattice} (F : ComplexLatticeFrame L)
    (α : L.carrier) (n : ℤ) (p q : FockN (rank L)) :
    uForm (rank L) (frameOscillator F α n p) q =
      -uForm (rank L) p (frameOscillator F α (-n) q) := by

theorem uForm_dressingField_adjoint (γ : L.carrier) (t : ℤ) (y y' : FockN (rank L)) :
    uForm (rank L) (dressingField F γ t y) y' = uForm (rank L) y (dressingField F γ (-t) y') := by

/-- Charged-state adjointness for the reconstructed state-field map itself. -/
theorem latForm_frameStateField_charge_adjoint (hEven : L.Even) (γ : L.carrier)
    (n : ℤ) (x y : LatModuleGen L) :
    latForm L (mode (frameStateField F hEven (latChargeGen L γ)) n x) y =
      (-1 : ℂ) ^ (LatGen.halfNorm L γ) *
        latForm L x (mode (frameStateField F hEven (latChargeGen L γ))
          (2 * LatGen.halfNorm L γ - n - 2) y) := by

theorem latForm_frameStateField_current_adjoint (hEven : L.Even) (a : Fin (rank L))
    (n : ℤ) (x y : LatModuleGen L) :
    latForm L (mode (frameStateField F hEven (frameWordState L [(0, a)] 0)) n x) y =
      -latForm L x
        (mode (frameStateField F hEven (frameWordState L [(0, a)] 0)) (-n) y) := by

/-- The tested generators generate the full reconstructed lattice vertex algebra. -/
theorem frameStateField_generated_by_charges_currents (hEven : L.Even)
    (S : Submodule ℂ (LatModuleGen L))
    (hCharge : ∀ α, latChargeGen L α ∈ S)
    (hCurrent : ∀ a, frameWordState L [(0, a)] 0 ∈ S)
    (hClosed : ∀ u ∈ S, ∀ v ∈ S, ∀ n : ℤ,
      mode (frameStateField F hEven u) n v ∈ S) : S = ⊤ := by

/-- Full FHL identity on the Vnat factor embedded as tensor generators. -/
theorem moonshineMatterForm_vnatGenerator_fhl (wt : ℤ) (a : Vnat)
    (ha : a ∈ VnatVOA.grading wt) (n : ℤ) (x y : moonshineMatterSpace) :
    moonshineMatterForm
      (mode (moonshineMatterNearVOA.Y (a ⊗ₜ[ℂ] lorentzianNearVOA.vacuum)) n x) y =
      (-1 : ℂ) ^ wt * ∑ j ∈ Finset.range (wt.toNat + 1), ((j.factorial : ℂ)⁻¹) *
        moonshineMatterForm x
          (mode (moonshineMatterNearVOA.Y
            (((VnatVOA.L 1 ^ j) a) ⊗ₜ[ℂ] lorentzianNearVOA.vacuum))
            (2 * wt - n - 2 - (j : ℤ)) y) := by
\end{leancode}

\begin{dictionary}
\chseventeenL{frameWordState L [(0, a)] 0} $=h_a(-1)\chseventeenvac$.  The charged formula is \eqref{ch17:eq:fhl} for the primary $e^\gamma$ (only
$j=0$), with \chseventeenL{halfNorm L γ} $=\Delta(\gamma)$ the weight; the current formula is \eqref{ch17:eq:fhl} for the weight-one primary
$h_a(-1)\chseventeenvac$.  For $\Vnat$ the cutoff \chseventeenL{wt.toNat + 1} is the CFT-type one, correct because $\Vnat$ is bounded below;
the proof transports the FHL invariance of \chseventeenL{natForm} (\chseventeenL{SimpleNat}, Chapter~\ref{ch:14}).  The analogous charged and
current statements on $\chseventeenM$ (\chseventeenL{moonshineMatterForm_chargeGenerator_adjoint},
\chseventeenL{moonshineMatterForm_currentGenerator_adjoint}) are in the same file.  \chseventeenL{frameStateField_generated_by_charges_currents}
is the generation hypothesis \chseventeenL{hGen} for $V_L$, in the form ``closed submodule containing the generators is everything''.
\end{dictionary}

\begin{proof}[Proof on paper]
(i) Fock form: $F$ pairs monomials diagonally with weights making $x_{j,a}\cdot$ and $-(j+1)\partial/\partial x_{j,a}$ adjoint; the
Sugawara modes are quadratic in oscillators, so the two signs cancel.  (ii) The adjoint of $E^-(\gamma,z)E^+(\gamma,z)$ is
$E^-(\gamma,z^{-1})E^+(\gamma,z^{-1})$ after using skew-adjointness of the oscillators; proved by induction on polynomials from the
commutator of oscillators with the dressing (\chseventeenL{directionOscillator_dressingField_apply}).  (iii) Combine (ii) with the
pairing of $e^\beta$ against $e^{-\beta}$, the cocycle identity for $\chi(\alpha+\beta)$ and the $z^{\langle\gamma,\beta\rangle}$ shift; the sign
$(-1)^{\Delta(\gamma)}$ comes out of the cocycle.  The reconstructed field of a word state is identified with the explicit current/charged
field, and generation follows by induction on words (a word state is a product of a current mode with a shorter word state).
(iv) Tensor with $\Vnat$: modes of $\chseventeenvac\otimes b$ act as $1\otimes b_{(n)}$, modes of $a\otimes\chseventeenvac$ as $a_{(n)}\otimes1$,
and $L_1$ acts factorwise (\chseventeenL{moonshineMatterVnatVirasoroOne_pow}).
\end{proof}

\begin{physics}
The light-cone currents are anti-self-conjugate, the vertex operators $e^\gamma$ are conjugate to themselves up to the sign
$(-1)^{h}$ of a weight-$h$ primary, and the Monster CFT's fields are BPZ-self-conjugate.
\end{physics}

\begin{status}
Proved, no hypotheses beyond evenness and a frame (discharged for $\chseventeenII$).  Printed in
\leanpath{VOA/\allowbreak Moonshine/\allowbreak MoonshineCNAudit3.lean} (Fock adjoints) and \leanpath{VOA/\allowbreak Moonshine/\allowbreak MoonshineCNAudit12.lean}
(dressing, charged, current, tensor generators); standard axioms.
\end{status}

\begin{theorem}[FHL invariance of the matter form, reduced to its generators]\label{ch17:thm:matterfhl}
(i) The states $a\otimes\chseventeenvac$ ($a\in\Vnat$), $\chseventeenvac\otimes e^\alpha$ and $\chseventeenvac\otimes h_a(-1)\chseventeenvac$ generate $\chseventeenM$.
(ii) Each of them is adjoint to its opposite modes for $B=B_\natural\otimes B_L$.
(iii) Conditional on the opposite iterate identity for $\chseventeenM$, every state of $\chseventeenM$ is; and then the descended form
on the Monster Lie algebra is invariant.  (The condition is discharged by Theorem~\ref{ch17:thm:oppborcherds}(ii); the
unconditional statements \chseventeenL{moonshineMatterOppositeModes_adjoint} and \chseventeenL{monsterLieBilinearForm_invariant} are in the
Chapter~\ref{ch:18} file \leanpath{VOA/\allowbreak Moonshine/\allowbreak MoonshineCNMonsterLieInvariance.lean}.)
\end{theorem}

\noindent{\raggedright Files \leanpath{VOA/\allowbreak Moonshine/\allowbreak MoonshineCNMatterOpposite.lean},
\leanpath{VOA/\allowbreak Moonshine/\allowbreak MoonshineCNMatterAdjointReduction.lean}.\par}
\begin{leancode}
/-- Charges, currents and the embedded Vnat factor generate the actual tensor algebra. -/
theorem moonshineMatter_generated_by_tensor_generators (S : Submodule ℂ moonshineMatterSpace)
    (hVnat : ∀ a : Vnat, a ⊗ₜ[ℂ] lorentzianNearVOA.vacuum ∈ S)
    (hCharge : ∀ α, VnatVOA.vacuum ⊗ₜ[ℂ] latChargeGen lorentzianLattice α ∈ S)
    (hCurrent : ∀ a, VnatVOA.vacuum ⊗ₜ[ℂ]
      frameWordState lorentzianLattice [(0, a)] 0 ∈ S)
    (hClosed : ∀ u ∈ S, ∀ v ∈ S, ∀ n : ℤ,
      mode (moonshineMatterNearVOA.Y u) n v ∈ S) : S = ⊤ := by

/-- All actual matter fields are adjoint, conditional only on their opposite iterate law. -/
theorem moonshineMatterOppositeModes_adjoint_of_iterate
    (hIter : moonshineMatterNearVOA.OppositeIterateIdentity
      moonshineMatterVirasoroOne_locallyNilpotent)
    (u : moonshineMatterSpace) (n : ℤ) (x y : moonshineMatterSpace) :
    moonshineMatterForm (mode (moonshineMatterNearVOA.Y u) n x) y =
      moonshineMatterForm x (moonshineMatterOppositeModes n u y) := by

/-- The exact remaining hypothesis for the Monster Lie form, with all generators discharged. -/
theorem monsterLieBilinearForm_invariant_of_oppositeIterate
    (hIter : moonshineMatterNearVOA.OppositeIterateIdentity
      moonshineMatterVirasoroOne_locallyNilpotent)
    (u v w : monsterLieAlgebra) :
    monsterLieBilinearForm ⁅u, v⁆ w = monsterLieBilinearForm u ⁅v, w⁆ :=
\end{leancode}

\begin{dictionary}
\chseventeenL{monsterLieAlgebra} (the weight-one primaries of $\chseventeenM$ modulo spurious states) and \chseventeenL{monsterLieBilinearForm}
(the descended form) are defined in Chapter~\ref{ch:18}.  The single hypothesis \chseventeenL{hIter} is a closed proposition with no
free variables; it is \emph{proved} as \chseventeenL{moonshineMatterNearVOA_oppositeIterateIdentity} (Chapter~\ref{ch:18} file, a
one-line application of \chseventeenL{oppositeIterateIdentity_of_locallyNilpotent}).  The other two generator-adjointness theorems
of \leanpath{VOA/\allowbreak Moonshine/\allowbreak MoonshineCNMatterAdjointReduction.lean} (\chseventeenL{moonshineMatterOppositeModes_charge_adjoint},
\chseventeenL{…_current_adjoint}, \chseventeenL{…_vnat_adjoint}) restate Theorem~\ref{ch17:thm:genadjoint}(iv) in opposite-mode form, using the
primarity computations of \leanpath{VOA/\allowbreak Moonshine/\allowbreak MoonshineCNGeneratorPrimary.lean} ($L_1e^\alpha=0$, $L_1h_a(-1)\chseventeenvac=0$, the
weights of the generators, and $L_1$ acting factorwise).  \textbf{A stale docstring:} the module docstring of
\texttt{MatterAdjointReduction} calls the opposite iterate identity ``the sole residual hypothesis''; it has since been
proved.  The file \texttt{MatterOpposite} also contains the earlier reduction
\chseventeenL{monsterLieBilinearForm_invariant_of_oppositeAdjoint}, conditional on adjointness of the primaries' zero modes (a
hypothesis now also proved, as \chseventeenL{moonshineMatterForm_primary_zeroMode_skew} in Chapter~\ref{ch:18}).
\end{dictionary}

\begin{proof}[Proof on paper]
(i) Pull a submodule back along $b\mapsto\chseventeenvac\otimes b$, apply \chseventeenL{frameStateField_generated_by_charges_currents}, then
use $(a\otimes\chseventeenvac)_{(-1)}(\chseventeenvac\otimes b)=a\otimes b$.  (ii) Theorem~\ref{ch17:thm:genadjoint}(iv) together with the
identification of the opposite modes of the generators via their weights and $L_1$-towers.  (iii)
Theorem~\ref{ch17:thm:oppborcherds}(iii), then specialise to weight-one primaries at mode $0$, where
$u^{\mathrm{op}}_{(0)}=-u_{(0)}$, giving $B(u_{(0)}b,c)=-B(b,u_{(0)}c)$, the invariance of $B([u,v],w)$.
\end{proof}

\begin{physics}
The Monster Lie algebra's invariant form is the restriction of the string's BPZ inner product to physical states:
$\langle[u,v],w\rangle=\langle u,[v,w]\rangle$ because the zero mode of a physical current is anti-self-adjoint.
\end{physics}

\begin{status}
(i), (ii) proved with no hypotheses beyond the closure data displayed in (i); (iii) proved conditional on \chseventeenL{hIter}, which is itself proved; hence the final
invariance is unconditional (Chapter~\ref{ch:18}).  Printed in \leanpath{VOA/\allowbreak Moonshine/\allowbreak MoonshineCNAudit13.lean}; the
discharging theorem in \leanpath{VOA/\allowbreak Moonshine/\allowbreak MoonshineCNAudit16.lean}; standard axioms.
\end{status}

%======================================================================
\section{Momentum sectors and momentum reversal}\label{ch17:sec:sector}
%======================================================================

The isotropic model of \S\ref{ch17:sec:hyp} has to be connected with the actual sectors $\chseventeenM^\alpha$.  Two files of this
chapter do the linear algebra; the comparison of Virasoro actions and the transfer of the no-ghost theorem are in
Chapter~\ref{ch:18}.

\begin{theorem}[The sector equivalence]\label{ch17:thm:sectorequiv}
For every $\alpha\in\chseventeenII$ there is a linear isomorphism $\Psi_\alpha:\Vnat\otimes\chseventeenF^{\mathrm{iso}}\to\chseventeenM^\alpha$,
$v\otimes p\mapsto v\otimes\Phi(p)e^\alpha$, where $\Phi:\chseventeenF^{\mathrm{iso}}\to\chseventeenF_2$ is the algebra isomorphism changing the
isotropic variables into the complex-frame variables ($a_{j}\mapsto\frac12(x_{j,1}-ix_{j,2})$,
$b_j\mapsto x_{j,1}+ix_{j,2}$); and $\Psi_\alpha$ intertwines the natural actions of $\Aut(\Vnat)$.
\end{theorem}

\noindent File \leanpath{VOA/\allowbreak Moonshine/\allowbreak MoonshineCNMatterSectorEquiv.lean}.
\begin{leancode}
/-- The actual momentum-sector equivalence, including the complex-frame change of variables. -/
def moonshineMatterSectorEquiv (α : lorentzianLattice.carrier) :
    (Vnat ⊗[ℂ] hypFock) ≃ₗ[ℂ] moonshineMatterMomentum α :=

theorem moonshineMatterSectorEquiv_tmul (α : lorentzianLattice.carrier) (v : Vnat) (p : hypFock) :
    (moonshineMatterSectorEquiv α (v ⊗ₜ p) : moonshineMatterSpace) =
      v ⊗ₜ latStateGen lorentzianLattice (hypLorentzianPolynomialEquiv p) α := rfl

theorem moonshineMatterSectorEquiv_equivariant (g : AutVnat) (α : lorentzianLattice.carrier)
    (u : Vnat ⊗[ℂ] hypFock) :
    (moonshineMatterSectorEquiv α (vnatSectorAction g u) : moonshineMatterSpace) =
      moonshineMatterAut g (moonshineMatterSectorEquiv α u) := by
\end{leancode}

\begin{dictionary}
\chseventeenL{hypLorentzianPolynomialEquiv : hypFock ≃ₐ[ℂ] FockN (rank lorentzianLattice)} is defined in the Chapter~\ref{ch:18}
file \leanpath{VOA/\allowbreak Moonshine/\allowbreak MoonshineCNSectorFrame.lean} (images of $a_j$, $b_j$ as in the theorem; with this, family
$a$ is the current of $(0,1)$, so $\kappa=x$, $\lambda=y$ at $\alpha=(x,y)$).  \chseventeenL{AutVnat} is the automorphism group of $\Vnat$
(Chapter~\ref{ch:20}); \chseventeenL{vnatSectorAction g} $=g\otimes1$ and \chseventeenL{moonshineMatterAut g} $=g\otimes1$ on $\chseventeenM$
(Chapter~\ref{ch:18} files).  The equation is in the ambient space $\chseventeenM$ (coercion \chseventeenL{( … : moonshineMatterSpace)}).
\textbf{Not asserted here}: that $\Psi_\alpha$ intertwines Virasoro actions (the module docstring says so explicitly; this is
\texttt{SectorVirasoroComparison}, Chapter~\ref{ch:18}).
\end{dictionary}

\begin{proof}[Proof on paper]
$\chseventeenM^\alpha=\Vnat\otimes\chseventeenF_2e^\alpha$ and $p\mapsto pe^\alpha$ is an isomorphism $\chseventeenF_2\to\chseventeenF_2e^\alpha$; compose with $\Phi$.
$g\otimes1$ commutes with anything acting on the second factor.
\end{proof}

\begin{physics}
At fixed spacetime momentum the string Fock space is ``internal CFT $\otimes$ light-cone oscillators''; the Monster acts on
the internal factor only.
\end{physics}

\begin{status}
Proved; printed in \leanpath{VOA/\allowbreak Moonshine/\allowbreak MoonshineCNAudit10.lean}; standard axioms.
\end{status}

\begin{definition}[Momentum reversal]\label{ch17:def:reversal}
$\Theta=1\otimes\theta$ on $\chseventeenM$, where $\theta(p\otimes e^\alpha)=\sigma(p)\otimes e^{-\alpha}$ is the lift of $-1$ (Chapter~\ref{ch:04},
Definition~\ref{ch04:def:theta}) and $\sigma$ the Fock involution changing the sign of every oscillator
(Definition~\ref{ch04:def:sigma}).
\end{definition}

\noindent File \leanpath{VOA/\allowbreak Moonshine/\allowbreak MoonshineCNMomentumReversal.lean}.
\begin{leancode}
def moonshineMatterReversal : moonshineMatterSpace ≃ₗ[ℂ] moonshineMatterSpace :=
  (thetaLat lorentzianLattice).lTensor Vnat
\end{leancode}

\begin{dictionary}
\chseventeenL{thetaLat} is defined in \leanpath{VOA/\allowbreak Moonshine/\allowbreak LatticeInvolution.lean} for every integral lattice; its definition does not
use positivity.  \chseventeenL{moonshineMatterReversal_tmul} confirms $v\otimes pe^\alpha\mapsto v\otimes\sigma(p)e^{-\alpha}$ (no cocycle sign
in this convention).
\end{dictionary}

\begin{physics}
Charge conjugation of the light-cone bosons, $X^\pm\mapsto-X^\pm$: it sends momentum $\alpha$ to $-\alpha$.
\end{physics}

\begin{status}
Definition; printed in \leanpath{VOA/\allowbreak Moonshine/\allowbreak MoonshineCNAudit10.lean}.
\end{status}

\begin{theorem}[Momentum reversal and the cross-momentum pairing]\label{ch17:thm:momentumform}
$\Theta$ is an involution, maps $\chseventeenM^\alpha$ to $\chseventeenM^{-\alpha}$ and commutes with every $L_n$.  For polynomials
$p,q$: $B_L(\theta(pe^\alpha),qe^\alpha)=\chi(\alpha)\,B_{\mathcal F}(\sigma p,q)$, and for the sector states
$B(\Theta\Psi_\alpha(a\otimes p),\Psi_\alpha(b\otimes q))=\chi(\alpha)\,B_{\mathcal F}(\sigma\Phi p,\Phi q)\,B_\natural(a,b)$.
\end{theorem}

\noindent File \leanpath{VOA/\allowbreak Moonshine/\allowbreak MoonshineCNMomentumReversal.lean}.
\begin{leancode}
theorem moonshineMatterReversal_involutive : Function.Involutive moonshineMatterReversal := by

theorem moonshineMatterReversal_momentum (α : lorentzianLattice.carrier)
    (u : moonshineMatterSpace) (hu : u ∈ moonshineMatterMomentum α) :
    moonshineMatterReversal u ∈ moonshineMatterMomentum (-α) := by

theorem moonshineMatterReversal_virasoro (n : ℤ) (u : moonshineMatterSpace) :
    moonshineMatterReversal (moonshineMatterNearVOA.L n u) =
      moonshineMatterNearVOA.L n (moonshineMatterReversal u) := by

/-- Exact cross-momentum normalization for arbitrary oscillator polynomials. -/
theorem lorentzianReversal_form (α : lorentzianLattice.carrier)
    (p q : FockN (rank lorentzianLattice)) :
    lorentzianBilinearForm (thetaLat lorentzianLattice (latStateGen lorentzianLattice p α))
      (latStateGen lorentzianLattice q α) =
      chSign lorentzianLattice α * uForm (rank lorentzianLattice)
        (sigmaN (rank lorentzianLattice) p) q := by

/-- The tensor-sector pairing after momentum reversal, before identifying its polynomial part. -/
theorem moonshineMatterReversal_sector_pairing (α : lorentzianLattice.carrier)
    (a b : Vnat) (p q : hypFock) :
    moonshineMatterForm
      (moonshineMatterReversal (moonshineMatterSectorEquiv α (a ⊗ₜ p)))
      (moonshineMatterSectorEquiv α (b ⊗ₜ q)) =
      chSign lorentzianLattice α *
        uForm (rank lorentzianLattice)
          (sigmaN (rank lorentzianLattice) (hypLorentzianPolynomialEquiv p))
          (hypLorentzianPolynomialEquiv q) * natForm a b := by
\end{leancode}

\begin{dictionary}
\chseventeenL{chSign lorentzianLattice α} $=\chi(\alpha)=(-1)^{\Delta(\alpha)}\varepsilon(\alpha,\alpha)$, a sign; it is \emph{not} replaced by $1$, since the
library's cocycle is not the coordinate cocycle (Definition~\ref{ch17:def:stdcocycle}).  \chseventeenL{sigmaN} is $\sigma$.  The last
theorem is a formula for arbitrary polynomials, not only ground states; identifying
$B_{\mathcal F}(\sigma\Phi p,\Phi q)$ with the isotropic form $B^{\mathrm{iso}}(p,q)$ (up to a constant) is a further theorem, in
Chapter~\ref{ch:18} (\texttt{SectorFormComparison}).
\end{dictionary}

\begin{proof}[Proof on paper]
$\theta$ is an involution reversing momentum (Chapter~\ref{ch:04}); $\sigma$ commutes with the Fock Virasoro modes (they are
quadratic) and sends the frame oscillator of $\alpha$ to that of $-\alpha$ (\chseventeenL{lorentzianReversal_frameOscillator}), which is what the
momentum terms of $L_n$ require.  The pairing formulas are the definition of $B_L$ at momenta $-\alpha$, $\alpha$.
\end{proof}

\begin{physics}
The BPZ pairing pairs a state of momentum $\alpha$ with one of momentum $-\alpha$ (the out-state is the conjugate of the in-state).
To obtain a form on a single momentum sector, as the no-ghost theorem needs, one composes with charge conjugation; the
result is the sector contravariant form.
\end{physics}

\begin{status}
Proved; printed in \leanpath{VOA/\allowbreak Moonshine/\allowbreak MoonshineCNAudit10.lean}; standard axioms.
\end{status}

\begin{warning}
\textbf{Dead end: comparing forms on the same momentum.}  The matter form vanishes on two states of the same nonzero
momentum (it pairs $\alpha$ with $-\alpha$).  It is proved in Lean that \emph{no} map $F$ from $\Vnat\otimes\chseventeenF^{\mathrm{iso}}$ into a
single sector $\chseventeenM^\alpha$, $\alpha\neq0$, can satisfy $B(Fu,Fv)=c\cdot B^{\mathrm{iso}}$-pairing$(u,v)$ for a nonzero constant $c$
(\chseventeenL{moonshineMatterSector_no_same_momentum_isometry}, Chapter~\ref{ch:18} file
\leanpath{VOA/\allowbreak Moonshine/\allowbreak MoonshineCNSectorFormConvention.lean}; witness $\chseventeenvac\otimes1$, whose sector norm is $1$).  The
correct comparison inserts the reversal $\Theta$, as in Theorem~\ref{ch17:thm:momentumform}.
\end{warning}

%======================================================================
\section{Two Lie-theoretic lemmas}\label{ch17:sec:lie}
%======================================================================

Two files assigned to this chapter are pure Lie theory, used by the Borcherds--Kac--Moody part of Chapter~\ref{ch:18}
(freeness of the positive part via Verma modules).  They do not involve vertex algebras.

\begin{theorem}[A Lie algebra embeds in its enveloping algebra]\label{ch17:thm:enveloping}
For every complex Lie algebra $\mathfrak g$ (of any dimension), the canonical map $\iota:\mathfrak g\to U(\mathfrak g)$ is injective.
\end{theorem}

\noindent File \leanpath{VOA/\allowbreak Moonshine/\allowbreak MoonshineCNEnvelopingEmbedding.lean}.
\begin{leancode}
/-- The canonical map of a Lie algebra into its universal enveloping algebra is injective. -/
theorem envelopingEmbedding_injective :
    Function.Injective (UniversalEnvelopingAlgebra.ι ℂ : L → UniversalEnvelopingAlgebra ℂ L) := by
\end{leancode}

\begin{dictionary}
\chseventeenL{L} is a section variable: any type with \chseventeenL{[LieRing L] [LieAlgebra ℂ L]}.  Mathlib v4.30 does not provide this
injectivity for infinite-dimensional Lie algebras (it is a consequence of PBW).
\end{dictionary}

\begin{proof}[Proof on paper]
Choose a basis $(y_a)$ of $\mathfrak g$, well-ordered.  The repository's PBW representation (\leanpath{VOA/\allowbreak Foundations/\allowbreak PBWTheorem/LieRep.lean})
makes $\mathfrak g$ act on the space of finitely supported functions on multisets of indices, with $y_a\cdot\{\}=\{a\}$.  Since
$\iota$ followed by this representation applied to the bottom vector is injective ($\sum c_ay_a\mapsto\sum c_a\{a\}$), $\iota$ is
injective.
\end{proof}

\begin{physics}
The symmetry generators of a theory are faithfully represented among the products of generators; in physics this is
always assumed tacitly.
\end{physics}

\begin{status}
Proved, no hypotheses; printed in \leanpath{VOA/\allowbreak Moonshine/\allowbreak MoonshineCNAudit17.lean}; standard axioms.
\end{status}

\begin{theorem}[The induced vector map]\label{ch17:thm:inducedvector}
Let $\mathfrak p\subseteq\mathfrak g$ be a subalgebra, $S\subseteq\mathfrak g$ a subspace on which $\mathfrak p$ acts by $\sigma$ such that
$[x,u]-\sigma(x)u\in\mathfrak p$ for $x\in\mathfrak p$, $u\in S$.  Then $[a\otimes u]\mapsto[a\,u\otimes1]$ is a well-defined
$\mathfrak g$-equivariant map from $U(\mathfrak g)\otimes_{U(\mathfrak p)}S$ to $U(\mathfrak g)\otimes_{U(\mathfrak p)}\C$; and if it and the PBW
multiplication for $\sigma$ are injective, then right multiplication $U(\mathfrak q)\otimes S\to U(\mathfrak q)$ is injective for any subalgebra
$\mathfrak q$ containing $S$ (no complement condition is assumed; in the application $\mathfrak q$ is a complement of $\mathfrak p$).
\end{theorem}

\noindent File \leanpath{VOA/\allowbreak Moonshine/\allowbreak MoonshineCNInducedVectorMap.lean}.
\begin{leancode}
variable (σ : p →ₗ⁅ℂ⁆ End ℂ S)
variable (hσ : ∀ (x : p) (u : S), ⁅(x : g), (u : g)⁆ - (σ x u : g) ∈ p)

/-- The induced map `[a ⊗ u] ↦ [a u ⊗ 1]`. -/
def inducedVectorMap : relativeInducedSpace p σ →ₗ[ℂ] trivialInducedSpace p :=

/-- The induced map intertwines the two induced actions of the whole Lie algebra. -/
theorem inducedVectorMap_equivariant (z : g) (q : relativeInducedSpace p σ) :
    inducedVectorMap p S σ hσ (relativeInducedRep p σ z q) =
      relativeInducedRep p (0 : p →ₗ⁅ℂ⁆ End ℂ ℂ) z (inducedVectorMap p S σ hσ q) := by

/-- Injectivity of the induced map and of the PBW multiplication for `σ` give injectivity of
right multiplication `U(q) ⊗ S → U(q)`. -/
theorem generatorMultiplication_injective_of_inducedVectorMap
    (hΦ : Function.Injective (inducedVectorMap p S σ hσ))
    (hV : Function.Injective (relativeVermaMultiplication p q σ)) :
    Function.Injective
      (envelopingGeneratorMultiplication (S.comap q.toSubmodule.subtype : Submodule ℂ q)) := by
\end{leancode}

\begin{dictionary}
\chseventeenL{relativeInducedSpace p σ} $=U(\mathfrak g)\otimes_{U(\mathfrak p)}S$ and \chseventeenL{relativeVermaMultiplication} the PBW map
$U(\mathfrak q)\otimes S\to U(\mathfrak g)\otimes_{U(\mathfrak p)}S$ are defined in Chapter~\ref{ch:18} files
(\texttt{RelativeInduction}, \texttt{VermaMultiplication}).  \chseventeenL{hσ} is the compatibility making the map well defined (it is
used in the definition's proof obligation).  \chseventeenL{hΦ}, \chseventeenL{hV} are genuine hypotheses of this reduction lemma; they are
discharged where it is applied (Chapter~\ref{ch:18}).  Further variables: \chseventeenL{(q : LieSubalgebra ℂ g) (hSq : S ≤ q.toSubmodule)}.
\end{dictionary}

\begin{proof}[Proof on paper]
Well-definedness: $a x\otimes u-a\otimes\sigma(x)u\mapsto a(xu-\sigma(x)u)\otimes1=a([x,u]-\sigma(x)u)\otimes1+au\,x\otimes1$, and both terms
vanish in the trivially induced module because elements of $\mathfrak p$ placed last kill $1$.  Equivariance is left
multiplication.  Through the PBW maps, the induced map becomes right multiplication $a\otimes u\mapsto au$, so injectivity transfers.
\end{proof}

\begin{physics}
Pure representation theory: it is used to show that the positive part of the Monster Lie algebra is free (no relations
among the generators beyond those forced by Serre-type conditions), the algebraic shadow of the free-field structure of
the no-ghost spectrum.
\end{physics}

\begin{status}
Proved with the displayed hypotheses; printed in \leanpath{VOA/\allowbreak Moonshine/\allowbreak MoonshineCNAudit17.lean}; standard axioms.
\end{status}

%======================================================================
\section{Accounting for the 91 files}\label{ch17:sec:files}
%======================================================================

Every file assigned to this chapter by \texttt{Notes/Textbook/assignment.txt} (the rule for \texttt{ch17}) is listed here with
the place where it is documented (``ind.'' $=$ headline declarations presented individually; ``coll.'' $=$ described
collectively).  All file names are \texttt{VOA/Moonshine/MoonshineCN\textit{Name}.lean}.

\begin{center}
\small
\begin{longtable}{p{0.36\textwidth}p{0.08\textwidth}p{0.48\textwidth}}
\toprule
files & how & where \\ \midrule
\endhead
\texttt{Lattice} & ind. & \S\ref{ch17:sec:lattice}, Def.~\ref{ch17:def:lattice}--\ref{ch17:def:cocycle} \\
\texttt{Oscillators}, \texttt{ComplexFrame} & ind. & \S\ref{ch17:sec:frame}; \chseventeenL{toNearVOA} in \S\ref{ch17:sec:target} \\
\texttt{VertexTarget} & ind. & Def.~\ref{ch17:def:target}, Thm~\ref{ch17:thm:notbounded} \\
\texttt{Grading} & ind. & Thm~\ref{ch17:thm:grading} \\
\texttt{Exponential}, \texttt{ChargedFields} & ind. & Def.~\ref{ch17:def:exponentials}, \ref{ch17:def:charged} \\
\texttt{ChargedLocality} & ind. & Thm~\ref{ch17:thm:chargedlocal} \\
\texttt{HalfCommutation}, \texttt{PositiveContraction}, \texttt{NegativeContraction}, \texttt{NegativeExpansion}, \texttt{NegativeCollapse} & coll. & table in \S\ref{ch17:sec:construction} \\
\texttt{Currents}, \texttt{Translation} & ind. & Def.~\ref{ch17:def:currents} \\
\texttt{CurrentCharge}, \texttt{CreationTranslation}, \texttt{ChargeTranslation} & coll. & table in \S\ref{ch17:sec:construction} \\
\texttt{Reconstruction}, \texttt{StateFields} & ind. & Thm~\ref{ch17:thm:reconstruct} \\
\texttt{Conformal}, \texttt{Virasoro} & ind. & Def.~\ref{ch17:def:omega}, Thm~\ref{ch17:thm:virasoro} \\
\texttt{Sugawara} & coll. & proof of Thm~\ref{ch17:thm:virasoro}, table \\
\texttt{ConformalTranslation}, \texttt{GradedFields} & ind. & Thm~\ref{ch17:thm:virasoro} \\
\texttt{Realization} & ind. & Thm~\ref{ch17:thm:realization} \\
\texttt{ConformalBasis} & ind. & Thm~\ref{ch17:thm:dualbasis} \\
\texttt{Momentum} & ind. & Def.~\ref{ch17:def:momentum}, Thm~\ref{ch17:thm:momentum} \\
\texttt{HyperbolicFock}, \texttt{HyperbolicSugawara}, \texttt{HyperbolicCommutant}, \texttt{HyperbolicVirasoro}, \texttt{HyperbolicMixedModes}, \texttt{HyperbolicModeGrading}, \texttt{HyperbolicForm}, \texttt{HyperbolicSugawaraAdjoint}, \texttt{HyperbolicOneFamily}, \texttt{HyperbolicFiltration}, \texttt{HyperbolicPBWMap}, \texttt{HyperbolicPBWEquiv}, \texttt{HyperbolicLightCone}, \texttt{HyperbolicSwap} & ind. & \S\ref{ch17:sec:hyp}, Def.~\ref{ch17:def:hypfock}--Thm~\ref{ch17:thm:hyppbw} \\
\texttt{HyperbolicNegativeModes}, \texttt{HyperbolicNegativeVirasoro}, \texttt{HyperbolicPositiveCreation}, \texttt{HyperbolicCentralScalar}, \texttt{HyperbolicFilteredWords}, \texttt{HyperbolicPBWCoordinates} & coll. & end of \S\ref{ch17:sec:hyp} \\
\texttt{TensorFields}, \texttt{TensorConformal}, \texttt{TensorGrading}, \texttt{TensorRealization} & ind. & Def.~\ref{ch17:def:tensor}, Thm~\ref{ch17:thm:tensor}, \ref{ch17:thm:c26} \\
\texttt{TensorMomentum} & ind. & Thm~\ref{ch17:thm:mattermomentum} \\
\texttt{LatticeForm}, \texttt{TensorForm}, \texttt{TensorNondegenerate} & ind. & Def.~\ref{ch17:def:nearform}, \ref{ch17:def:matterform}, Thm~\ref{ch17:thm:matterform} \\
\texttt{TensorNegative}, \texttt{TensorPBWMap}, \texttt{TensorPBWEquiv}, \texttt{TensorLightCone}, \texttt{TensorLightConeKernels} & ind. & Def.~\ref{ch17:def:tensornegative}, Thm~\ref{ch17:thm:tensorlightcone} \\
\texttt{NearGoddard}, \texttt{NearTranslation}, \texttt{NearNthProduct}, \texttt{NearIterate}, \texttt{NearBorcherds}, \texttt{NearSkew}, \texttt{NearConsequences} & ind. & \S\ref{ch17:sec:near} \\
\texttt{NearMatterIdentities} & coll. & after Thm~\ref{ch17:thm:nearconseq} \\
\texttt{ChargeCutoffWitness} & ind. & Thm~\ref{ch17:thm:cutoffwitness} \\
\texttt{NearNilpotence}, \texttt{ChargeNilpotence}, \texttt{MatterNilpotence} & ind. & Thm~\ref{ch17:thm:nilpotence} \\
\texttt{NearOppositeModes}, \texttt{MatterOpposite} & ind. & Def.~\ref{ch17:def:opposite}, Thm~\ref{ch17:thm:matterfhl} \\
\texttt{NearAdjointClosure} & ind. & Def.~\ref{ch17:def:oppiterate}, Thm~\ref{ch17:thm:oppborcherds} \\
\texttt{NearOppositeIterate} & ind. & Thm~\ref{ch17:thm:oppborcherds} \\
\texttt{NearOppositeComul}, \texttt{NearOppositeLeft}, \texttt{NearOppositeRight} & coll. & proof of Thm~\ref{ch17:thm:oppborcherds} \\
\texttt{FockAdjoint}, \texttt{DressingAdjoint}, \texttt{ChargedAdjoint}, \texttt{CurrentAdjoint}, \texttt{TensorGeneratorAdjoint} & ind. & Thm~\ref{ch17:thm:genadjoint} \\
\texttt{GeneratorPrimary} & coll. & Thm~\ref{ch17:thm:cutoffwitness}, \ref{ch17:thm:matterfhl} \\
\texttt{MatterAdjointReduction} & ind. & Thm~\ref{ch17:thm:matterfhl} \\
\texttt{MatterSectorEquiv} & ind. & Thm~\ref{ch17:thm:sectorequiv} \\
\texttt{MomentumReversal} & ind. & Def.~\ref{ch17:def:reversal}, Thm~\ref{ch17:thm:momentumform} \\
\texttt{EnvelopingEmbedding}, \texttt{InducedVectorMap} & ind. & \S\ref{ch17:sec:lie} \\
\bottomrule
\end{longtable}
\end{center}

\noindent The count is $4$ (lattice, frame, target) $+2$ (grading, momentum) $+20$ (construction) $+2$ (realisation, basis)
$+20$ (\texttt{Hyperbolic*}) $+13$ (\texttt{Tensor*}) $+2$ (\texttt{LatticeForm}, \texttt{FockAdjoint}) $+15$ (\texttt{Near*})
$+3$ (\texttt{ChargeCutoffWitness}, \texttt{ChargeNilpotence}, \texttt{ChargedAdjoint}) $+2$ (\texttt{CurrentAdjoint},
\texttt{DressingAdjoint}) $+4$ (\texttt{Matter*}) $+1$ (\texttt{GeneratorPrimary}) $+1$ (\texttt{MomentumReversal}) $+2$ (Lie
lemmas) $=91$.  None of them is an audit file; the audits that print their declarations
(\texttt{MoonshineCNAudit}, \texttt{Audit2}, \texttt{3}, \texttt{5}, \texttt{6}, \texttt{7}, \texttt{9}, \texttt{10},
\texttt{12}, \texttt{13}, \texttt{16}, \texttt{17}) belong to Chapter~\ref{ch:16}.

%======================================================================
\section{Recorded traps, dead ends and documentation defects}\label{ch17:sec:traps}
%======================================================================

\begin{warning}
\textbf{1. No positive-definiteness witness.}  The positive-definite lattice construction cannot be instantiated at
$\chseventeenII$ (\chseventeenL{lorentzian_not_posDef}); the library duplicated it with a complex-frame parameter instead of
introducing a false hypothesis.  \textbf{2. Two cocycles.}  \chseventeenL{lorentzianChargedField} (cocycle $(-1)^{x_1y_0}$) and the
fields of \chseventeenL{lorentzianNearVOA} (FLM cocycle) are different families; both are local, they are \emph{not} claimed equal, and
no gauge equivalence is formalised.  Consequently the sign
$\chi(\alpha)=$ \chseventeenL{chSign} appears in every cross-momentum pairing and must not be replaced by $1$.  \textbf{3. The
total-weight FHL cutoff is false} in a near-VOA (Theorem~\ref{ch17:thm:cutoffwitness}); the cutoff is the level above the
sector ground energy.  \textbf{4. Same-momentum form comparison is impossible} (warning in \S\ref{ch17:sec:sector}).
\textbf{5. Two contravariant forms with opposite oscillator signs.} The isotropic form makes oscillators self-adjoint, the
lattice form skew-adjoint (\S\ref{ch17:sec:hyp}); they are compared only through the reversal $\Theta$ and an explicit
change of variables.  \textbf{6. The sign trap} $(m,n)\leftrightarrow(m,-n)$ between Borcherds' coordinates and the repository's
(\S\ref{ch17:sec:conventions}).  \textbf{7. \chseventeenL{NearInvariantForm} is only Virasoro adjointness}, not FHL invariance;
FHL invariance of the matter form is a separate theorem (\S\ref{ch17:sec:fhl}).  \textbf{8. A primary-only shortcut is
insufficient}: restricting the opposite iterate identity to weight-one primaries does not close the generator argument,
because generation passes through products of arbitrary weights; the general identity
was needed and is proved.  \textbf{9. An overstated obstruction}: it was at first believed that the VOA proof of the
contragredient closure needed the cutoff $p+1$ and the module type \chseventeenL{OrdVModule}; closer inspection found
that $p+1$ is used only as a box bound and a nilpotence certificate, both replaceable by arbitrary nilpotence exponents
(\S\ref{ch17:sec:fhl}).
\end{warning}

\paragraph{Documentation defects.}  None of the following affects a Lean statement.
(D1) The docstrings of \chseventeenL{EvenLatticeVertexRealization} (``An inhabitant \dots\ is still required'') and
\chseventeenL{LorentzianVertexTarget} (``The exact remaining existence statement'') in
\leanpath{VOA/\allowbreak Moonshine/\allowbreak MoonshineCNVertexTarget.lean} are stale: \chseventeenL{lorentzianVertexTarget_inhabited} proves the target.
(D2) The module docstring of \leanpath{VOA/\allowbreak Moonshine/\allowbreak MoonshineCNMatterAdjointReduction.lean} (``That identity is the sole
residual hypothesis'') and the docstring of \chseventeenL{OppositeIterateIdentity} (``The precise remaining opposite-field identity'')
are stale since \chseventeenL{oppositeIterateIdentity_of_locallyNilpotent}.  (D3) A route note speaks of a ``grading by $\chseventeenII$ with finite-dimensional pieces''; the momentum pieces of $\chseventeenM$ are
infinite-dimensional (each is $\Vnat\otimes\chseventeenF_2$), only the joint (momentum, weight) pieces are finite-dimensional
(the Lean files state this correctly).  (D4) A blueprint note labels the
zero modes as $\kappa=(e,r)=y$, $\lambda=(f,r)=x$ with $a\leftrightarrow e=(1,0)$; in the Lean identification
(\chseventeenL{hypLorentzianSector_zeroCharges}) family $a$ is the current of $(0,1)$ and $\kappa=x$.  The products $\kappa\lambda=xy$ agree,
so no theorem is affected.  (D5) Minor redundancy: \chseventeenL{UnboundedConformalVertexAlgebra} duplicates the fields of
\chseventeenL{NearVOA}, and \chseventeenL{nearVirasoroMode} duplicates \chseventeenL{NearVOA.L}; both are harmless.

%======================================================================
\section{Machine verification for this chapter}\label{ch17:sec:audit}
%======================================================================

The following is the build record of the release (Chapter~\ref{ch:15}, Section~\ref{ch15:sec:record}), cross-checked
against the audit files' imports and \texttt{\#print axioms} lines in the source files of the release.

\begin{plaincode}
audit module (VOA/Moonshine/)   files of this chapter it imports            #print axioms lines
MoonshineCNAudit.lean           Lattice Oscillators VertexTarget            53
MoonshineCNAudit2.lean          ComplexFrame ... Realization Momentum (25)  383
MoonshineCNAudit3.lean          Tensor{Fields,Conformal,Grading,Realization,
                                Momentum,Form,Nondegenerate} FockAdjoint
                                LatticeForm                                 105
MoonshineCNAudit5.lean          Hyperbolic{Fock,Commutant,Sugawara,Form}    117
MoonshineCNAudit6.lean          Hyperbolic{NegativeModes,...,Swap} (10)     98
MoonshineCNAudit7.lean          Hyperbolic{CentralScalar,...} Tensor{Negative,
                                PBWMap,PBWEquiv,LightCone,LightConeKernels} 94
MoonshineCNAudit9.lean          Near{Goddard,...,MatterIdentities} (8)      105
MoonshineCNAudit10.lean         MatterSectorEquiv MomentumReversal          83
MoonshineCNAudit12.lean         DressingAdjoint ChargedAdjoint CurrentAdjoint
                                TensorGeneratorAdjoint                      66
MoonshineCNAudit13.lean         NearNilpotence ChargeNilpotence MatterNilpotence
                                NearOppositeModes NearAdjointClosure
                                MatterOpposite GeneratorPrimary
                                MatterAdjointReduction ChargeCutoffWitness  60
MoonshineCNAudit16.lean         NearOpposite{Comul,Left,Right,Iterate}      43
MoonshineCNAudit17.lean         EnvelopingEmbedding InducedVectorMap        141
\end{plaincode}

\noindent (The line counts include declarations of other chapters' files imported by the same audit.)  By the build record of the release, each audit builds green
and every printed axiom list is a subset of \chseventeenL{[propext, Classical.choice, Quot.sound]}, with zero \chseventeenL{sorryAx}; the
project axiom \chseventeenL{VOA.pbw} (now a theorem) is not involved.  The umbrella \lean{VOA.Moonshine} builds green (9,504 jobs), and \texttt{run\_axiom\_audits.py} re-runs every audit (Chapter~\ref{ch:15}).  On that
record every headline of this chapter is \texttt{\#print axioms}-clean; it is unconditional
except where a hypothesis is displayed in its Lean statement, and every such hypothesis is either a mathematical premise of a
general theorem (evenness, a frame, nonzero momentum, a Witt relation for an arbitrary operator family, generation and
generator adjointness) or a conditional reduction whose hypothesis is proved elsewhere (\chseventeenL{hIter}).

%======================================================================
\section{Referee's checklist}\label{ch17:sec:checklist}
%======================================================================

\begin{center}
\small
\begin{longtable}{p{0.26\textwidth}p{0.32\textwidth}p{0.19\textwidth}p{0.1\textwidth}}
\toprule
declaration & intended meaning & hypotheses & check \\ \midrule
\endhead
\chseventeenL{lorentzianLattice} & $\chseventeenII$, Gram $\left(\begin{smallmatrix}0&1\\1&0\end{smallmatrix}\right)$ & — (definition) & Def~\ref{ch17:def:lattice}\\
\chseventeenL{lorentzian_even}, \chseventeenL{lorentzian_selfDual}, \chseventeenL{lorentzian_not_posDef}, \chseventeenL{lorentzian_gram_det} & even, unimodular, indefinite, rank 2 & none & Thm~\ref{ch17:thm:latticeprops}\\
\chseventeenL{lorentzianEnergy_unbounded_below} & energies unbounded below & none & Thm~\ref{ch17:thm:energy}\\
\chseventeenL{ComplexLatticeFrame}, \chseventeenL{lorentzianComplexFrame} & bilinear orthonormal frame; constructed for $\chseventeenII$ & — & Def~\ref{ch17:def:frame}\\
\chseventeenL{frameOscillator_commutator} & lattice Heisenberg relations & a frame & Thm~\ref{ch17:thm:frame}\\
\chseventeenL{NearVOA} & VA + Virasoro axioms, no finiteness/lower bound & — (structure) & Def~\ref{ch17:def:nearvoa}\\
\chseventeenL{EvenLatticeVertexRealization}, \chseventeenL{LorentzianVertexTarget} & what $\chseventeenVII$ must be & — & Def~\ref{ch17:def:target}\\
\chseventeenL{lorentzian_grading_not_bounded_below} & no grading with physical weights is bounded below & any grading & Thm~\ref{ch17:thm:notbounded}\\
\chseventeenL{latticeWeightGrading_isInternal}, \chseventeenL{lorentzianWeightZero_not_finiteDimensional} & grading internal for all $L$; $V_0$ infinite-dim. & none & Thm~\ref{ch17:thm:grading}\\
\chseventeenL{chargedField_local} & charged fields mutually local, both signs & frame, factor system & Thm~\ref{ch17:thm:chargedlocal}\\
\chseventeenL{frameStateField_exists} & reconstruction of $Y$ & even & Thm~\ref{ch17:thm:reconstruct}\\
\chseventeenL{frameStateField_virasoro} & Virasoro, $c=\operatorname{rank}L$ & even & Thm~\ref{ch17:thm:virasoro}\\
\chseventeenL{complexFrameLatticeRealization} & $V_L$ for indefinite even $L$ & even, nondeg., frame & Thm~\ref{ch17:thm:realization}\\
\chseventeenL{lorentzianNearVOA}, \chseventeenL{lorentzianNearVOA_centralCharge} & $\chseventeenVII$, $c=2$ & \textbf{none} & Thm~\ref{ch17:thm:realization}\\
\chseventeenL{lorentzianNearVOA_conformal_dualBasis} & $\omega=\frac12(u(-1)v(-1)+v(-1)u(-1))\chseventeenvac$ & none & Thm~\ref{ch17:thm:dualbasis}\\
\chseventeenL{latticeMomentumGrading_isInternal}, \chseventeenL{lorentzianNearVOA_momentum} & momentum grading, conserved & none & Thm~\ref{ch17:thm:momentum}\\
\chseventeenL{hypOscillator_comm} & light-cone Heisenberg relations & none & Thm~\ref{ch17:thm:hypheis}\\
\chseventeenL{hypCommutant_scalar} & commutant is scalar & commutation & Thm~\ref{ch17:thm:hypcommutant}\\
\chseventeenL{hypSugawara_virasoro} & isotropic Sugawara, $c=2$ & none & Thm~\ref{ch17:thm:hypvirasoro}\\
\chseventeenL{hypForm_sugawara_adjoint}, \chseventeenL{hypForm_weight_nondegenerate} & contravariant, levelwise nondegenerate & none & Thm~\ref{ch17:thm:hypform}\\
\chseventeenL{hypOneFamilyQuotientEquiv} & one-family quotient $\cong\C$ & none & Thm~\ref{ch17:thm:onefamily}\\
\chseventeenL{hypPBWEquiv}, \chseventeenL{hypFock_lightConeSplit_nonzero} & light-cone splitting & $\kappa\ne0$ / $(\kappa,\lambda)\ne0$ & Thm~\ref{ch17:thm:hyppbw}\\
\chseventeenL{nearTensorVOA} & tensor product of near-VOAs & none & Thm~\ref{ch17:thm:tensor}\\
\chseventeenL{moonshineMatterNearVOA_centralCharge} & $\Vnat\otimes\chseventeenVII$ has $c=26$ & \textbf{none} & Thm~\ref{ch17:thm:c26}\\
\chseventeenL{moonshineMatterMomentum_isInternal}, \chseventeenL{moonshineMatterMomentum_mode} & $\chseventeenII$-grading of $\chseventeenM$ & none & Thm~\ref{ch17:thm:mattermomentum}\\
\chseventeenL{moonshineMatterForm_nondegenerate}, \chseventeenL{moonshineMatterForm_virasoro_adjoint} & nondegenerate, $L_n^\dagger=L_{-n}$ & none & Thm~\ref{ch17:thm:matterform}\\
\chseventeenL{hypTensor_lightConeSplit_nonzero} & tensor light-cone splitting & Witt relation \chseventeenL{hD}, $(\kappa,\lambda)\ne0$ & Thm~\ref{ch17:thm:tensorlightcone}\\
\chseventeenL{goddard_uniqueness}, \chseventeenL{jacobi}, \chseventeenL{smul_eq_skew} & Goddard, Borcherds, skew for near-VOAs & none & Thms~\ref{ch17:thm:goddard}, \ref{ch17:thm:nearborcherds}\\
\chseventeenL{primary_comm}, \chseventeenL{zeroMode_preserves_spurious}, \chseventeenL{skew_zero} & physical-state bracket facts & weight-1 primary & Thm~\ref{ch17:thm:nearconseq}\\
\chseventeenL{lorentzianVirasoroOne_weightZero_not_zero} & total-weight cutoff fails & none & Thm~\ref{ch17:thm:cutoffwitness}\\
\chseventeenL{moonshineMatterVirasoroOne_cutoff} & $L_1^{N-\Delta(\alpha)+1}u=0$ & weight $N$, momentum $\alpha$ & Thm~\ref{ch17:thm:nilpotence}\\
\chseventeenL{oppositeIterateIdentity_of_locallyNilpotent} & opposite iterate identity & $L_1$ loc.\ nilpotent & Thm~\ref{ch17:thm:oppborcherds}\\
\chseventeenL{adjoint_of_locallyNilpotent_generators} & adjoint generators $\Rightarrow$ all fields adjoint & nilp., generation, gen.\ adjointness & Thm~\ref{ch17:thm:oppborcherds}\\
\chseventeenL{latForm_frameStateField_charge_adjoint}, \chseventeenL{moonshineMatterForm_vnatGenerator_fhl} & generators FHL-adjoint & even (discharged) & Thm~\ref{ch17:thm:genadjoint}\\
\chseventeenL{monsterLieBilinearForm_invariant_of_oppositeIterate} & invariance of Monster Lie form & \chseventeenL{hIter} (proved, Ch.~18) & Thm~\ref{ch17:thm:matterfhl}\\
\chseventeenL{moonshineMatterSectorEquiv_equivariant} & sector $\cong\Vnat\otimes\chseventeenF^{\mathrm{iso}}$, $\Aut$-equivariant & none & Thm~\ref{ch17:thm:sectorequiv}\\
\chseventeenL{moonshineMatterReversal_sector_pairing} & cross-momentum pairing formula & none & Thm~\ref{ch17:thm:momentumform}\\
\chseventeenL{envelopingEmbedding_injective} & $\mathfrak g\hookrightarrow U(\mathfrak g)$ & none & Thm~\ref{ch17:thm:enveloping}\\
\chseventeenL{generatorMultiplication_injective_of_inducedVectorMap} & reduction of freeness to induced map & \chseventeenL{hΦ}, \chseventeenL{hV} & Thm~\ref{ch17:thm:inducedvector}\\
\bottomrule
\end{longtable}
\end{center}

All rows are \texttt{\#print axioms}-clean (standard axioms only).  Rows marked ``none'' are unconditional; the other
hypotheses are premises of general theorems, discharged for $\chseventeenII$, $\chseventeenM$ and $\Vnat$ where the theorems are
applied, except \chseventeenL{hIter}, which is proved in Chapter~\ref{ch:18}.

%======================================================================
\section{Notes and references}\label{ch17:sec:notes}
%======================================================================

The lattice vertex algebras of even lattices, including indefinite ones, and the use of $\Vnat\otimes V_{\chseventeenII}$ are due
to R.~E.~Borcherds, \emph{Vertex algebras, Kac--Moody algebras, and the Monster}, Proc.\ Natl.\ Acad.\ Sci.\ USA 83 (1986)
3068--3071, and \emph{Monstrous moonshine and monstrous Lie superalgebras}, Invent.\ Math.\ 109 (1992) 405--444 (\S\S2--5:
the vertex algebra of $\chseventeenII$, the no-ghost theorem, the Monster Lie algebra).  The construction of lattice vertex
operators and the cocycle follows I.~Frenkel, J.~Lepowsky and A.~Meurman, \emph{Vertex Operator Algebras and the Monster},
Academic Press, 1988 (Chapters~7--8); the vertex operators $Y(e^\alpha,z)$ go back to I.~B.~Frenkel and V.~G.~Kac, \emph{Basic
representations of affine Lie algebras and dual resonance models}, Invent.\ Math.\ 62 (1980) 23--66, and G.~Segal, \emph{Unitary
representations of some infinite-dimensional groups}, Comm.\ Math.\ Phys.\ 80 (1981) 301--342.  Near-VOAs (``Definition
5.1.3(c)'') are from T.~Gannon, \emph{Moonshine beyond the Monster}, Cambridge University Press, 2006, \S5.1.  The invariant
bilinear form and the contragredient formula are from I.~B.~Frenkel, Y.-Z.~Huang and J.~Lepowsky, \emph{On axiomatic
approaches to vertex operator algebras and modules}, Mem.\ Amer.\ Math.\ Soc.\ 104 (1993), no.~494, and H.~Li, \emph{Symmetric
invariant bilinear forms on vertex operator algebras}, J.\ Pure Appl.\ Algebra 96 (1994) 279--297.  Goddard uniqueness,
Dong's lemma, reconstruction and the Borcherds identity are presented as in V.~Kac, \emph{Vertex Algebras for Beginners},
2nd ed., University Lecture Series 10, AMS, 1998 (\S\S3--4), and E.~Frenkel and D.~Ben-Zvi, \emph{Vertex Algebras and Algebraic
Curves}, 2nd ed., AMS, 2004.

The no-ghost theorem is due to R.~C.~Brower, \emph{Spectrum-generating algebra and no-ghost theorem for the dual model},
Phys.\ Rev.\ D 6 (1972) 1655--1662, and P.~Goddard and C.~B.~Thorn, \emph{Compatibility of the dual Pomeron with unitarity and
the absence of ghosts in the dual resonance model}, Phys.\ Lett.\ B 40 (1972) 235--238; the light-cone splitting of
\S\ref{ch17:sec:hyp} is the Goddard--Thorn decomposition in the form used by Borcherds (1992, \S5).  The BRST
interpretation is I.~B.~Frenkel, H.~Garland and G.~J.~Zuckerman, \emph{Semi-infinite cohomology and string theory}, Proc.\
Natl.\ Acad.\ Sci.\ USA 83 (1986) 8442--8446.  For the physics (critical dimension, light-cone gauge, physical states,
BPZ conjugation) see J.~Polchinski, \emph{String Theory}, Vol.~1, Cambridge University Press, 1998, Chapters~1--4 and~8; and
for the Monster string, J.~A.~Harvey and G.~Moore, \emph{Algebras, BPS states, and strings}, Nucl.\ Phys.\ B 463 (1996)
315--368 (arXiv:hep-th/9510182), and S.~Carnahan, \emph{Monstrous Lie algebras and generalized moonshine}, arXiv:1208.6254.
A survey of the whole proof is T.~Gannon, \emph{Monstrous moonshine: the first twenty-five years}, Bull.\ London Math.\ Soc.\
38 (2006) 1--33 (arXiv:math/0402345).

The Lean development was recorded in the authors' working notes, which are not part of the release.  The positive-definite construction it adapts is documented in
Chapters~\ref{ch:03} and~\ref{ch:L05}; the vertex-algebra theory it ports, in Chapter~\ref{ch:L03}; invariant forms, in
Chapters~\ref{ch:04} and~\ref{ch:L04}.

\endgroup
